\documentclass[11pt,a4paper]{article}
\usepackage[margin=2.4cm]{geometry}
\usepackage{amsmath,amssymb,amsthm,mathtools}
\usepackage[protrusion=true,expansion=false]{microtype}
\usepackage{booktabs,enumitem}
\usepackage{bm}
\usepackage{xcolor}
\usepackage[colorlinks=true,linkcolor=blue!50!black,citecolor=blue!50!black,urlcolor=blue!50!black]{hyperref}
\emergencystretch=3em
\hbadness=10000
\hfuzz=2pt
\newtheorem{theorem}{Theorem}[section]
\newtheorem{lemma}[theorem]{Lemma}
\newtheorem{proposition}[theorem]{Proposition}
\newtheorem{corollary}[theorem]{Corollary}
\theoremstyle{definition}
\newtheorem{definition}[theorem]{Definition}
\newtheorem{construction}[theorem]{Construction}
\newtheorem{example}[theorem]{Example}
\theoremstyle{remark}
\newtheorem{remark}[theorem]{Remark}
\newcommand{\N}{\ensuremath{\mathbb{N}}}
\newcommand{\Z}{\ensuremath{\mathbb{Z}}}
\newcommand{\Q}{\ensuremath{\mathbb{Q}}}
\newcommand{\R}{\ensuremath{\mathbb{R}}}
\providecommand{\C}{}\renewcommand{\C}{\ensuremath{\mathbb{C}}}
\newcommand{\F}{\ensuremath{\mathbb{F}}}
\newcommand{\G}{\ensuremath{\mathbb{G}}}
\newcommand{\Fp}{\ensuremath{\mathbb{F}_{p}}}
\newcommand{\Fq}{\ensuremath{\mathbb{F}_{q}}}
\newcommand{\E}{\ensuremath{\mathbb{E}}}
\newcommand{\ip}[2]{\langle #1,\, #2 \rangle}
\DeclareMathOperator{\ord}{ord}
\DeclareMathOperator{\lcm}{lcm}
\DeclareMathOperator{\negl}{negl}
\DeclareMathOperator{\poly}{poly}
\DeclareMathOperator{\Hom}{Hom}
\DeclareMathOperator{\img}{im}
\DeclareMathOperator{\Tr}{Tr}
\newcommand{\legendre}[2]{\left(\frac{#1}{#2}\right)}
\DeclareMathOperator{\sgn}{sgn}
\DeclareMathOperator{\rank}{rank}
\DeclareMathOperator{\End}{End}
\DeclareMathOperator{\Aut}{Aut}
\DeclareMathOperator{\id}{id}
\DeclareMathOperator{\Frob}{Frob}
\DeclareMathOperator{\disc}{disc}
\DeclareMathOperator{\charf}{char}
\providecommand{\abs}[1]{\lvert #1 \rvert}
\providecommand{\norm}[1]{\lVert #1 \rVert}
\providecommand{\ceil}[1]{\lceil #1 \rceil}
\providecommand{\floor}[1]{\lfloor #1 \rfloor}
\providecommand{\indicator}{\mathbf{1}}
\providecommand{\mathbbm}[1]{\mathbf{#1}}
\newcommand{\Fr}{\ensuremath{\mathbb{F}_{r}}}
\newcommand{\Adv}{\ensuremath{\mathcal{A}}}
\usepackage[nameinlink]{cleveref}

\title{\textbf{\Huge Math Guide}\\[6pt]\large Mathematical foundations for Halo 2 and Zcash Orchard}
\author{m@rek.onl}
\date{}
\begin{document}
\maketitle
\thispagestyle{empty}
\begin{abstract}
This is the first volume of a four-part series --- \emph{Math Guide}, \emph{Crypto Guide}, \emph{Halo~2 Guide}, \emph{Orchard Guide} --- developing the cryptography of Zcash's Orchard protocol from the ground up. Assuming only elementary algebra and calculus, it builds the mathematics the later volumes rely on: sets and notation; the integers and modular arithmetic; groups; rings and fields; polynomials; finite fields and the number theory of the discrete logarithm; linear algebra and inner products; roots of unity and evaluation domains; elliptic curves and their group law; and the probability, asymptotics, and model of computation needed to state cryptographic claims precisely.
\end{abstract}
\clearpage
\tableofcontents
\clearpage
\section{Notation: sets, functions, and relations}\label{sec:sets}
This short section fixes the small amount of set-theoretic language we use
throughout the series. A reader comfortable with elementary mathematics may
skim it and refer back as needed.

\subsection{Sets}\label{sets:subsec:sets}
A \emph{set} is a collection of distinct objects, its \emph{elements}; one writes
\(x \in A\) for membership and \(x \notin A\) for its negation. A set may be
given by listing, \(\{a,b,c\}\), or by \emph{set-builder} notation
\(\{x \in A : P(x)\}\) (``the \(x\) in \(A\) such that \(P(x)\)''). One writes
\(A \subseteq B\) when every element of \(A\) lies in \(B\), and \(A = B\) when
\(A \subseteq B\) and \(B \subseteq A\). The usual operations are union
\(A \cup B\), intersection \(A \cap B\), difference \(A \setminus B\), and the
empty set \(\varnothing\). The \emph{Cartesian product} is
\(A \times B = \{(a,b) : a \in A,\ b \in B\}\), and \(A^n\) is the set of
\(n\)-tuples over \(A\). For a finite set \(A\), \(\lvert A\rvert\) denotes its
number of elements. The standard number systems are the naturals \(\N\), the
integers \(\Z\), the rationals \(\Q\), the reals \(\R\), and the complex numbers
\(\C\), with \(\N \subset \Z \subset \Q \subset \R \subset \C\); we construct the prime field
\(\Fp\) in \S\ref{sec:intmod} and treat general finite fields \(\Fq\) in
\S\ref{sec:finitefields}.

\subsection{Functions}\label{sets:subsec:functions}
A \emph{function} \(f \colon A \to B\) assigns to each \(a \in A\) a unique value
\(f(a) \in B\); \(A\) is the \emph{domain} and \(B\) the \emph{codomain}. It is
\emph{injective} (one-to-one) if \(f(a) = f(a')\) implies \(a = a'\);
\emph{surjective} (onto) if every \(b \in B\) equals \(f(a)\) for some \(a\); and
\emph{bijective} if both. The \emph{composition} of \(f \colon A \to B\) and
\(g \colon B \to C\) is \((g \circ f)(a) = g(f(a))\). A function has a two-sided
\emph{inverse} \(f^{-1} \colon B \to A\) (with \(f^{-1}\circ f = \id_A\) and
\(f \circ f^{-1} = \id_B\)) if and only if it is bijective. For \(S \subseteq A\)
the \emph{image} is \(f(S) = \{f(a) : a \in S\}\), and for \(T \subseteq B\) the
\emph{preimage} is \(f^{-1}(T) = \{a \in A : f(a) \in T\}\) (defined whether or
not \(f\) is invertible).

\subsection{Relations and equivalences}\label{sets:subsec:relations}
A \emph{binary relation} on \(A\) is a subset \(R \subseteq A \times A\); one writes
\(a \sim b\) for \((a,b) \in R\). It is an \emph{equivalence relation} if it is
\emph{reflexive} (\(a \sim a\)), \emph{symmetric} (\(a \sim b \Rightarrow b \sim
a\)), and \emph{transitive} (\(a \sim b\) and \(b \sim c\) imply \(a \sim c\)).
The \emph{equivalence class} of \(a\) is \([a] = \{x \in A : x \sim a\}\); the
distinct classes \emph{partition} \(A\) into disjoint nonempty pieces whose union
is \(A\), and the set of classes is the \emph{quotient} \(A/{\sim}\). We use this
construction twice: to form the integers modulo \(n\) (\S\ref{sec:intmod}) and the cosets of a
subgroup (\S\ref{sec:groups}).

\subsection{Logic and proof conventions}\label{sets:subsec:logic}
We employ the symbols \(\forall\) (``for all''), \(\exists\) (``there exists''), \(\exists!\)
(``there exists a unique''), and ``\(\colon\)'' or ``s.t.'' for ``such that''.
Negation swaps the quantifiers: \(\neg(\forall x\, P(x))\) is \(\exists x\,\neg
P(x)\), and \(\neg(\exists x\, P(x))\) is \(\forall x\,\neg P(x)\). The order of
unlike quantifiers matters: \(\forall x\,\exists y\,P(x,y)\) (the \(y\) may depend
on \(x\)) is weaker than \(\exists y\,\forall x\,P(x,y)\) (one \(y\) works for
all). We prove an implication \(P \Rightarrow Q\) directly (assume \(P\), derive \(Q\)), by
contraposition (assume \(\neg Q\), derive \(\neg P\)), or by contradiction; we prove a statement
\(\forall x\,P(x)\) by taking an arbitrary \(x\), and \(\exists x\,P(x)\) by
exhibiting a witness. We prove statements about all natural numbers by
induction.
\section{The integers and modular arithmetic}\label{sec:intmod}

The integers \(\Z = \{\dots,-2,-1,0,1,2,\dots\}\) are the arithmetical bedrock on which every algebraic and cryptographic structure in this monograph ultimately rests. The finite fields underlying elliptic curves, the scalar arithmetic of Halo~2, and the polynomial commitments of Orchard all reduce, in the end, to computations with congruence classes of integers. This section develops the theory of the integers from a small set of assumptions and arrives at the ring \(\Z/n\Z\) of integers modulo \(n\), its group of units, and the algorithms that make arithmetic in it effective. We keep the treatment constructive throughout: every existence statement comes, wherever possible, with an algorithm that produces the object in question, since cryptography requires not merely that inverses and representatives exist but that one can compute them efficiently.

We take the basic algebraic properties of \(\Z\) for granted: that it is a commutative ring with the usual addition and multiplication, that it has no zero divisors (if \(ab=0\) then \(a=0\) or \(b=0\)), and that it carries a total order \(\le\) compatible with the arithmetic. We assume axiomatically the single deep property that follows.

\begin{theorem}[Well-ordering principle]\label{intmod:thm:wop}
Every nonempty subset \(S\subseteq\N=\{0,1,2,\dots\}\) of the natural numbers has a least element: there exists \(m\in S\) such that \(m\le s\) for all \(s\in S\).
\end{theorem}

The well-ordering principle is logically equivalent to the principle of mathematical induction; we use both freely. Nearly every existence result below is, at bottom, an application of well-ordering: one exhibits a nonempty set of natural numbers and extracts its minimum.

\subsection{Divisibility}\label{intmod:sub:divisibility}

\begin{definition}\label{intmod:def:divides}
Let \(a,b\in\Z\). One says that \(a\) \emph{divides} \(b\), written \(a\mid b\), if there exists an integer \(c\in\Z\) with \(b=ac\). In this case \(a\) is called a \emph{divisor} or \emph{factor} of \(b\), and \(b\) is a \emph{multiple} of \(a\). If no such \(c\) exists one writes \(a\nmid b\).
\end{definition}

Several immediate consequences of the definition deserve to be recorded; we shall use them constantly and usually without explicit citation.

\begin{proposition}\label{intmod:prop:divprops}
Let \(a,b,c,d\in\Z\). Then:
\begin{enumerate}
\item \(a\mid a\) (reflexivity), \(1\mid a\), and \(a\mid 0\).
\item If \(a\mid b\) and \(b\mid c\) then \(a\mid c\) (transitivity).
\item If \(a\mid b\) and \(a\mid c\) then \(a\mid (bx+cy)\) for all \(x,y\in\Z\); in particular \(a\mid (b\pm c)\).
\item If \(a\mid b\) and \(b\neq 0\) then \(|a|\le |b|\).
\item If \(a\mid b\) and \(b\mid a\) then \(a=\pm b\).
\item \(0\mid b\) if and only if \(b=0\).
\end{enumerate}
\end{proposition}

\begin{proof}
(1) is immediate: \(a=a\cdot 1\), \(a=1\cdot a\), \(0=a\cdot 0\). For (2), write \(b=ae\) and \(c=bf\); then \(c=a(ef)\), so \(a\mid c\). For (3), write \(b=ae\) and \(c=af\); then \(bx+cy=a(ex+fy)\). For (4), write \(b=ac\) with \(b\neq 0\), so \(c\neq 0\) and hence \(|c|\ge 1\); then \(|b|=|a|\,|c|\ge |a|\). For (5), if either is \(0\) then by (6) so is the other and the claim holds; otherwise (4) gives \(|a|\le|b|\) and \(|b|\le|a|\), so \(|a|=|b|\) and \(a=\pm b\). For (6), \(0\mid b\) means \(b=0\cdot c=0\); conversely \(0\mid 0\).
\end{proof}

The relation \(\mid\) is thus reflexive and transitive; it is a partial order when restricted to the natural numbers (where (5) becomes \(a=b\)), but on all of \(\Z\) it fails to be antisymmetric precisely because of signs. The \emph{units} of \(\Z\)—the integers dividing \(1\)—are exactly \(\pm 1\), and two integers that divide each other are called \emph{associates}; in \(\Z\) the associates of \(a\) are \(\pm a\).

\subsection{The division algorithm}\label{intmod:sub:divalg}

Divisibility is the special case of exact division. The fundamental fact that makes the arithmetic of \(\Z\) tractable is that one can always divide \emph{with remainder}, and that the quotient and remainder are uniquely determined once a convention for the remainder's range is fixed.

\begin{theorem}[Division algorithm]\label{intmod:thm:division}
Let \(a\in\Z\) and \(b\in\Z\) with \(b>0\). Then there exist unique integers \(q\) (the \emph{quotient}) and \(r\) (the \emph{remainder}) such that
\[
a = bq + r, \qquad 0 \le r < b.
\]
\end{theorem}

\begin{proof}
\emph{Existence.} Consider the set
\[
S = \{\, a - bk : k\in\Z,\ a-bk \ge 0 \,\} \subseteq \N.
\]
First, \(S\neq\varnothing\). If \(a\ge 0\) then \(a=a-b\cdot 0\in S\). If \(a<0\) then, since \(b\ge 1\), taking \(k=a\) gives \(a-ba=a(1-b)\ge 0\) (a product of a nonpositive and a nonpositive factor), so \(a-ba\in S\). Thus \(S\) is a nonempty subset of \(\N\), and by the well-ordering principle (Theorem~\ref{intmod:thm:wop}) it has a least element \(r=a-bq\) for some \(q\in\Z\). By construction \(r\ge 0\). If \(r\ge b\), then \(r-b = a-b(q+1)\ge 0\) would lie in \(S\) and be strictly smaller than \(r\), contradicting minimality. Hence \(0\le r<b\).

\emph{Uniqueness.} Suppose \(a=bq+r=bq'+r'\) with \(0\le r,r'<b\). Then \(b(q-q')=r'-r\), so \(b\mid (r'-r)\). But \(-b < r'-r < b\), and the only multiple of \(b\) strictly between \(-b\) and \(b\) is \(0\). Hence \(r'=r\), and then \(b(q-q')=0\) with \(b\neq 0\) forces \(q=q'\).
\end{proof}

\begin{remark}\label{intmod:rem:negdivisor}
The hypothesis \(b>0\) is only for definiteness. For general \(b\neq 0\) one obtains unique \(q,r\) with \(a=bq+r\) and \(0\le r<|b|\), by applying the theorem to \(|b|\) and adjusting the sign of \(q\). Note that the remainder convention \(0\le r<|b|\) differs from the truncation behaviour of many programming languages, where the remainder of a negative dividend may be negative; for the number-theoretic results below the convention above is the appropriate one.
\end{remark}

We denote the remainder \(r\) produced by the division algorithm by \(a \bmod b\), and the quotient \(q\) by \(\lfloor a/b\rfloor\) when \(b>0\) (the \emph{floor} of the rational number \(a/b\)). With this notation, \(b\mid a\) holds precisely when \(a\bmod b=0\).

\subsection{The greatest common divisor and the Euclidean algorithm}\label{intmod:sub:gcd}

\begin{definition}\label{intmod:def:commondiv}
A \emph{common divisor} of integers \(a\) and \(b\) is an integer \(d\) with \(d\mid a\) and \(d\mid b\). The \emph{greatest common divisor} of \(a\) and \(b\), not both zero, is the largest such \(d\); it is denoted \(\gcd(a,b)\). By convention \(\gcd(0,0)=0\). Integers \(a\) and \(b\) are \emph{coprime} (or \emph{relatively prime}) if \(\gcd(a,b)=1\).
\end{definition}

The greatest common divisor is well defined when \((a,b)\neq(0,0)\): the set of common divisors is nonempty (it contains \(1\)) and bounded above (every common divisor of, say, the nonzero \(a\) has absolute value at most \(|a|\) by Proposition~\ref{intmod:prop:divprops}(4)), so it has a largest element, which is positive since \(1\) is a common divisor. We establish a much stronger characterisation shortly: the gcd is not merely the numerically largest common divisor but is divisible by every common divisor. The key computational engine is the following observation.

\begin{lemma}\label{intmod:lem:gcdreduce}
Let \(a,b\in\Z\). For any \(q\in\Z\),
\[
\gcd(a,b) = \gcd(b,\,a-qb).
\]
In particular, if \(a=bq+r\) then \(\gcd(a,b)=\gcd(b,r)\), and \(\gcd(a,0)=|a|\).
\end{lemma}

\begin{proof}
It suffices to show the two pairs \((a,b)\) and \((b,a-qb)\) have exactly the same common divisors; equal sets of common divisors have equal greatest elements (and the same coprimality status). If \(d\mid a\) and \(d\mid b\), then \(d\mid(a-qb)\) by Proposition~\ref{intmod:prop:divprops}(3), so \(d\) is a common divisor of \(b\) and \(a-qb\). Conversely, if \(d\mid b\) and \(d\mid(a-qb)\), then \(d\mid\bigl((a-qb)+qb\bigr)=a\), so \(d\) is a common divisor of \(a\) and \(b\). The two sets of common divisors coincide. Finally \(\gcd(a,0)=|a|\) because the common divisors of \(a\) and \(0\) are exactly the divisors of \(a\), the largest of which in absolute value is \(|a|\) (and \(|a|>0\) when \(a\neq 0\); the case \(a=0\) is the convention \(\gcd(0,0)=0\)).
\end{proof}

Iterating the lemma with the division algorithm yields the oldest nontrivial algorithm in mathematics.

\begin{theorem}[Euclidean algorithm]\label{intmod:thm:euclid}
Let \(a,b\) be integers with \(b\neq 0\). Define a sequence of remainders by \(r_0=|a|\), \(r_1=|b|\), and, as long as \(r_{i}\neq 0\), let \(r_{i+1}=r_{i-1}\bmod r_i\), so that
\[
r_{i-1}=r_i q_i + r_{i+1}, \qquad 0\le r_{i+1}<r_i.
\]
The sequence \(r_1>r_2>\cdots\) of positive remainders is strictly decreasing, so some \(r_{n+1}=0\); the last nonzero remainder satisfies \(r_n=\gcd(a,b)\).
\end{theorem}

\begin{proof}
The remainders \(r_1,r_2,\dots\) are nonnegative integers and, while nonzero, strictly decreasing because \(r_{i+1}=r_{i-1}\bmod r_i<r_i\). A strictly decreasing sequence of nonnegative integers cannot be infinite (again by well-ordering: otherwise the set of values would have no least element), so the algorithm terminates with some \(r_{n+1}=0\) and \(r_n>0\). By Lemma~\ref{intmod:lem:gcdreduce}, each step preserves the gcd:
\[
\gcd(a,b)=\gcd(r_0,r_1)=\gcd(r_1,r_2)=\cdots=\gcd(r_n,r_{n+1})=\gcd(r_n,0)=r_n,
\]
using \(\gcd(|a|,|b|)=\gcd(a,b)\) (signs do not affect common divisors).
\end{proof}

\begin{example}\label{intmod:exa:euclid}
Compute \(\gcd(1071,462)\):
\[
\begin{aligned}
1071 &= 2\cdot 462 + 147, \\
462 &= 3\cdot 147 + 21, \\
147 &= 7\cdot 21 + 0.
\end{aligned}
\]
The last nonzero remainder is \(21\), so \(\gcd(1071,462)=21\). Indeed \(1071=21\cdot 51\) and \(462=21\cdot 22\), with \(\gcd(51,22)=1\).
\end{example}

\begin{remark}\label{intmod:rem:euclidcomplexity}
The Euclidean algorithm is efficient. The number of division steps for inputs \(a>b>0\) is \(O(\log b)\); the worst case occurs at consecutive Fibonacci numbers (Lamé's theorem), where each quotient \(q_i\) equals \(1\) and the remainders shrink as slowly as possible. With schoolbook arithmetic on \(k\)-bit inputs the whole computation costs \(O(k^2)\) bit operations. This efficiency is what renders gcd-based cryptographic operations—above all modular inversion, below—practical at the field sizes used in Halo~2 and Orchard.
\end{remark}

\subsection{The Bézout identity and the extended Euclidean algorithm}\label{intmod:sub:bezout}

The Euclidean algorithm computes the gcd, but the back-substitution of its equations yields something more powerful: the gcd is an \emph{integer linear combination} of the inputs. This is among the most useful facts in elementary number theory.

\begin{theorem}[Bézout's identity]\label{intmod:thm:bezout}
Let \(a,b\in\Z\), not both zero, and let \(d=\gcd(a,b)\). Then there exist integers \(u,v\in\Z\) with
\[
au + bv = d.
\]
Moreover \(d\) is the \emph{smallest} positive integer expressible in the form \(ax+by\) with \(x,y\in\Z\), and the set of all such combinations is exactly the set of multiples of \(d\):
\[
\{\, ax+by : x,y\in\Z \,\} = d\Z = \{\, kd : k\in\Z \,\}.
\]
\end{theorem}

\begin{proof}
Let \(I=\{\,ax+by:x,y\in\Z\,\}\). This set contains \(a\) and \(b\) and is closed under addition and under multiplication by any integer; concretely, if \(s=ax_1+by_1\) and \(t=ax_2+by_2\) lie in \(I\), then for any \(m,n\in\Z\) one has \(ms+nt=a(mx_1+nx_2)+b(my_1+ny_2)\in I\). Since \(a,b\) are not both zero, \(I\) contains a positive element (one of \(\pm a,\pm b\)), so the set \(I\cap\N_{>0}\) of positive elements of \(I\) is nonempty. By well-ordering it has a least element \(g=au+bv>0\).

We claim that \(g\) divides every element of \(I\). Indeed, let \(c\in I\) and write \(c=gq+r\) with \(0\le r<g\) by the division algorithm. Then \(r=c-gq\in I\) because \(I\) is closed under the operations just described. By minimality of \(g\) among positive elements of \(I\), and \(0\le r<g\), necessarily \(r=0\). Hence \(g\mid c\) for every \(c\in I\); in particular \(g\mid a\) and \(g\mid b\), so \(g\) is a common divisor of \(a\) and \(b\). Conversely, any common divisor \(e\) of \(a\) and \(b\) divides \(au+bv=g\) by Proposition~\ref{intmod:prop:divprops}(3). Therefore \(g\) is a common divisor that is divisible by every common divisor; being positive, it is in particular the \emph{greatest} common divisor, so \(g=d\). The displayed equation \(au+bv=d\) follows, \(d\) is the least positive element of \(I\), and \(I=g\Z=d\Z\) because every element of \(I\) is a multiple of \(g\), as shown, and conversely every multiple \(kg=a(ku)+b(kv)\) lies in \(I\).
\end{proof}

\begin{remark}\label{intmod:rem:strongestgcd}
The proof yields a characterisation of the gcd that is independent of the order \(\le\): \(d=\gcd(a,b)\) is, up to sign, the unique common divisor of \(a\) and \(b\) that is divisible by every common divisor. This is the definition that generalises to arbitrary commutative rings, where ``largest'' may make no sense but ``divisible by every common divisor'' still does. We call the subsets \(I\subseteq\Z\) closed under addition and under multiplication by arbitrary integers \emph{ideals}; the content of Bézout's identity is that every ideal of \(\Z\) has the form \(d\Z\) for a single generator \(d\), i.e.\ \(\Z\) is a \emph{principal ideal domain}. Ideals recur in the construction of quotient rings.
\end{remark}

The Bézout coefficients \(u,v\) are not unique. If \(au+bv=d\), then for every \(k\in\Z\),
\[
a\Bigl(u+k\tfrac{b}{d}\Bigr)+b\Bigl(v-k\tfrac{a}{d}\Bigr)=d,
\]
and these are all the solutions. To compute a particular pair \((u,v)\) efficiently, we augment the Euclidean algorithm so that it carries each remainder as an explicit combination of \(a\) and \(b\).

\begin{theorem}[Extended Euclidean algorithm]\label{intmod:thm:exteuclid}
There is an algorithm that, on input \(a,b\in\Z\) not both zero, outputs \(d=\gcd(a,b)\) together with integers \(u,v\) satisfying \(au+bv=d\), using the same number of division steps as the Euclidean algorithm.
\end{theorem}

\begin{proof}
Alongside the remainder sequence of Theorem~\ref{intmod:thm:euclid}, maintain two auxiliary sequences \((s_i)\) and \((t_i)\) with the invariant
\[
r_i = a\,s_i + b\,t_i \qquad\text{for all } i. \tag{\(\ast\)}\label{intmod:eq:eeainv}
\]
Initialise (taking \(a,b\ge 0\) for clarity; we restore signs at the end)
\[
r_0=a,\ s_0=1,\ t_0=0; \qquad r_1=b,\ s_1=0,\ t_1=1,
\]
so \eqref{intmod:eq:eeainv} holds for \(i=0,1\). At step \(i\ge 1\), having \(r_{i-1}=r_i q_i + r_{i+1}\) from the division algorithm, set
\[
s_{i+1}=s_{i-1}-q_i s_i, \qquad t_{i+1}=t_{i-1}-q_i t_i.
\]
Then
\[
a\,s_{i+1}+b\,t_{i+1} = (a s_{i-1}+b t_{i-1}) - q_i\,(a s_i + b t_i) = r_{i-1}-q_i r_i = r_{i+1},
\]
so induction preserves \eqref{intmod:eq:eeainv}. When the algorithm terminates with \(r_{n+1}=0\) and \(r_n=d=\gcd(a,b)\), the invariant for \(i=n\) reads \(d=a s_n + b t_n\); output \(u=s_n\), \(v=t_n\). Termination and the value of \(d\) are exactly as in Theorem~\ref{intmod:thm:euclid}, and each step does only a constant amount of extra work.
\end{proof}

In practice it is convenient to lay the computation out as a table with columns \((r_i,q_i,s_i,t_i)\), computing each new row from the previous two.

\begin{example}\label{intmod:exa:exteuclid}
Run the extended algorithm on \(a=1071\), \(b=462\) (continuing Example~\ref{intmod:exa:euclid}). The quotients are \(q_1=2,q_2=3,q_3=7\).
\[
\begin{array}{c|c|c|c}
i & r_i & s_i & t_i \\\hline
0 & 1071 & 1 & 0 \\
1 & 462 & 0 & 1 \\
2 & 147 & 1 & -2 \\
3 & 21 & -3 & 7 \\
4 & 0 & 22 & -51
\end{array}
\]
For instance, row \(2\) is row \(0\) minus \(q_1=2\) times row \(1\): \(s_2=1-2\cdot0=1\), \(t_2=0-2\cdot1=-2\), and indeed \(1071\cdot1+462\cdot(-2)=1071-924=147\). The last nonzero remainder is \(r_3=21=\gcd(1071,462)\), with
\[
1071\cdot(-3) + 462\cdot 7 = -3213 + 3234 = 21.
\]
Thus \((u,v)=(-3,7)\). The final row \((s_4,t_4)=(22,-51)=(b/d,-a/d)\) up to sign records the homogeneous solution: \(1071\cdot 22 + 462\cdot(-51)=0\).
\end{example}

A first, decisive application of Bézout's identity is Euclid's lemma, the property of primes that ultimately forces unique factorisation.

\begin{proposition}[Euclid's lemma]\label{intmod:prop:euclidlemma}
Let \(a,b,c\in\Z\).
\begin{enumerate}
\item If \(c\mid ab\) and \(\gcd(c,a)=1\), then \(c\mid b\).
\item If \(\gcd(a,b)=1\), \(a\mid n\), and \(b\mid n\), then \(ab\mid n\).
\item If \(\gcd(a,b)=1\) and \(\gcd(a,c)=1\), then \(\gcd(a,bc)=1\).
\end{enumerate}
\end{proposition}

\begin{proof}
(1) By Bézout there are \(u,v\) with \(cu+av=1\). Multiply by \(b\): \(b=cub+avb=c(ub)+(ab)v\). Since \(c\mid ab\), \(c\) divides both terms on the right, hence \(c\mid b\).

(2) Write \(n=ak\). Since \(b\mid n=ak\) and \(\gcd(b,a)=1\), part (1) gives \(b\mid k\), say \(k=bm\). Then \(n=a(bm)=(ab)m\), so \(ab\mid n\).

(3) By Bézout, \(ax_1+by_1=1\) and \(ax_2+cy_2=1\) for suitable integers. Multiplying,
\[
1=(ax_1+by_1)(ax_2+cy_2)=a\underbrace{(ax_1x_2+x_1cy_2+by_1x_2)}_{\in\Z}+bc\,(y_1y_2),
\]
which exhibits \(1\) as an integer combination of \(a\) and \(bc\). Hence \(\gcd(a,bc)\) divides \(1\), so it equals \(1\).
\end{proof}
\subsection{Primes and the fundamental theorem of arithmetic}\label{intmod:sub:primes}

\begin{definition}\label{intmod:def:prime}
An integer \(p>1\) is \emph{prime} if its only positive divisors are \(1\) and \(p\). An integer \(n>1\) that is not prime is \emph{composite}; thus \(n\) is composite iff \(n=ab\) for some integers \(1<a,b<n\). The integer \(1\) is a \emph{unit} and is by convention neither prime nor composite.
\end{definition}

The exclusion of \(1\) from the primes is not arbitrary: it is exactly what makes factorisation into primes unique, since otherwise one could insert arbitrarily many factors of \(1\). The decisive property of primes is the prime case of Euclid's lemma.

\begin{lemma}[Euclid's lemma for primes]\label{intmod:lem:primedivides}
Let \(p\) be prime and \(a,b\in\Z\). If \(p\mid ab\) then \(p\mid a\) or \(p\mid b\). More generally, if \(p\mid a_1a_2\cdots a_k\) then \(p\mid a_i\) for some \(i\).
\end{lemma}

\begin{proof}
The only positive divisors of \(p\) are \(1\) and \(p\), so \(\gcd(p,a)\in\{1,p\}\). If \(\gcd(p,a)=p\) then \(p\mid a\) and the claim holds; otherwise \(\gcd(p,a)=1\), and Proposition~\ref{intmod:prop:euclidlemma}(1) gives \(p\mid b\). The general statement follows by induction on \(k\): from \(p\mid a_1(a_2\cdots a_k)\) it follows that \(p\mid a_1\) or \(p\mid a_2\cdots a_k\), and the inductive hypothesis handles the latter.
\end{proof}

\begin{theorem}[Fundamental theorem of arithmetic]\label{intmod:thm:fta}
Every integer \(n>1\) can be written as a product of primes,
\[
n = p_1 p_2 \cdots p_k,
\]
and this factorisation is unique up to the order of the factors. Equivalently, there is a unique expression
\[
n = p_1^{e_1} p_2^{e_2}\cdots p_r^{e_r}, \qquad p_1<p_2<\cdots<p_r \text{ prime},\ e_i\ge 1.
\]
\end{theorem}

\begin{proof}
\emph{Existence.} Suppose, for contradiction, that some integer \(>1\) is not a product of primes, and let \(n\) be the least such (well-ordering). Then \(n\) is not prime (a prime is trivially a product of one prime), so \(n=ab\) with \(1<a,b<n\). By minimality, each of \(a,b\) is a product of primes; concatenating these factorisations expresses \(n\) as a product of primes, a contradiction. Hence every \(n>1\) factors into primes.

\emph{Uniqueness.} Suppose
\[
p_1 p_2 \cdots p_k = q_1 q_2 \cdots q_\ell
\]
are two factorisations of the same integer into primes. The argument proceeds by induction on \(k\). If \(k=0\) the product is the empty product \(1\), forcing \(\ell=0\) as well (a nonempty product of primes is \(>1\)). For the inductive step, \(p_1\) divides the right-hand product, so by Lemma~\ref{intmod:lem:primedivides} \(p_1\mid q_j\) for some \(j\). Since \(q_j\) is prime, its only divisor \(>1\) is itself, so \(p_1=q_j\). Cancelling this common factor (legitimate as \(\Z\) has no zero divisors) yields
\[
p_2\cdots p_k = q_1\cdots \widehat{q_j}\cdots q_\ell,
\]
where the hat denotes omission. By the inductive hypothesis these two factorisations agree up to order, with \(k-1=\ell-1\); restoring \(p_1=q_j\) shows the original factorisations agree up to order. Collecting equal primes into prime powers gives the canonical form, whose uniqueness is immediate from the uniqueness of the multiset of primes.
\end{proof}

\begin{remark}\label{intmod:rem:ftaroles}
Both halves of the theorem are needed and have different characters. Existence is a statement about the order structure (well-ordering) and would hold in any reasonable factorisation; uniqueness rests squarely on Euclid's lemma, hence on Bézout, hence on the division algorithm. There are number systems with a division-like structure failing Euclid's lemma in which factorisation into irreducibles is \emph{not} unique; the integers avoid this pathology precisely because they form a Euclidean domain.
\end{remark}

\subsection{Congruences and the ring \texorpdfstring{\(\Z/n\Z\)}{Z/nZ}}\label{intmod:sub:congruences}

The development now passes from divisibility to its quotient, the arithmetic of remainders. Fix an integer \(n\ge 1\), the \emph{modulus}.

\begin{definition}\label{intmod:def:congruent}
Let \(n\ge 1\) and \(a,b\in\Z\). The integer \(a\) is \emph{congruent to} \(b\) \emph{modulo} \(n\), written
\[
a\equiv b \pmod n,
\]
if \(n\mid (a-b)\), equivalently if \(a\) and \(b\) leave the same remainder upon division by \(n\).
\end{definition}

The two formulations agree: if \(a=nq_1+r_1\) and \(b=nq_2+r_2\) with \(0\le r_1,r_2<n\), then \(a-b=n(q_1-q_2)+(r_1-r_2)\), and \(n\mid(a-b)\) iff \(n\mid(r_1-r_2)\) iff \(r_1=r_2\) (since \(|r_1-r_2|<n\)).

\begin{proposition}\label{intmod:prop:congequiv}
Fix \(n\ge 1\). Congruence modulo \(n\) is an equivalence relation on \(\Z\) that is moreover compatible with addition and multiplication: if \(a\equiv a'\) and \(b\equiv b'\pmod n\), then
\[
a+b\equiv a'+b' \pmod n, \qquad ab\equiv a'b' \pmod n.
\]
Consequently \(a^k\equiv (a')^k\pmod n\) for all \(k\ge 0\), and one may substitute congruent values freely in any polynomial expression with integer coefficients.
\end{proposition}

\begin{proof}
Reflexivity (\(n\mid 0\)), symmetry (\(n\mid(a-b)\Rightarrow n\mid(b-a)\)), and transitivity (\(n\mid(a-b)\) and \(n\mid(b-c)\) give \(n\mid(a-c)\) by Proposition~\ref{intmod:prop:divprops}(3)) are immediate. For compatibility, write \(a-a'=n s\) and \(b-b'=n t\). Then
\[
(a+b)-(a'+b')=n(s+t),
\]
and
\[
ab-a'b' = ab - a'b + a'b - a'b' = (a-a')b + a'(b-b') = n(sb + a' t),
\]
so both \(n\mid\bigl((a+b)-(a'+b')\bigr)\) and \(n\mid(ab-a'b')\). The statement about powers and polynomials follows by induction from the multiplicative case.
\end{proof}

Because congruence is an equivalence relation, it partitions \(\Z\) into disjoint \emph{congruence classes} (or \emph{residue classes}).

\begin{definition}\label{intmod:def:residueclass}
The \emph{congruence class} of \(a\) modulo \(n\) is
\[
[a]_n = \{\, b\in\Z : b\equiv a \pmod n \,\} = \{\, a+kn : k\in\Z \,\} = a+n\Z.
\]
We write the set of all congruence classes as
\[
\Z/n\Z = \{\, [a]_n : a\in\Z \,\}.
\]
\end{definition}

By the division algorithm every integer is congruent modulo \(n\) to exactly one of \(0,1,\dots,n-1\), namely its remainder; hence the \(n\) classes \([0]_n,[1]_n,\dots,[n-1]_n\) are distinct and exhaust \(\Z/n\Z\). The set \(\{0,1,\dots,n-1\}\) is the standard \emph{system of representatives}, called the \emph{least nonnegative residues}.

\begin{theorem}\label{intmod:thm:znring}
For each \(n\ge 1\), the operations
\[
[a]_n + [b]_n := [a+b]_n, \qquad [a]_n\cdot[b]_n := [ab]_n
\]
are well defined and make \(\Z/n\Z\) a commutative ring with additive identity \([0]_n\) and multiplicative identity \([1]_n\). It has exactly \(n\) elements.
\end{theorem}

\begin{proof}
Well-definedness is precisely the compatibility of Proposition~\ref{intmod:prop:congequiv}: the sum and product of classes do not depend on the chosen representatives. Each ring axiom for \(\Z/n\Z\) (associativity and commutativity of \(+\) and \(\cdot\), distributivity, existence of \(0\) and \(1\), additive inverses \(-[a]_n=[-a]_n\)) is inherited directly from the corresponding axiom in \(\Z\) by passing to classes; for instance associativity of multiplication reads \(([a][b])[c]=[abc]=[a]([b][c])\). The count of exactly \(n\) classes appeared above.
\end{proof}

\begin{remark}\label{intmod:rem:quotientring}
Structurally, \(\Z/n\Z\) is the quotient of the ring \(\Z\) by the ideal \(n\Z=\{kn:k\in\Z\}\) (Remark~\ref{intmod:rem:strongestgcd}). The map \(\pi\colon\Z\to\Z/n\Z\), \(\pi(a)=[a]_n\), is a surjective ring homomorphism with kernel \(n\Z\); this is the prototype of the quotient constructions used throughout the algebra to come. One often drops the brackets and the subscript, writing \(a\) for \([a]_n\) when the modulus is clear, and uses ``\(\Z/n\Z\)'' or ``\(\Z_n\)'' interchangeably; \(\Z_n\) denotes this ring (not the \(n\)-adic integers, which do not appear here).
\end{remark}

A basic structural question is when \(\Z/n\Z\) has no zero divisors—equivalently, when it is a field.

\begin{proposition}\label{intmod:prop:zndomain}
For \(n\ge 2\), the ring \(\Z/n\Z\) has no zero divisors if and only if \(n\) is prime; in that case, as we show below, it is a field, written \(\Fp\) when \(n=p\).
\end{proposition}

\begin{proof}
If \(n=ab\) with \(1<a,b<n\) is composite, then \([a]_n\neq[0]_n\neq[b]_n\) but \([a]_n[b]_n=[n]_n=[0]_n\), so there are zero divisors. Conversely, if \(n=p\) is prime and \([a]_p[b]_p=[0]_p\), then \(p\mid ab\), so by Lemma~\ref{intmod:lem:primedivides} \(p\mid a\) or \(p\mid b\), i.e.\ \([a]_p=[0]_p\) or \([b]_p=[0]_p\). That a finite commutative ring without zero divisors and with \(1\neq 0\) is a field follows from the discussion of units below (Corollary~\ref{intmod:cor:fpfield}).
\end{proof}

\subsection{Units modulo \texorpdfstring{\(n\)}{n} and modular inverses}\label{intmod:sub:units}

In any ring, the elements with multiplicative inverses are the most important. In \(\Z/n\Z\) the Euclidean machinery characterises them and computes them.

\begin{definition}\label{intmod:def:unit}
An element \([a]_n\in\Z/n\Z\) is a \emph{unit} (or is \emph{invertible}) if there exists \([x]_n\) with \([a]_n[x]_n=[1]_n\), i.e.\ \(ax\equiv 1\pmod n\). Such an \([x]_n\) is the \emph{(modular) inverse} of \([a]_n\), written \([a]_n^{-1}\). We write the set of units as
\[
(\Z/n\Z)^\times = \{\, [a]_n : [a]_n \text{ is a unit} \,\}.
\]
\end{definition}

\begin{theorem}\label{intmod:thm:unitiff}
Let \(n\ge 1\) and \(a\in\Z\). The class \([a]_n\) is a unit in \(\Z/n\Z\) if and only if \(\gcd(a,n)=1\). When \(\gcd(a,n)=1\), the inverse is unique, and the extended Euclidean algorithm computes a representative: if \(au+nv=1\), then \([a]_n^{-1}=[u]_n\).
\end{theorem}

\begin{proof}
Suppose \(\gcd(a,n)=1\). By Bézout (Theorem~\ref{intmod:thm:bezout}) there are \(u,v\) with \(au+nv=1\). Reducing modulo \(n\), \(au\equiv 1\pmod n\), so \([u]_n\) is an inverse of \([a]_n\). The extended Euclidean algorithm (Theorem~\ref{intmod:thm:exteuclid}) produces such a \(u\) explicitly.

Conversely, if \([a]_n\) is a unit, say \(ax\equiv 1\pmod n\), then \(ax-1=nk\) for some \(k\), i.e.\ \(ax-nk=1\). Any common divisor of \(a\) and \(n\) divides the left side, hence divides \(1\); thus \(\gcd(a,n)=1\).

Uniqueness: if \(ax\equiv 1\) and \(ax'\equiv 1\pmod n\), then \(x\equiv x(ax')\equiv(xa)x'\equiv x'\pmod n\), so \([x]_n=[x']_n\).
\end{proof}

\begin{corollary}\label{intmod:cor:fpfield}
If \(p\) is prime then every nonzero class in \(\Z/p\Z\) is a unit, so \(\Z/p\Z=\Fp\) is a field. More generally, \(\Z/n\Z\) is a field if and only if \(n\) is prime.
\end{corollary}

\begin{proof}
If \(p\) is prime and \([a]_p\neq[0]_p\) then \(p\nmid a\), so \(\gcd(a,p)=1\) (the only positive divisors of \(p\) are \(1,p\)), and Theorem~\ref{intmod:thm:unitiff} makes \([a]_p\) a unit. Thus every nonzero element is invertible and \(\Fp\) is a field. If \(n\) is composite, \(n=ab\) with \(1<a,b<n\), then \([a]_n\neq[0]_n\) yet \(\gcd(a,n)=a>1\), so \([a]_n\) is not a unit; a ring with a nonzero non-unit (equivalently, with zero divisors, by Proposition~\ref{intmod:prop:zndomain}) is not a field. The case \(n=1\) gives the zero ring, not a field by convention.
\end{proof}

\begin{example}\label{intmod:exa:modinverse}
Consider the computation of \([17]_{43}^{-1}\). The extended Euclidean algorithm on \((43,17)\):
\[
\begin{aligned}
43 &= 2\cdot 17 + 9, \\
17 &= 1\cdot 9 + 8, \\
9 &= 1\cdot 8 + 1, \\
8 &= 8\cdot 1 + 0,
\end{aligned}
\]
so \(\gcd(43,17)=1\). Back-substituting:
\[
1 = 9 - 8 = 9 - (17-9) = 2\cdot 9 - 17 = 2(43-2\cdot 17) - 17 = 2\cdot 43 - 5\cdot 17.
\]
Hence \(-5\cdot 17 \equiv 1\pmod{43}\), and \([17]_{43}^{-1}=[-5]_{43}=[38]_{43}\). Check: \(17\cdot 38 = 646 = 15\cdot 43 + 1\), so \(17\cdot 38\equiv 1\pmod{43}\).
\end{example}

The set of units forms a group under multiplication, a fact we use heavily in what follows.

\begin{proposition}\label{intmod:prop:unitgroup}
Under multiplication of classes, \((\Z/n\Z)^\times\) is an abelian group. Its order is denoted \(\varphi(n)\) (\emph{Euler's totient}): the number of integers in \(\{1,\dots,n\}\) coprime to \(n\).
\end{proposition}

\begin{proof}
Multiplication of classes is associative and commutative (Theorem~\ref{intmod:thm:znring}) with identity \([1]_n\). It is closed on units: if \([a]_n,[b]_n\) are units with inverses \([a]_n^{-1},[b]_n^{-1}\), then \([ab]_n\) has inverse \([b]_n^{-1}[a]_n^{-1}\), since \((ab)(b^{-1}a^{-1})\equiv 1\). Every unit has an inverse by definition, and that inverse is itself a unit (it is invertible, with inverse the original element). Hence \((\Z/n\Z)^\times\) is an abelian group. By Theorem~\ref{intmod:thm:unitiff} its elements are exactly the classes \([a]_n\) with \(\gcd(a,n)=1\), \(1\le a\le n\), of which there are \(\varphi(n)\) by definition.
\end{proof}

\begin{remark}\label{intmod:rem:totientcrypto}
The totient \(\varphi(n)\) governs the multiplicative order structure modulo \(n\). For prime \(p\), \(\varphi(p)=p-1\), so \(\Fp^\times\) has order \(p-1\); for a prime power, \(\varphi(p^k)=p^{k-1}(p-1)\). Lagrange's theorem (developed in the group-theory section) then gives Euler's theorem \(a^{\varphi(n)}\equiv 1\pmod n\) for \(\gcd(a,n)=1\), and in particular Fermat's little theorem \(a^{p-1}\equiv 1\pmod p\) for \(p\nmid a\). These facts underlie modular exponentiation, the cyclic structure of \(\Fp^\times\), and the choice of scalar field sizes in the elliptic-curve constructions of Halo~2 and Orchard. The Chinese Remainder Theorem yields the multiplicativity of \(\varphi\), which then computes \(\varphi(n)\) for general \(n\).
\end{remark}

\section{Groups}\label{sec:groups}

The notion of a \emph{group} is the foundational abstraction of modern algebra. A group axiomatises the bare minimum needed to discuss \emph{symmetry} and \emph{invertible composition}: a set together with one associative way of combining its elements, possessing an identity and inverses. This short list of axioms is rich enough to support a deep structure theory, and almost every object encountered later---the integers modulo a prime under addition, the nonzero residues under multiplication, the points of an elliptic curve, the symmetries of a finite field---is a group. This section develops group theory from first principles, proving the central structural results in full: Lagrange's theorem, the classification of cyclic groups, the quotient of an abelian group by a subgroup, and the basic theory of homomorphisms and isomorphisms.

\subsection{Definition and axioms}

Before we can discuss groups, we must make precise the meaning of ``combining'' two elements of a set.

\begin{definition}\label{groups:def-binop}
Let \(S\) be a set. A \emph{binary operation} on \(S\) is a function
\[
\ast \colon S \times S \to S, \qquad (a,b) \mapsto a \ast b .
\]
The element \(a \ast b\) denotes the image of the pair \((a,b)\). The defining feature is \emph{closure}: for all \(a,b \in S\), the element \(a \ast b\) again lies in \(S\).
\end{definition}

Closure is built into the very statement that \(\ast\) is a function \emph{into} \(S\); we name it explicitly nonetheless because, when we later examine a subset \(H \subseteq S\), the question of whether \(\ast\) restricts to a binary operation on \(H\) is exactly the question of whether \(H\) is closed under \(\ast\).

\begin{definition}\label{groups:def-group}
A \emph{group} is a pair \((G, \ast)\) consisting of a set \(G\) and a binary operation \(\ast\) on \(G\) satisfying the following three axioms.
\begin{enumerate}
\item[(G1)] \emph{Associativity.} For all \(a,b,c \in G\),
\[
(a \ast b) \ast c = a \ast (b \ast c).
\]
\item[(G2)] \emph{Identity element.} There exists an element \(e \in G\) such that for all \(a \in G\),
\[
e \ast a = a \ast e = a.
\]
\item[(G3)] \emph{Inverses.} For each \(a \in G\) there exists an element \(b \in G\) such that
\[
a \ast b = b \ast a = e,
\]
where \(e\) is an identity element as in (G2).
\end{enumerate}
A set with a binary operation satisfying only (G1) is called a \emph{semigroup}; if it also satisfies (G2) it is a \emph{monoid}. Thus a group is a monoid in which every element is invertible.
\end{definition}

We frequently use the phrase ``the group \(G\)'', with the operation understood. The number of elements of \(G\), written \(|G|\), is called the \emph{order} of the group; if \(|G|\) is finite \(G\) is a \emph{finite group}, otherwise an \emph{infinite group}.

The axioms assert the \emph{existence} of an identity and of inverses, but a priori not their uniqueness. Uniqueness is a theorem, and a useful one, since it licenses the notation \(e\) and \(a^{-1}\).

\begin{proposition}\label{groups:prop-unique-id-inv}
Let \((G,\ast)\) be a group.
\begin{enumerate}
\item The identity element is unique.
\item Each \(a \in G\) has a unique inverse, which is denoted \(a^{-1}\).
\item \emph{(Cancellation.)} For all \(a,x,y \in G\), if \(a \ast x = a \ast y\) then \(x = y\), and if \(x \ast a = y \ast a\) then \(x = y\).
\item For all \(a,b \in G\), \((a^{-1})^{-1} = a\) and \((a \ast b)^{-1} = b^{-1} \ast a^{-1}\).
\end{enumerate}
\end{proposition}

\begin{proof}
(1) Suppose \(e\) and \(e'\) both satisfy (G2). Then, using that \(e'\) is an identity, \(e = e \ast e'\); and using that \(e\) is an identity, \(e \ast e' = e'\). Hence \(e = e'\).

(2) Suppose \(b\) and \(b'\) are both inverses of \(a\). Then
\[
b = b \ast e = b \ast (a \ast b') = (b \ast a) \ast b' = e \ast b' = b',
\]
using (G2), the inverse property of \(b'\), associativity (G1), and the inverse property of \(b\). Thus the inverse is unique.

(3) Suppose \(a \ast x = a \ast y\). Applying \(a^{-1}\) on the left,
\[
x = e \ast x = (a^{-1} \ast a) \ast x = a^{-1} \ast (a \ast x) = a^{-1} \ast (a \ast y) = (a^{-1} \ast a) \ast y = e \ast y = y.
\]
The right-cancellation law is symmetric.

(4) Since \(a \ast a^{-1} = a^{-1} \ast a = e\), the element \(a\) is an inverse of \(a^{-1}\); by uniqueness \((a^{-1})^{-1} = a\). For the second identity, compute
\[
(a \ast b) \ast (b^{-1} \ast a^{-1}) = a \ast (b \ast b^{-1}) \ast a^{-1} = a \ast e \ast a^{-1} = a \ast a^{-1} = e,
\]
and symmetrically \((b^{-1} \ast a^{-1}) \ast (a \ast b) = e\). By uniqueness of inverses, \((a\ast b)^{-1} = b^{-1} \ast a^{-1}\).
\end{proof}

\paragraph{Notation.} Two notational conventions dominate. In \emph{multiplicative} notation one writes the operation as juxtaposition \(ab\) (or \(a \cdot b\)), the identity as \(1\), and the inverse as \(a^{-1}\); powers are \(a^n = \underbrace{a \cdots a}_{n}\) for \(n \ge 1\), \(a^0 = 1\), and \(a^{-n} = (a^{-1})^n\). In \emph{additive} notation, reserved almost exclusively for commutative groups, one writes \(a + b\), the identity as \(0\), the inverse as \(-a\), and \(na\) for the \(n\)-fold sum. The usual exponent laws \(a^m a^n = a^{m+n}\) and \((a^m)^n = a^{mn}\) hold for all integers \(m,n\) (proved by induction together with Proposition~\ref{groups:prop-unique-id-inv}). Henceforth \(ab\) usually denotes \(a \ast b\) and we drop the symbol \(\ast\) unless clarity demands it.
\subsection{Abelian groups}

\begin{definition}\label{groups:def-abelian}
A group \((G,\ast)\) is \emph{abelian} (or \emph{commutative}) if in addition
\[
a \ast b = b \ast a \qquad \text{for all } a,b \in G.
\]
A group that is not abelian we call \emph{non-abelian}.
\end{definition}

The term honours Niels Henrik Abel. Commutativity is a strong simplifying assumption: in an abelian group one may freely reorder products, the formula \((ab)^{-1} = b^{-1}a^{-1} = a^{-1}b^{-1}\) loses its order-sensitivity, and---as we show below---the structure theory becomes especially clean. Convention reserves additive notation for abelian groups precisely because \(+\) connotes commutativity.

\subsection{Examples of groups}

We now exhibit a generous supply of groups. The reader should check each example against the three axioms; we comment only on the points that are not immediate.

\begin{example}[Integers under addition]\label{groups:exa-Z}
The set \(\Z\) of integers with the operation \(+\) is an abelian group. Addition is associative and commutative, \(0\) is the identity, and the inverse of \(n\) is \(-n\). The same holds for \((\Q,+)\), \((\R,+)\), and \((\C,+)\). All of these are infinite. Note \((\Z, \cdot)\) is \emph{not} a group: \(2\) has no multiplicative inverse in \(\Z\).
\end{example}

\begin{example}[Nonzero elements of a field under multiplication]\label{groups:exa-field-mult}
For \(\Q\), \(\R\), or \(\C\), the set of \emph{nonzero} elements, denoted \(\Q^\times\), \(\R^\times\), \(\C^\times\), forms an abelian group under multiplication, with identity \(1\) and inverse \(a^{-1} = 1/a\). We must exclude the element \(0\) because \(0\) has no multiplicative inverse. More generally, for any field \(\F\) the nonzero elements \(\F^\times = \F \setminus \{0\}\) form an abelian group under multiplication; this fact is part of the very definition of a field.
\end{example}

\begin{example}[Integers and units modulo \(n\)]\label{groups:exa-Zn}
\((\Z/n\Z, +)\) is an abelian group of order \(n\) (Theorem~\ref{intmod:thm:znring}), and the invertible classes form the abelian group \((\Z/n\Z)^\times\) of order \(\varphi(n)\) under multiplication (\Cref{intmod:thm:unitiff,intmod:prop:unitgroup}). When \(n = p\) is prime, every nonzero class is a unit, so \((\Z/p\Z)^\times = \Fp^\times\) has order \(p-1\); this is Example~\ref{groups:exa-field-mult} for the finite field \(\Fp\).
\end{example}

\begin{example}[Trivial group and small groups]\label{groups:exa-trivial}
The one-element set \(\{e\}\) with the only possible operation is a group, the \emph{trivial group}. The group \(\Z/2\Z = \{0,1\}\) is the smallest nontrivial group. The \emph{Klein four-group} \(V = \{e,a,b,c\}\), in which every nonidentity element squares to \(e\) and the product of any two distinct nonidentity elements is the third, is abelian of order \(4\) and is \emph{not} isomorphic to \(\Z/4\Z\) (it has no element of order \(4\)).
\end{example}

\subsection{Order of an element}

\begin{definition}\label{groups:def-order-elt}
Let \(G\) be a group with identity \(e\) and let \(a \in G\). The \emph{order} of \(a\), written \(\ord(a)\), is the least positive integer \(n\) such that \(a^n = e\), if such an \(n\) exists; in that case \(a\) has \emph{finite order}. If no positive power of \(a\) equals \(e\), then \(a\) has \emph{infinite order} and we write \(\ord(a) = \infty\).
\end{definition}

The theory of cyclic groups below reconciles the two uses of the word ``order''---for an element and for a group: \(\ord(a)\) equals the order of the smallest subgroup containing \(a\).

\begin{proposition}\label{groups:prop-order-divides}
Let \(a \in G\) have finite order \(d = \ord(a)\). Then for any integer \(m\),
\[
a^m = e \iff d \mid m.
\]
Consequently \(a^m = a^k\) iff \(m \equiv k \pmod d\), and the distinct powers of \(a\) are exactly \(e = a^0, a^1, \dots, a^{d-1}\).
\end{proposition}

\begin{proof}
If \(d \mid m\), write \(m = dq\); then \(a^m = (a^d)^q = e^q = e\). Conversely, suppose \(a^m = e\). By the division algorithm write \(m = dq + r\) with \(0 \le r < d\). Then
\[
e = a^m = (a^d)^q a^r = a^r.
\]
Since \(0 \le r < d\) and \(d\) is the \emph{least} positive integer with \(a^d = e\), it follows that \(r = 0\), i.e.\ \(d \mid m\). The remaining claims follow: \(a^m = a^k \iff a^{m-k} = e \iff d \mid (m-k)\), and the \(d\) elements \(a^0,\dots,a^{d-1}\) are pairwise distinct because their exponents are incongruent modulo \(d\).
\end{proof}

\begin{example}\label{groups:exa-orders}
In \((\Z/n\Z,+)\) the order of the element \(1\) is \(n\), since \(\underbrace{1+\cdots+1}_{k} = k\) is \(0\) iff \(n \mid k\). The order of an element \(a\) is \(n/\gcd(a,n)\) (Corollary~\ref{groups:cor-order-power}). In \(\C^\times\), the element \(i\) has order \(4\): \(i^1=i, i^2=-1, i^3=-i, i^4=1\). In \(\Z\) under addition, every nonzero element has infinite order.
\end{example}

\subsection{Subgroups and the subgroup test}

\begin{definition}\label{groups:def-subgroup}
A subset \(H\) of a group \((G,\ast)\) is a \emph{subgroup}, written \(H \le G\), if \(H\) is itself a group under the restriction of \(\ast\) to \(H \times H\). Explicitly: \(H\) is nonempty, closed under \(\ast\) (so \(\ast\) restricts to a binary operation on \(H\)), contains the identity \(e\) of \(G\), and contains \(a^{-1}\) whenever it contains \(a\).
\end{definition}

Every group \(G\) has at least two subgroups: the \emph{trivial subgroup} \(\{e\}\) and \(G\) itself. We call a subgroup other than \(G\) \emph{proper}. Note that the identity of a subgroup necessarily coincides with the identity of \(G\) (apply cancellation in \(G\) to \(e_H \ast e_H = e_H = e_H \ast e\)), and inverses computed in \(H\) agree with those in \(G\); thus the definition is not as redundant as it may appear.

Verifying all the group axioms for a candidate subgroup is wasteful, since associativity is inherited for free. The following test reduces the verification to a single condition.

\begin{proposition}[One-step subgroup test]\label{groups:prop-subgroup-test}
Let \(G\) be a group and \(H \subseteq G\). Then \(H \le G\) if and only if
\begin{enumerate}
\item \(H \neq \varnothing\), and
\item for all \(a,b \in H\), the element \(a b^{-1} \in H\).
\end{enumerate}
\end{proposition}

\begin{proof}
($\Rightarrow$) If \(H\) is a subgroup it is nonempty and, given \(a,b \in H\), it contains \(b^{-1}\) and hence \(ab^{-1}\).

($\Leftarrow$) Assume (1) and (2). Pick any \(a \in H\) (possible by (1)). Taking \(b = a\) in (2) gives \(a a^{-1} = e \in H\). Taking \(a = e\) and \(b\) arbitrary in \(H\) gives \(e b^{-1} = b^{-1} \in H\), so \(H\) is closed under inverses. Finally, for \(a,b \in H\) we now know \(b^{-1} \in H\), so applying (2) to the pair \((a, b^{-1})\) yields \(a (b^{-1})^{-1} = ab \in H\); thus \(H\) is closed under the operation. Associativity is inherited from \(G\). Hence \(H\) is a group, i.e.\ \(H \le G\).
\end{proof}

\begin{remark}\label{groups:rem-finite-test}
For a \emph{finite} nonempty subset \(H\) closed under the operation, closure alone forces \(H \le G\): given \(a \in H\), the powers \(a, a^2, a^3, \dots\) all lie in \(H\) and cannot all be distinct (as \(H\) is finite), so \(a^i = a^j\) for some \(i<j\), giving \(a^{j-i} = e \in H\) and \(a^{-1} = a^{j-i-1} \in H\). Thus for finite subsets, closure under multiplication is the only condition one must verify.
\end{remark}

\begin{example}\label{groups:exa-subgroups}
The even integers \(2\Z = \{2k : k \in \Z\}\) form a subgroup of \((\Z,+)\); more generally \(n\Z \le \Z\) for every \(n \ge 0\), and these are \emph{all} the subgroups of \(\Z\) (Proposition~\ref{groups:prop-subgroups-of-Z}).
\end{example}

\begin{proposition}\label{groups:prop-intersection}
The intersection of any family \(\{H_i\}_{i \in I}\) of subgroups of \(G\) is a subgroup of \(G\).
\end{proposition}

\begin{proof}
Let \(H = \bigcap_i H_i\). Each \(H_i\) contains \(e\), so \(e \in H\) and \(H \neq \varnothing\). If \(a,b \in H\) then \(a,b \in H_i\) for every \(i\), so \(ab^{-1} \in H_i\) for every \(i\) by the subgroup test, whence \(ab^{-1} \in H\). By Proposition~\ref{groups:prop-subgroup-test}, \(H \le G\).
\end{proof}

This lets us define the subgroup ``generated'' by an arbitrary subset as the smallest subgroup containing it.

\begin{definition}\label{groups:def-generated}
Let \(S \subseteq G\). The \emph{subgroup generated by \(S\)}, written \(\langle S \rangle\), is the intersection of all subgroups of \(G\) that contain \(S\). By Proposition~\ref{groups:prop-intersection} it is a subgroup, and it is the smallest subgroup containing \(S\). Concretely, \(\langle S \rangle\) consists of all finite products \(s_1^{\epsilon_1} \cdots s_k^{\epsilon_k}\) with \(s_i \in S\) and \(\epsilon_i \in \{+1,-1\}\) (the empty product being \(e\)).
\end{definition}

\subsection{Cyclic groups, generators, and their classification}

The simplest groups are those generated by a single element.

\begin{definition}\label{groups:def-cyclic}
A group \(G\) is \emph{cyclic} if there exists \(g \in G\) with \(G = \langle g \rangle\), where
\[
\langle g \rangle = \{ g^k : k \in \Z \}
\]
is the subgroup generated by \(g\). We call such a \(g\) a \emph{generator} of \(G\). Every cyclic group is abelian, since \(g^m g^n = g^{m+n} = g^n g^m\).
\end{definition}

The description \(\langle g \rangle = \{g^k : k \in \Z\}\) matches Definition~\ref{groups:def-generated}: the right-hand side is closed under products and inverses and contains \(g\), and any subgroup containing \(g\) must contain all its powers.

\begin{proposition}\label{groups:prop-cyclic-order}
Let \(G = \langle g \rangle\) be cyclic.
\begin{enumerate}
\item If \(g\) has infinite order, then the powers \(g^k\) (\(k \in \Z\)) are pairwise distinct, so \(G\) is infinite and \(|G| = \infty\).
\item If \(g\) has finite order \(d\), then \(G = \{e, g, g^2, \dots, g^{d-1}\}\) has exactly \(d\) elements, so \(|G| = \ord(g) = d\).
\end{enumerate}
In both cases \(|\langle g \rangle| = \ord(g)\).
\end{proposition}

\begin{proof}
(1) If \(g\) has infinite order and \(g^i = g^j\) with \(i \neq j\), then \(g^{i-j} = e\) with \(i - j \neq 0\), contradicting infinite order; so all powers are distinct. (2) is Proposition~\ref{groups:prop-order-divides}: the distinct powers are exactly \(g^0,\dots,g^{d-1}\).
\end{proof}

\begin{example}\label{groups:exa-cyclic}
\((\Z, +)\) is infinite cyclic, generated by \(1\) (and also by \(-1\)). \((\Z/n\Z, +)\) is cyclic of order \(n\), generated by \(1\). The group \((\Z/8\Z)^\times = \{1,3,5,7\}\) is \emph{not} cyclic: every element squares to \(1\), so no element has order \(4\). By contrast \((\Z/5\Z)^\times = \{1,2,3,4\}\) is cyclic, generated by \(2\): \(2^1=2,\,2^2=4,\,2^3=3,\,2^4=1\).
\end{example}

We now determine \emph{which} powers of a generator are themselves generators, and classify the subgroups of a cyclic group.

\begin{corollary}\label{groups:cor-order-power}
Let \(g\) have finite order \(n\), and let \(k \in \Z\). Then
\[
\ord(g^k) = \frac{n}{\gcd(n,k)}.
\]
In particular \(g^k\) generates \(\langle g \rangle\) if and only if \(\gcd(n,k)=1\), so a cyclic group of order \(n\) has exactly \(\varphi(n)\) generators.
\end{corollary}

\begin{proof}
Let \(d = \gcd(n,k)\) and write \(k = d k'\), \(n = d n'\) with \(\gcd(n',k')=1\). Then \((g^k)^m = g^{km} = e\) iff \(n \mid km\), i.e.\ iff \(n' \mid k' m\); since \(\gcd(n',k')=1\), this holds iff \(n' \mid m\). The least positive such \(m\) is \(n' = n/d\), so \(\ord(g^k) = n/\gcd(n,k)\). Now \(g^k\) generates \(\langle g\rangle\) iff \(\ord(g^k) = n\) iff \(\gcd(n,k)=1\). The number of \(k\) with \(1 \le k \le n\) and \(\gcd(n,k)=1\) is \(\varphi(n)\).
\end{proof}

\begin{theorem}[Classification of subgroups of cyclic groups]\label{groups:thm-cyclic-subgroups}
Let \(G = \langle g\rangle\) be cyclic.
\begin{enumerate}
\item Every subgroup of \(G\) is cyclic.
\item If \(G\) is infinite, its subgroups are exactly \(\langle g^m\rangle\) for \(m \ge 0\); distinct \(m \ge 0\) give distinct subgroups, and \(\langle g^m\rangle \cong (\Z,+)\) for \(m \ge 1\).
\item If \(|G| = n\), then for each positive divisor \(d \mid n\) there is exactly one subgroup of order \(d\), namely \(\langle g^{n/d}\rangle\), and these are all the subgroups. Hence the subgroups of \(G\) are in bijection with the positive divisors of \(n\).
\end{enumerate}
\end{theorem}

\begin{proof}
(1) Let \(H \le G\). If \(H = \{e\}\) it is cyclic. Otherwise \(H\) contains some \(g^k\) with \(k \neq 0\); since \(H\) is closed under inverses it contains \(g^{-k}\) too, so \(H\) contains a positive power of \(g\). Let \(m\) be the \emph{least} positive integer with \(g^m \in H\). We claim that \(H = \langle g^m\rangle\). Clearly \(\langle g^m\rangle \subseteq H\). Conversely take any \(g^t \in H\); by the division algorithm \(t = mq + r\) with \(0 \le r < m\), so
\[
g^r = g^{t - mq} = g^t (g^m)^{-q} \in H .
\]
By minimality of \(m\) and \(0 \le r < m\), it follows that \(r = 0\), so \(g^t = (g^m)^q \in \langle g^m\rangle\). Hence \(H = \langle g^m\rangle\) is cyclic.

(2) If \(G\) is infinite, then \(g\) has infinite order, so the subgroups \(\langle g^m\rangle\) (\(m \ge 1\)) are all infinite cyclic, and \(\langle g^0\rangle = \{e\}\). They are distinct because \(\langle g^m\rangle\) has \(g^m\) as its least positive power but does not contain \(g^j\) for \(0<j<m\). Part (1) shows these are all subgroups.

(3) Suppose \(|G| = n\), so \(\ord(g) = n\). Let \(d \mid n\). The element \(g^{n/d}\) has order \(\ord(g^{n/d}) = n/\gcd(n, n/d) = n/(n/d) = d\) by Corollary~\ref{groups:cor-order-power}, so \(\langle g^{n/d}\rangle\) is a subgroup of order \(d\). For uniqueness, let \(H \le G\) with \(|H| = d\). By (1), \(H = \langle g^m\rangle\) for the least positive \(m\) with \(g^m \in H\); and by Lagrange's theorem (proved below, but a direct argument suffices here) \(|H| = \ord(g^m) = n/\gcd(n,m) = d\), so \(\gcd(n,m) = n/d\). Writing \(m = (n/d) m'\), every element of \(H\) is a power of \(g^{n/d}\) raised to a power, hence \(H \subseteq \langle g^{n/d}\rangle\); comparing orders gives \(H = \langle g^{n/d}\rangle\). Thus there is exactly one subgroup of each order \(d \mid n\), and by (1) every subgroup arises this way.
\end{proof}

\begin{remark}\label{groups:rem-Zn-subgroups}
Specialised to \(G = \Z/n\Z\), part (3) states: for each \(d \mid n\) there is a unique subgroup of order \(d\), namely \(\langle n/d \rangle = \{0, n/d, 2n/d, \dots\}\), which is itself isomorphic to \(\Z/d\Z\). For \(n = 12\), the subgroup lattice mirrors the divisor lattice \(1 \mid 2 \mid 4,\ 1\mid 2 \mid 6,\ 3 \mid 6,\ \dots\) of \(12\).
\end{remark}

\begin{proposition}\label{groups:prop-subgroups-of-Z}
Every subgroup of \((\Z,+)\) has the form \(n\Z\) for a unique integer \(n \ge 0\).
\end{proposition}

\begin{proof}
This is Theorem~\ref{groups:thm-cyclic-subgroups}(2) read additively; we give the direct argument. Let \(H \le \Z\). If \(H = \{0\}\) then \(H = 0\Z\). Otherwise \(H\) contains a nonzero element and hence (being closed under negation) a positive one; let \(n\) be the least positive element of \(H\). Then \(n\Z \subseteq H\). For any \(h \in H\), division gives \(h = nq + r\) with \(0 \le r < n\), and \(r = h - nq \in H\); minimality forces \(r = 0\), so \(h \in n\Z\). Thus \(H = n\Z\). Uniqueness is clear since one recovers \(n\) as the least positive element of \(H\).
\end{proof}

\subsection{Cosets and Lagrange's theorem}

\begin{definition}\label{groups:def-coset}
Let \(H \le G\) and \(a \in G\). The \emph{left coset} of \(H\) by \(a\) is
\[
aH := \{ ah : h \in H \},
\]
and the \emph{right coset} is \(Ha := \{ha : h \in H\}\). We denote the set of left cosets \(G/H\), and the number of left cosets, \([G : H] := |G/H|\), is the \emph{index} of \(H\) in \(G\).
\end{definition}

Cosets are the blocks of a partition, and they all have the same size. This is the mechanism underlying Lagrange's theorem.

\begin{lemma}\label{groups:lem-cosets-partition}
Let \(H \le G\). Then:
\begin{enumerate}
\item \(a \in aH\) for every \(a\); the cosets cover \(G\).
\item Two left cosets are either equal or disjoint: \(aH = bH\) if and only if \(a^{-1} b \in H\), and otherwise \(aH \cap bH = \varnothing\).
\item Every left coset has the same cardinality as \(H\); the map \(h \mapsto ah\) is a bijection \(H \to aH\).
\end{enumerate}
The same statements hold for right cosets, with \(b a^{-1} \in H\) in (2).
\end{lemma}

\begin{proof}
(1) \(a = ae \in aH\) since \(e \in H\). Hence the union of all cosets is \(G\).

(2) Suppose \(aH \cap bH \neq \varnothing\), say \(ah_1 = bh_2\) with \(h_1,h_2 \in H\). Then \(a^{-1}b = h_1 h_2^{-1} \in H\). Conversely if \(a^{-1}b = h \in H\), then \(b = ah\), so for any \(h' \in H\), \(bh' = a(hh') \in aH\), giving \(bH \subseteq aH\); by symmetry (\(b^{-1}a = h^{-1} \in H\)) we obtain \(aH \subseteq bH\), so \(aH = bH\). Thus two cosets intersect only if they coincide; equivalently, distinct cosets are disjoint.

(3) The map \(\lambda_a \colon H \to aH\), \(h \mapsto ah\), is surjective by definition of \(aH\) and injective by left-cancellation (\(ah = ah' \Rightarrow h = h'\)). Hence \(|aH| = |H|\).
\end{proof}

\begin{theorem}[Lagrange]\label{groups:thm-lagrange}
Let \(G\) be a \emph{finite} group and \(H \le G\). Then \(|H|\) divides \(|G|\), and
\[
|G| = [G:H]\,|H|.
\]
\end{theorem}

\begin{proof}
By Lemma~\ref{groups:lem-cosets-partition}, the distinct left cosets of \(H\) partition \(G\) into \([G:H]\) blocks, each of cardinality \(|H|\). Summing the sizes of the blocks,
\[
|G| = \sum_{\text{distinct cosets } aH} |aH| = [G:H] \cdot |H|.
\]
In particular \(|H|\) divides \(|G|\).
\end{proof}

The corollaries are immediate but far-reaching.

\begin{corollary}\label{groups:cor-order-elt-divides}
In a finite group \(G\), the order of every element divides \(|G|\).
\end{corollary}

\begin{proof}
For \(a \in G\), the cyclic subgroup \(\langle a\rangle\) has order \(\ord(a)\) by Proposition~\ref{groups:prop-cyclic-order}. By Lagrange, \(\ord(a) = |\langle a\rangle|\) divides \(|G|\).
\end{proof}

\begin{corollary}\label{groups:cor-aG-e}
If \(G\) is finite of order \(N\), then \(a^N = e\) for every \(a \in G\).
\end{corollary}

\begin{proof}
Let \(d = \ord(a)\). By Corollary~\ref{groups:cor-order-elt-divides}, \(d \mid N\), say \(N = dm\). Then \(a^N = (a^d)^m = e^m = e\).
\end{proof}

\begin{corollary}[Euler, Fermat]\label{groups:cor-euler-fermat}
For \(\gcd(a,n)=1\) one has \(a^{\varphi(n)} \equiv 1 \pmod n\). In particular, for a prime \(p\) and any \(a\) not divisible by \(p\), \(a^{p-1} \equiv 1 \pmod p\), and \(a^p \equiv a \pmod p\) for all \(a\).
\end{corollary}

\begin{proof}
Apply Corollary~\ref{groups:cor-aG-e} to the group \((\Z/n\Z)^\times\) of order \(\varphi(n)\): the class \([a]\) satisfies \([a]^{\varphi(n)} = [1]\). For \(n = p\) prime, \(\varphi(p) = p-1\), yielding Fermat's little theorem; multiplying \(a^{p-1}\equiv 1\) by \(a\) gives \(a^p \equiv a\), which holds even when \(p \mid a\) since both sides are then \(0\).
\end{proof}

\begin{corollary}\label{groups:cor-prime-cyclic}
Every group of prime order \(p\) is cyclic, hence isomorphic to \(\Z/p\Z\), and has no proper nontrivial subgroups.
\end{corollary}

\begin{proof}
Let \(|G| = p\) and pick \(a \neq e\). Then \(\langle a\rangle\) has order \(>1\) dividing \(p\), so \(|\langle a\rangle| = p = |G|\), forcing \(\langle a\rangle = G\). Thus \(G\) is cyclic. Any subgroup has order dividing \(p\), hence order \(1\) or \(p\). The isomorphism with \(\Z/p\Z\) follows from the classification of cyclic groups (Theorem~\ref{groups:thm-cyclic-classification}).
\end{proof}

We close the subsection with the quotient construction in the abelian case, which is all the later development requires: it underlies the quotient rings of \cref{sec:ringsfields} and, through them, the construction of \(\Z/n\Z\) and of polynomial quotients.

\begin{proposition}[Quotient of an abelian group]\label{groups:prop-quotient-abelian}
Let \(A\) be an abelian group, written additively, and let \(H \le A\). On the set \(A/H\) of cosets, the operation
\[
(a + H) + (b + H) := (a + b) + H
\]
is well defined and makes \(A/H\) an abelian group, the \emph{quotient group} of \(A\) by \(H\), with identity \(H = 0 + H\) and inverse \(-(a+H) = (-a) + H\). If \(A\) is finite, \(|A/H| = [A:H] = |A|/|H|\).
\end{proposition}

\begin{proof}
Well-definedness is the only point of substance. Suppose \(a + H = a' + H\) and \(b + H = b' + H\); by Lemma~\ref{groups:lem-cosets-partition}(2) (written additively) this means \(a' - a \in H\) and \(b' - b \in H\). Then \((a' + b') - (a + b) = (a' - a) + (b' - b) \in H\), since \(H\) is closed under addition, so \((a'+b') + H = (a+b) + H\): the sum of cosets does not depend on the chosen representatives. The group axioms now pass from \(A\) to \(A/H\) coset-wise: associativity and commutativity read \(((a+H)+(b+H))+(c+H) = (a+b+c)+H\) in either grouping, the coset \(H\) is the identity, and \((-a)+H\) inverts \(a+H\). The cardinality statement is Lagrange's theorem (Theorem~\ref{groups:thm-lagrange}).
\end{proof}

\begin{remark}\label{groups:rem-quotient-nonabelian}
For a non-abelian group \(G\) the same construction succeeds exactly for the subgroups \(N \le G\) whose left and right cosets coincide (\(gN = Ng\) for all \(g\)), the \emph{normal} subgroups; commutativity makes every subgroup of an abelian group normal, which is why no extra hypothesis appeared above. We need only the abelian case.
\end{remark}

\subsection{Homomorphisms, kernels, and images}

We compare groups by means of maps that respect their operations.

\begin{definition}\label{groups:def-hom}
Let \((G,\ast)\) and \((G',\circ)\) be groups. A \emph{group homomorphism} is a function \(\varphi \colon G \to G'\) such that
\[
\varphi(a \ast b) = \varphi(a) \circ \varphi(b) \qquad \text{for all } a,b \in G.
\]
We denote the set of all homomorphisms \(G \to G'\) by \(\Hom(G,G')\).
\end{definition}

\begin{proposition}\label{groups:prop-hom-basic}
Let \(\varphi \colon G \to G'\) be a homomorphism, with identities \(e \in G\), \(e' \in G'\). Then
\begin{enumerate}
\item \(\varphi(e) = e'\);
\item \(\varphi(a^{-1}) = \varphi(a)^{-1}\) for all \(a\);
\item \(\varphi(a^n) = \varphi(a)^n\) for all \(n \in \Z\);
\item the image \(\varphi(G)\) is a subgroup of \(G'\).
\end{enumerate}
\end{proposition}

\begin{proof}
(1) From \(\varphi(e) = \varphi(e \ast e) = \varphi(e)\circ\varphi(e)\); cancelling \(\varphi(e)\) in \(G'\) gives \(e' = \varphi(e)\).

(2) \(\varphi(a)\circ\varphi(a^{-1}) = \varphi(a \ast a^{-1}) = \varphi(e) = e'\), and likewise on the other side, so \(\varphi(a^{-1})\) is the inverse of \(\varphi(a)\).

(3) For \(n \ge 0\) by induction using the homomorphism property; for \(n < 0\) combine with (2).

(4) Nonempty since \(e' = \varphi(e) \in \varphi(G)\). If \(x = \varphi(a)\), \(y = \varphi(b)\) lie in \(\varphi(G)\), then \(x y^{-1} = \varphi(a)\varphi(b)^{-1} = \varphi(ab^{-1}) \in \varphi(G)\). By the subgroup test, \(\varphi(G) \le G'\).
\end{proof}

\begin{definition}\label{groups:def-kernel}
The \emph{kernel} of a homomorphism \(\varphi \colon G \to G'\) is
\[
\ker\varphi := \{ a \in G : \varphi(a) = e' \} = \varphi^{-1}(\{e'\}),
\]
and its \emph{image} is \(\img\varphi := \varphi(G) = \{\varphi(a) : a \in G\}\).
\end{definition}

\begin{proposition}\label{groups:prop-kernel-injective}
Let \(\varphi \colon G \to G'\) be a homomorphism. Then \(\ker\varphi \le G\), and \(\varphi\) is injective if and only if \(\ker\varphi = \{e\}\).
\end{proposition}

\begin{proof}
\(\ker\varphi\) is nonempty (\(e \in \ker\varphi\) by Proposition~\ref{groups:prop-hom-basic}(1)). If \(a,b \in \ker\varphi\), then \(\varphi(ab^{-1}) = \varphi(a)\varphi(b)^{-1} = e'(e')^{-1} = e'\), so \(ab^{-1} \in \ker\varphi\); thus \(\ker\varphi \le G\).

If \(\varphi\) is injective, then \(\varphi(a)=e'=\varphi(e)\) forces \(a=e\), so \(\ker\varphi = \{e\}\). Conversely, suppose \(\ker\varphi = \{e\}\) and \(\varphi(a) = \varphi(b)\). Then \(\varphi(ab^{-1}) = \varphi(a)\varphi(b)^{-1} = e'\), so \(ab^{-1} \in \ker\varphi = \{e\}\), giving \(ab^{-1} = e\), i.e.\ \(a = b\). Hence \(\varphi\) is injective.
\end{proof}

Kernels are exactly the subgroups by which one may quotient: for an abelian \(G\), Proposition~\ref{groups:prop-quotient-abelian} forms \(G/\ker\varphi\), and in general the kernel satisfies the normality condition of Remark~\ref{groups:rem-quotient-nonabelian}.

\subsection{Isomorphisms}

\begin{definition}\label{groups:def-iso}
A homomorphism \(\varphi \colon G \to G'\) that is bijective is an \emph{isomorphism}; we then write \(G \cong G'\) and call \(G,G'\) \emph{isomorphic}. An isomorphism from \(G\) to itself is an \emph{automorphism}. A homomorphism is a \emph{monomorphism} if injective and an \emph{epimorphism} if surjective.
\end{definition}

\begin{proposition}\label{groups:prop-iso-inverse}
If \(\varphi\colon G \to G'\) is an isomorphism, so is its set-theoretic inverse \(\varphi^{-1}\colon G' \to G\). The relation \(\cong\) is an equivalence relation on any collection of groups.
\end{proposition}

\begin{proof}
Let \(x,y \in G'\), say \(x = \varphi(a)\), \(y = \varphi(b)\). Then \(\varphi^{-1}(xy) = \varphi^{-1}(\varphi(a)\varphi(b)) = \varphi^{-1}(\varphi(ab)) = ab = \varphi^{-1}(x)\varphi^{-1}(y)\), so \(\varphi^{-1}\) is a homomorphism, and it is bijective. Reflexivity of \(\cong\) uses the identity automorphism, symmetry uses \(\varphi^{-1}\), and transitivity uses that a composite of isomorphisms is an isomorphism (the composite of homomorphisms is a homomorphism, and the composite of bijections is a bijection).
\end{proof}

Isomorphic groups are ``the same'' for all group-theoretic purposes: every property expressible in terms of the operation transfers across an isomorphism. For this reason we speak of classifying groups \emph{up to isomorphism}.

\begin{theorem}[Classification of cyclic groups]\label{groups:thm-cyclic-classification}
Every cyclic group is isomorphic to exactly one of the following:
\begin{enumerate}
\item the infinite cyclic group \((\Z, +)\), if it is infinite;
\item the finite cyclic group \((\Z/n\Z, +)\), if it has order \(n\).
\end{enumerate}
In particular two cyclic groups are isomorphic if and only if they have the same order.
\end{theorem}

\begin{proof}
Let \(G = \langle g\rangle\). Define \(\varphi \colon \Z \to G\) by \(\varphi(k) = g^k\). This is a homomorphism from \((\Z,+)\): \(\varphi(k+\ell) = g^{k+\ell} = g^k g^\ell = \varphi(k)\varphi(\ell)\). It is surjective since \(G\) consists exactly of the powers of \(g\).

If \(g\) has infinite order, \(\varphi\) is injective by Proposition~\ref{groups:prop-cyclic-order}(1) (distinct exponents give distinct powers), so \(\varphi\) is an isomorphism \(\Z \cong G\).

If \(g\) has order \(n\), then \(\ker\varphi = \{k : g^k = e\} = n\Z\) by Proposition~\ref{groups:prop-order-divides}. The induced map \(\bar\varphi \colon \Z/n\Z \to G\), \([k] \mapsto g^k\), is well defined (if \([k]=[\ell]\) then \(n \mid k - \ell\), so \(g^k = g^\ell\)), a homomorphism, surjective, and injective (if \(g^k = e\) then \(n \mid k\), so \([k]=[0]\)). Hence \(\Z/n\Z \cong G\). Finally, \(\Z\) and \(\Z/n\Z\) are pairwise non-isomorphic for distinct \(n\) since isomorphic groups have equal order.
\end{proof}


\section{Rings and fields}\label{sec:ringsfields}

The preceding development treated groups: sets equipped with a single associative binary operation possessing an identity and inverses. Most of the structures of genuine interest in cryptography, however, support \emph{two} interacting operations, an addition and a multiplication, that together obey the familiar arithmetic laws of \(\Z\). The integers, the rationals, the residues modulo \(n\), polynomials, and matrices are all of this kind. The abstract structure that captures these examples is the \emph{ring}, and the most important special case --- the one in which division by every nonzero element is possible --- is the \emph{field}. Fields are the natural arena for linear algebra and for the theory of polynomials, and the finite fields \(\Fp\) and \(\Fq\) are the computational substrate of essentially all of the elliptic-curve and proof-system machinery developed later in this monograph. This section builds the theory from the ground up.

Throughout, we assume only the theory of abelian groups developed earlier. We write the additive structure of a ring additively, with identity element \(0\), and the multiplicative structure multiplicatively, with (when it exists) identity element \(1\).

\subsection{Rings and commutative rings}

\begin{definition}\label{ringsfields:def-ring}
A \emph{ring} is a set \(R\) equipped with two binary operations, \emph{addition} \(+\colon R\times R\to R\) and \emph{multiplication} \(\cdot\colon R\times R\to R\), satisfying the following axioms.
\begin{enumerate}
\item \textbf{Additive abelian group.} \((R,+)\) is an abelian group; that is, addition is associative and commutative, there is an element \(0\in R\) with \(a+0=a\) for all \(a\), and every \(a\in R\) has an additive inverse \(-a\) with \(a+(-a)=0\).
\item \textbf{Multiplicative associativity.} \((a\cdot b)\cdot c = a\cdot(b\cdot c)\) for all \(a,b,c\in R\).
\item \textbf{Multiplicative identity.} There is an element \(1\in R\) with \(1\cdot a = a\cdot 1 = a\) for all \(a\in R\).
\item \textbf{Distributivity.} For all \(a,b,c\in R\),
\[
a\cdot(b+c) = a\cdot b + a\cdot c
\qquad\text{and}\qquad
(a+b)\cdot c = a\cdot c + b\cdot c.
\]
\end{enumerate}
A ring \(R\) is \emph{commutative} if in addition \(a\cdot b = b\cdot a\) for all \(a,b\in R\).
\end{definition}

We customarily write \(ab\) for \(a\cdot b\), adopting the convention that multiplication binds more tightly than addition, so that \(ab+ac\) means \((ab)+(ac)\). The element \(0\) is the \emph{additive identity} or \emph{zero}, and \(1\) is the \emph{multiplicative identity} or \emph{unity}.

\begin{remark}\label{ringsfields:rem-conventions}
Conventions in the literature vary. Some authors do not require a multiplicative identity \(1\), reserving the term \emph{ring with unity} for the structure of Definition~\ref{ringsfields:def-ring} and using \emph{rng} (``ring without the \(i\) of identity'') for the weaker notion. In this monograph \emph{every ring has a \(1\)}, every ring homomorphism preserves it, and ``ring'' will frequently mean ``commutative ring'' in the examples that matter to us; we always state commutativity explicitly when a result depends on it.
\end{remark}

Before we produce examples, we record the elementary arithmetic consequences of the axioms. These are the rules we use constantly and without comment; it is worth establishing once that they all follow from distributivity.

\begin{proposition}\label{ringsfields:prop-arith}
Let \(R\) be a ring and let \(a,b\in R\). Then:
\begin{enumerate}
\item \(0\cdot a = a\cdot 0 = 0\);
\item \((-a)b = a(-b) = -(ab)\);
\item \((-a)(-b) = ab\);
\item more generally, for integers \(m,n\) (interpreting \(na\) as the \(n\)-fold sum in the additive group), \((ma)(nb) = (mn)(ab)\).
\end{enumerate}
\end{proposition}

\begin{proof}
(1) Since \(0+0=0\), distributivity gives \(0\cdot a = (0+0)\cdot a = 0\cdot a + 0\cdot a\). Adding \(-(0\cdot a)\) to both sides yields \(0 = 0\cdot a\). The argument for \(a\cdot 0\) is symmetric.

(2) One computes \(ab + (-a)b = (a + (-a))b = 0\cdot b = 0\) by distributivity and part (1). Hence \((-a)b\) is the additive inverse of \(ab\), i.e.\ \((-a)b = -(ab)\). The identity \(a(-b) = -(ab)\) is symmetric.

(3) Applying (2) twice, \((-a)(-b) = -\bigl(a(-b)\bigr) = -\bigl(-(ab)\bigr) = ab\), where the last step uses that negation is an involution in the abelian group \((R,+)\).

(4) This follows from (2) and repeated distributivity (a routine induction on \(|m|\) and \(|n|\)); we omit the bookkeeping.
\end{proof}

\begin{remark}\label{ringsfields:rem-zeroring}
If \(1 = 0\) in a ring \(R\), then for every \(a\in R\) one has \(a = 1\cdot a = 0\cdot a = 0\) by Proposition~\ref{ringsfields:prop-arith}(1), so \(R=\{0\}\). This single-element structure is the \emph{zero ring} (or \emph{trivial ring}); it is a perfectly legitimate ring. Whenever we wish to exclude it we impose the assumption \(1\neq 0\), and indeed the definition of a field below builds in the standing assumption \(1\neq 0\).
\end{remark}

\begin{example}\label{ringsfields:exa-rings}
We list the basic examples; the reader should verify the axioms in each case.
\begin{enumerate}
\item \textbf{The integers \(\Z\)} with ordinary addition and multiplication form a commutative ring. This is the prototype, and much of ring theory is an attempt to isolate which features of \(\Z\) survive in greater generality.
\item \textbf{The rationals \(\Q\), the reals \(\R\), and the complexes \(\C\)} are commutative rings (indeed fields, see \S\ref{ringsfields:subsec-fields}).
\item \textbf{The integers modulo \(n\), \(\Z/n\Z\).} For an integer \(n\geq 1\), the set of residue classes \(\{\,\overline 0,\overline 1,\dots,\overline{n-1}\,\}\) with addition and multiplication induced from \(\Z\) is a commutative ring, the additive group being the cyclic group of order \(n\). One writes \(\Z/n\Z\) (or \(\Z_n\) when no confusion with \(n\)-adic integers can arise). The case \(n=1\) gives the zero ring.
\item \textbf{Polynomial rings.} If \(R\) is a commutative ring, the set \(R[x]\) of polynomials \(\sum_{i=0}^d a_i x^i\) with coefficients \(a_i\in R\), under the usual addition and multiplication, is a commutative ring. The next section treats polynomials over a field systematically; here we note only that \(R[x]\) is again a commutative ring with the same \(0\) and \(1\) as \(R\) (viewed as constant polynomials).
\item \textbf{Matrix rings.} For \(n\geq 1\) the set \(M_n(R)\) of \(n\times n\) matrices over a ring \(R\), with matrix addition and matrix multiplication, is a ring. For \(n\geq 2\) it is \emph{not} commutative even when \(R\) is: for instance \(\left(\begin{smallmatrix}0&1\\0&0\end{smallmatrix}\right)\left(\begin{smallmatrix}0&0\\1&0\end{smallmatrix}\right) = \left(\begin{smallmatrix}1&0\\0&0\end{smallmatrix}\right)\) while \(\left(\begin{smallmatrix}0&0\\1&0\end{smallmatrix}\right)\left(\begin{smallmatrix}0&1\\0&0\end{smallmatrix}\right) = \left(\begin{smallmatrix}0&0\\0&1\end{smallmatrix}\right)\). This is the first important example of a noncommutative ring.
\item \textbf{Product rings.} If \(R\) and \(S\) are rings, the Cartesian product \(R\times S\) with componentwise operations \((r,s)+(r',s')=(r+r',s+s')\) and \((r,s)(r',s')=(rr',ss')\) is a ring, with zero \((0,0)\) and identity \((1,1)\). More generally any finite product \(R_1\times\cdots\times R_k\) is a ring.
\item \textbf{The zero ring} \(\{0\}\) of Remark~\ref{ringsfields:rem-zeroring}.
\end{enumerate}
\end{example}

\subsection{Units, zero divisors, and integral domains}\label{ringsfields:subsec-units}

The integers \(\Z\) enjoy two properties beyond the bare ring axioms that will guide the rest of this section: a product of nonzero integers is nonzero (\emph{cancellation}), and the only integers by which one may divide within \(\Z\) are \(\pm 1\). We now abstract both.

\begin{definition}\label{ringsfields:def-unit}
Let \(R\) be a ring. An element \(u\in R\) is a \emph{unit} (or is \emph{invertible}) if there exists \(v\in R\) with \(uv = vu = 1\). The element \(v\) is then unique (if \(uv=vu=1\) and \(uv'=v'u=1\) then \(v = v(uv') = (vu)v' = v'\)), is called the \emph{inverse} of \(u\), and is written \(u^{-1}\). We denote the set of units of \(R\) by \(R^{\times}\).
\end{definition}

\begin{proposition}\label{ringsfields:prop-units-group}
For any ring \(R\), the set \(R^{\times}\) of units is a group under multiplication, the \emph{group of units} (or \emph{multiplicative group}) of \(R\).
\end{proposition}

\begin{proof}
The identity \(1\) is a unit (it is its own inverse), so \(1\in R^{\times}\). If \(u,w\in R^{\times}\) with inverses \(u^{-1},w^{-1}\), then \((uw)(w^{-1}u^{-1}) = u(ww^{-1})u^{-1} = uu^{-1}=1\) and symmetrically \((w^{-1}u^{-1})(uw)=1\), so \(uw\in R^{\times}\) with \((uw)^{-1}=w^{-1}u^{-1}\); thus \(R^{\times}\) is closed under multiplication. Associativity is inherited from \(R\). Finally every \(u\in R^{\times}\) has the inverse \(u^{-1}\in R^{\times}\) (whose inverse is \(u\)). Hence \(R^{\times}\) is a group.
\end{proof}

\begin{example}\label{ringsfields:exa-units}
\begin{enumerate}
\item \(\Z^{\times} = \{1,-1\}\).
\item \(\Q^{\times}=\Q\setminus\{0\}\), and likewise \(\R^{\times}=\R\setminus\{0\}\), \(\C^{\times}=\C\setminus\{0\}\). In a field, \emph{every} nonzero element is a unit; this is the defining feature, taken up in \S\ref{ringsfields:subsec-fields}.
\item In \(\Z/n\Z\), the class \(\overline a\) is a unit if and only if \(\gcd(a,n)=1\) (Theorem~\ref{intmod:thm:unitiff}); thus \((\Z/n\Z)^{\times}\) consists of the residues coprime to \(n\), and its order is \(\varphi(n)\), Euler's totient.
\item \(M_n(R)^{\times}\) is the group of invertible matrices. When \(R\) is a field this is the \emph{general linear group} \(\mathrm{GL}_n\), consisting of the matrices of nonzero determinant.
\end{enumerate}
\end{example}

We now isolate the failure of cancellation.

\begin{definition}\label{ringsfields:def-zerodiv}
Let \(R\) be a ring. A nonzero element \(a\in R\) is a \emph{left zero divisor} if there exists a nonzero \(b\in R\) with \(ab = 0\); \emph{right zero divisor} is defined symmetrically with \(ba=0\). In a commutative ring the two notions coincide and one simply says \emph{zero divisor}. An element that is neither a left nor a right zero divisor (and is nonzero) is called \emph{regular} or \emph{cancellable}.
\end{definition}

The following observation justifies the terminology, making precise the sense in which zero divisors obstruct ``dividing out''.

\begin{proposition}\label{ringsfields:prop-cancel}
Let \(R\) be a ring and \(a\in R\) a nonzero element. Then \(a\) is not a left zero divisor if and only if it is \emph{left-cancellable}: \(ab = ac\) implies \(b=c\) for all \(b,c\in R\).
\end{proposition}

\begin{proof}
Suppose \(a\) is not a left zero divisor and \(ab=ac\). Then \(a(b-c)=ab-ac=0\) by distributivity and Proposition~\ref{ringsfields:prop-arith}(2). Since \(a\) is not a left zero divisor and \(a\neq 0\), the factor \(b-c\) must be \(0\), i.e.\ \(b=c\). Conversely, if \(a\) is a left zero divisor, say \(ab=0\) with \(b\neq 0\), then \(ab = a\cdot 0\) but \(b\neq 0\), so cancellation fails.
\end{proof}

\begin{remark}\label{ringsfields:rem-unit-not-zd}
A unit is never a zero divisor: if \(u\) is a unit and \(ub=0\), then \(b = u^{-1}(ub) = u^{-1}\cdot 0 = 0\). The converse fails in general --- in \(\Z\) every nonzero element is a non-zero-divisor, but only \(\pm1\) are units --- though it does hold in finite rings, as Theorem~\ref{ringsfields:thm-finite-domain} establishes.
\end{remark}

\begin{example}\label{ringsfields:exa-zerodiv}
\begin{enumerate}
\item In \(\Z/6\Z\), \(\overline 2\cdot\overline 3 = \overline 6 = \overline 0\), so \(\overline 2\) and \(\overline 3\) are zero divisors. More generally \(\overline a\neq\overline 0\) is a zero divisor in \(\Z/n\Z\) precisely when \(\gcd(a,n)>1\); thus in \(\Z/n\Z\) every nonzero element is either a unit or a zero divisor (Theorem~\ref{ringsfields:thm-finite-domain} revisits this dichotomy).
\item In the product ring \(R\times S\) with both factors nonzero, \((1,0)(0,1) = (0,0)\), so product rings always have zero divisors.
\item In \(M_2(\R)\), \(\left(\begin{smallmatrix}1&0\\0&0\end{smallmatrix}\right)\left(\begin{smallmatrix}0&0\\0&1\end{smallmatrix}\right) = \left(\begin{smallmatrix}0&0\\0&0\end{smallmatrix}\right)\), so singular nonzero matrices are zero divisors.
\end{enumerate}
\end{example}

The rings with no zero divisors are precisely those in which the cancellation law of \(\Z\) holds; in the commutative case they receive a name.

\begin{definition}\label{ringsfields:def-domain}
An \emph{integral domain} (or simply a \emph{domain}) is a commutative ring \(R\) with \(1\neq 0\) and no zero divisors: for all \(a,b\in R\), \(ab=0\) implies \(a=0\) or \(b=0\).
\end{definition}

Equivalently, by Proposition~\ref{ringsfields:prop-cancel}, an integral domain is a nonzero commutative ring in which every nonzero element is cancellable.

\begin{example}\label{ringsfields:exa-domain}
\(\Z\) is an integral domain --- this is the example that names the concept. Every field is an integral domain (Proposition~\ref{ringsfields:prop-field-domain}). The ring \(\Z/n\Z\) is an integral domain if and only if \(n\) is prime, by Example~\ref{ringsfields:exa-zerodiv}(1) together with the fact that for \(n=p\) prime, \(ab\equiv 0\pmod p\) forces \(p\mid a\) or \(p\mid b\) (Euclid's lemma). If \(R\) is an integral domain then so is the polynomial ring \(R[x]\): the leading coefficient of a product of nonzero polynomials is the product of the leading coefficients, hence nonzero. By contrast \(\Z/6\Z\), product rings, and \(M_n(R)\) for \(n\geq 2\) are not domains.
\end{example}

\subsection{Ideals and quotient rings}\label{ringsfields:subsec-ideals}

The formation of quotients of rings --- the ring-theoretic analogue of the quotient of an abelian group (Proposition~\ref{groups:prop-quotient-abelian}) --- requires the appropriate notion of substructure by which to quotient. For rings this is the \emph{ideal}. We treat this material briefly, as we need it chiefly to construct \(\Z/n\Z\) conceptually and to indicate how the finite fields \(\Fq\) would arise; a fuller treatment belongs to a course in commutative algebra.

\begin{definition}\label{ringsfields:def-ideal}
Let \(R\) be a commutative ring. A subset \(I\subseteq R\) is an \emph{ideal} if
\begin{enumerate}
\item \(I\) is an additive subgroup of \((R,+)\): \(0\in I\), and \(a-b\in I\) whenever \(a,b\in I\); and
\item \(I\) is closed under multiplication by arbitrary ring elements: \(ra\in I\) whenever \(r\in R\) and \(a\in I\) (the \emph{absorption} property).
\end{enumerate}
\end{definition}

In a noncommutative ring one distinguishes left, right, and two-sided ideals; for the present purposes commutative rings suffice, where all three coincide. Note that an ideal is almost never a subring: if an ideal \(I\) contains \(1\), then by absorption \(r = r\cdot 1\in I\) for every \(r\in R\), so \(I=R\). Thus the only ideal containing a unit is \(R\) itself.

\begin{example}\label{ringsfields:exa-ideal}
\begin{enumerate}
\item In any commutative ring \(R\), the subsets \(\{0\}\) and \(R\) are ideals, the \emph{trivial} and \emph{improper} ideals respectively.
\item For a fixed \(a\in R\), the set \((a) := aR = \{\,ar : r\in R\,\}\) is an ideal, the \emph{principal ideal generated by \(a\)}. In \(\Z\) the principal ideal \((n) = n\Z\) consists of all multiples of \(n\). In fact \emph{every} ideal of \(\Z\) is principal: an ideal \(I\neq\{0\}\) contains a least positive element \(n\), and the division algorithm shows \(I = n\Z\) (any \(a\in I\) gives \(a = qn + r\) with \(0\le r<n\), whence \(r = a-qn\in I\) forces \(r=0\)). A domain in which every ideal is principal is a \emph{principal ideal domain}.
\item More generally, for \(a_1,\dots,a_k\in R\), the set \((a_1,\dots,a_k) = \{\sum_i r_i a_i : r_i\in R\}\) is the ideal \emph{generated by} the \(a_i\).
\end{enumerate}
\end{example}

The significance of ideals is that they are exactly the subsets by which one may quotient.

\begin{definition}\label{ringsfields:def-quotient}
Let \(I\) be an ideal of a commutative ring \(R\). Since \((I,+)\) is a subgroup of the abelian group \((R,+)\), we may form the quotient group \(R/I\) of additive cosets \(a + I\) (Proposition~\ref{groups:prop-quotient-abelian}). We define multiplication of cosets by
\[
(a+I)(b+I) := ab + I .
\]
\end{definition}

\begin{proposition}\label{ringsfields:prop-quotient-well-defined}
With the operations of Definition~\ref{ringsfields:def-quotient}, \(R/I\) is a commutative ring, the \emph{quotient ring} of \(R\) by \(I\), with zero \(0+I = I\) and identity \(1+I\). In particular the multiplication is well defined.
\end{proposition}

\begin{proof}
We must check that the product of cosets does not depend on the chosen representatives. Suppose \(a+I = a'+I\) and \(b+I = b'+I\), so that \(a' = a + s\) and \(b' = b + t\) with \(s,t\in I\). Then
\[
a'b' = (a+s)(b+t) = ab + at + sb + st .
\]
The three terms \(at\), \(sb\), and \(st\) all lie in \(I\): \(at\in I\) and \(st\in I\) by absorption (\(t\in I\)), and \(sb\in I\) by absorption (\(s\in I\)). Hence \(a'b' - ab \in I\), i.e.\ \(a'b' + I = ab + I\), and multiplication is well defined. The ring axioms for \(R/I\) are then inherited coset-wise from those of \(R\): for instance distributivity reads \((a+I)\bigl((b+I)+(c+I)\bigr) = (a(b+c))+I = (ab+ac)+I = (a+I)(b+I)+(a+I)(c+I)\). Commutativity, associativity, and the identity laws follow identically.
\end{proof}

\begin{example}\label{ringsfields:exa-quotient-Zn}
Taking \(R=\Z\) and \(I = n\Z\), we recover the quotient ring \(\Z/n\Z\) of Example~\ref{ringsfields:exa-rings}(3) \emph{as a quotient ring}: its elements are the cosets \(a + n\Z\), i.e.\ the residue classes modulo \(n\), and the coset operations are exactly residue arithmetic. This is the conceptual origin of modular arithmetic, and it generalises: quotients of polynomial rings \(\Fp[x]/(f)\) by an irreducible polynomial \(f\) furnish the finite fields \(\Fq\) with \(q = p^k\). That construction lies off our path --- the curves of Halo~2 and Orchard work over prime fields --- so we record its existence and do not develop it.
\end{example}
\subsection{Fields}\label{ringsfields:subsec-fields}

We now introduce the central object of this section.

\begin{definition}\label{ringsfields:def-field}
A \emph{field} is a commutative ring \(F\) in which \(1\neq 0\) and every nonzero element is a unit. Equivalently, \(F\) is a set with two operations \(+\) and \(\cdot\) such that \((F,+)\) is an abelian group with identity \(0\), \((F\setminus\{0\},\cdot)\) is an abelian group with identity \(1\), and multiplication distributes over addition.
\end{definition}

Unwinding the definition: in a field one may add, subtract, and multiply as in any commutative ring, and additionally \emph{divide by any nonzero element}, with \(a/b := ab^{-1}\). The group \((F\setminus\{0\},\cdot) = F^{\times}\) is the \emph{multiplicative group} of the field.

\begin{example}\label{ringsfields:exa-fields}
\begin{enumerate}
\item \(\Q\), \(\R\), and \(\C\) are fields; \(\Z\) is not, since \(2\) has no inverse in \(\Z\).
\item For \(p\) prime, \(\Fp := \Z/p\Z\) is a field, the \emph{prime field} of characteristic \(p\) (Corollary~\ref{intmod:cor:fpfield}; Theorem~\ref{ringsfields:thm-finite-domain} below recovers this abstractly). These fields, and their finite extensions \(\Fq\) with \(q=p^k\), are fundamental to the cryptography developed later.
\item The field \(\Q(i) = \{a+bi : a,b\in\Q\}\subseteq\C\) and, more generally, number fields, are fields obtained by adjoining algebraic elements to \(\Q\).
\item The field of rational functions \(F(x) = \{\,f/g : f,g\in F[x],\ g\neq 0\,\}\) over a field \(F\) is a field; it is the field of fractions of \(F[x]\).
\end{enumerate}
\end{example}

\begin{proposition}\label{ringsfields:prop-field-domain}
Every field is an integral domain.
\end{proposition}

\begin{proof}
Let \(F\) be a field; by definition \(1\neq 0\) and \(F\) is commutative. Suppose \(ab=0\) with \(a\neq 0\). Since \(a\) is a unit, multiply by \(a^{-1}\) on the left: \(b = 1\cdot b = (a^{-1}a)b = a^{-1}(ab) = a^{-1}\cdot 0 = 0\). Hence \(a=0\) or \(b=0\), so \(F\) has no zero divisors.
\end{proof}

The converse is false (\(\Z\) is a domain but not a field), but it becomes true under finiteness --- a fact of central importance for finite fields.

\begin{theorem}\label{ringsfields:thm-finite-domain}
Every finite integral domain is a field.
\end{theorem}

\begin{proof}
Let \(R\) be a finite integral domain and let \(a\in R\) with \(a\neq 0\); we must produce a multiplicative inverse for \(a\). Consider the ``multiplication by \(a\)'' map
\[
\mu_a\colon R\to R,\qquad \mu_a(x) = ax .
\]
This map is injective: if \(ax = ay\) then \(a(x-y)=0\), and since \(a\neq 0\) and \(R\) is a domain, \(x-y=0\), i.e.\ \(x=y\) (this is cancellation, Proposition~\ref{ringsfields:prop-cancel}). An injective map from a finite set to itself is necessarily surjective. Hence \(\mu_a\) is onto, so there exists \(b\in R\) with \(ab = \mu_a(b) = 1\). Because \(R\) is commutative, \(ba=1\) as well, so \(b = a^{-1}\). As \(a\) was an arbitrary nonzero element, every nonzero element is a unit, and \(R\) is a field.
\end{proof}

Applied to \(\Z/n\Z\), the theorem recovers Corollary~\ref{intmod:cor:fpfield}: for prime \(p\) the finite domain \(\Z/p\Z\) (Example~\ref{ringsfields:exa-domain}) is a field, and there the extended Euclidean algorithm computed each inverse explicitly.

\subsection{The characteristic of a ring}\label{ringsfields:subsec-char}

There is a canonical way in which a copy of \(\Z\) maps into any ring, and the size of its image is a fundamental invariant.

\begin{definition}\label{ringsfields:def-char}
Let \(R\) be a ring. For an integer \(n\) and the identity \(1\in R\), write \(n\cdot 1\) for the \(n\)-fold sum \(1+1+\cdots+1\) (\(n\) terms) if \(n>0\), the element \(0\) if \(n=0\), and \(-\bigl((-n)\cdot 1\bigr)\) if \(n<0\). The \emph{characteristic} of \(R\), written \(\operatorname{char}(R)\), is the least positive integer \(n\) such that \(n\cdot 1 = 0\), if such an \(n\) exists; otherwise \(\operatorname{char}(R) = 0\).
\end{definition}

Equivalently, the map \(\iota\colon\Z\to R\) given by \(\iota(n) = n\cdot 1\) is a ring homomorphism, and \(\operatorname{char}(R)\) is the nonnegative generator of its kernel, an ideal of \(\Z\). Thus \(\operatorname{char}(R) = 0\) when \(\iota\) is injective (a copy of \(\Z\) sits inside \(R\)), and \(\operatorname{char}(R) = n\) when the kernel is \(n\Z\).

\begin{example}\label{ringsfields:exa-char}
\(\operatorname{char}(\Z) = \operatorname{char}(\Q) = \operatorname{char}(\R) = \operatorname{char}(\C) = 0\). For the residue rings, \(\operatorname{char}(\Z/n\Z) = n\), since \(n\cdot\overline 1 = \overline n = \overline 0\) and no smaller positive multiple of \(\overline 1\) vanishes. In particular \(\operatorname{char}(\Fp) = p\).
\end{example}

\begin{theorem}\label{ringsfields:thm-char-prime}
The characteristic of an integral domain (in particular of a field) is either \(0\) or a prime number.
\end{theorem}

\begin{proof}
Let \(R\) be an integral domain and suppose \(\operatorname{char}(R) = n > 0\). First, \(n\neq 1\): if \(n=1\) then \(1 = 1\cdot 1 = 0\), forcing \(R\) to be the zero ring, contrary to \(1\neq 0\) in a domain. Suppose, for contradiction, that \(n\) is composite, \(n = ab\) with \(1 < a,b < n\). Then in \(R\),
\[
0 = n\cdot 1 = (ab)\cdot 1 = (a\cdot 1)(b\cdot 1),
\]
using the distributive bookkeeping of Proposition~\ref{ringsfields:prop-arith}(4). Since \(R\) is a domain, either \(a\cdot 1 = 0\) or \(b\cdot 1 = 0\). But \(0<a<n\) and \(0<b<n\), contradicting the minimality of \(n=\operatorname{char}(R)\). Hence \(n\) has no nontrivial factorisation, i.e.\ \(n\) is prime.
\end{proof}

\begin{remark}\label{ringsfields:rem-prime-subfield}
A field \(F\) of characteristic \(0\) contains a copy of \(\Z\) (via \(\iota\)) and hence, dividing, a copy of \(\Q\); a field of characteristic \(p\) contains a copy of \(\Fp = \Z/p\Z\). In either case this smallest subfield is the \emph{prime subfield} of \(F\). Every field is therefore an extension either of \(\Q\) or of some \(\Fp\). This dichotomy --- characteristic zero versus characteristic \(p\) --- pervades the subject; the fields \(\Fq\) underlying the cryptography treated here all have prime characteristic \(p\), where the ``freshman's dream'' identity \((a+b)^p = a^p + b^p\) holds and gives rise to the Frobenius endomorphism, developed in the finite-fields section (\cref{finitefields:lem:freshman}).
\end{remark}

\section{Polynomials over a field}\label{sec:polynomials}

Throughout this section, \(\F\) denotes a field. Polynomials over a field are
among the central objects of the entire development: they are the engine behind
finite fields, error-correcting codes, and — most importantly for the present
purposes —
the polynomial commitment schemes and arithmetization that underlie Halo~2 and
Zcash Orchard. The single most exploited fact in that entire edifice is
elementary: \emph{a nonzero polynomial of degree \(d\) over a field has at most \(d\)
roots}. From it we extract Lagrange interpolation and the Schwartz--Zippel
principle, the probabilistic heart of every interactive proof we eventually
construct. This section develops all of this from the definition of a polynomial,
assuming only the theory of fields of the preceding section.

\subsection{The polynomial ring \texorpdfstring{\(\F[X]\)}{F[X]}}
\label{polynomials:subsec:ring}

We begin with the formal definition. The crucial pedagogical point — to which the
development returns at length in Subsection~\ref{polynomials:subsec:formal-vs-function} — is
that a polynomial is \emph{not} a function. It is a formal expression, a finite
list of coefficients, that one may \emph{evaluate} to obtain a function, but which
carries more information than that function in general.

\begin{definition}[Polynomial]\label{polynomials:def:polynomial}
A \emph{polynomial} over \(\F\) in the indeterminate \(X\) is a formal expression
\[
  f = a_0 + a_1 X + a_2 X^2 + \cdots + a_n X^n = \sum_{i=0}^{n} a_i X^i,
\]
where \(n \in \N\) (taking \(\N = \{0, 1, 2, \dots\}\)) and the
\emph{coefficients} \(a_0, a_1, \dots, a_n\) lie in \(\F\). Formally, we may
identify \(f\) with the sequence \((a_0, a_1, a_2, \dots)\) of its coefficients,
in which only finitely many entries are nonzero. Two polynomials are
\emph{equal} precisely when their coefficient sequences agree term by term,
i.e.\ \(\sum_i a_i X^i = \sum_i b_i X^i\) if and only if \(a_i = b_i\) for every
\(i \in \N\). We call the symbol \(X\) an \emph{indeterminate} or
\emph{formal variable}; it is not an element of \(\F\) and stands for nothing in
particular — it is merely a placeholder that organizes the coefficients. We denote the set
of all polynomials over \(\F\) in \(X\) by \(\F[X]\).
\end{definition}

\begin{remark}\label{polynomials:rem:sequence-model}
The sequence model is the cleanest way to render Definition~\ref{polynomials:def:polynomial}
fully rigorous: \(\F[X]\) is the set of all functions \(a \colon \N \to \F\) that
are \emph{eventually zero} (there is some \(N\) with \(a(i) = 0\) for all
\(i > N\)). The expression \(\sum_i a_i X^i\) is then merely suggestive notation
for the sequence \((a_0, a_1, \dots)\). We use the two
points of view interchangeably. The notation is convenient because, as we show below, the
arithmetic of these sequences is exactly the arithmetic the notation suggests.
\end{remark}

We now equip the ring \(\F[X]\) with addition and multiplication. Let
\(f = \sum_i a_i X^i\) and \(g = \sum_i b_i X^i\), where we pad the shorter
coefficient sequence with zeros so that both run over the same index set.

\begin{definition}[{Ring operations on \(\F[X]\)}]\label{polynomials:def:operations}
The \emph{sum} and \emph{product} of \(f = \sum_i a_i X^i\) and
\(g = \sum_i b_i X^i\) are
\[
  f + g = \sum_{i=0}^{\infty} (a_i + b_i)\, X^i,
  \qquad
  f \cdot g = \sum_{k=0}^{\infty} c_k X^k,
  \quad\text{where}\quad
  c_k = \sum_{i+j = k} a_i b_j = \sum_{i=0}^{k} a_i\, b_{k-i}.
\]
The coefficient \(c_k\) is the \emph{convolution} of the two coefficient
sequences. Both sums are finite: only finitely many \(a_i\) and \(b_j\) are
nonzero, so only finitely many \(c_k\) are nonzero, and the result is again a
polynomial.
\end{definition}

The product formula is exactly what we obtain by multiplying out
\(\bigl(\sum_i a_i X^i\bigr)\bigl(\sum_j b_j X^j\bigr)\) and collecting like
powers of \(X\), using the rule \(X^i \cdot X^j = X^{i+j}\). Stating
it as a convolution makes no appeal to "substituting a value for
\(X\)"; it is a purely formal manipulation of coefficient sequences.

\begin{theorem}[{\(\F[X]\) is a commutative ring}]\label{polynomials:thm:ring}
With the operations of Definition~\ref{polynomials:def:operations}, the set
\(\F[X]\) is a commutative ring with the following distinguished elements:
\begin{itemize}
  \item the \emph{zero polynomial} \(0\), the sequence \((0,0,0,\dots)\), which
        is the additive identity;
  \item the \emph{constant polynomial} \(1\), the sequence \((1,0,0,\dots)\),
        which is the multiplicative identity.
\end{itemize}
Moreover, the map \(c \mapsto (c,0,0,\dots)\) embeds \(\F\) into \(\F[X]\) as a
subring (the \emph{constant polynomials}), so that \(\F \subseteq \F[X]\).
\end{theorem}

\begin{proof}
We verify the ring axioms. \emph{Additive structure.} Addition is defined
coordinatewise, so it inherits commutativity, associativity, the identity
\((0,0,\dots)\), and inverses \((-a_0, -a_1, \dots)\) directly from the additive
group of \(\F\). Thus \((\F[X], +)\) is an abelian group.

\emph{Commutativity of multiplication.} With
\(f \cdot g = \sum_k \bigl(\sum_{i+j=k} a_i b_j\bigr) X^k\), swapping the roles
of \(f\) and \(g\) replaces \(a_i b_j\) by \(b_i a_j\); reindexing the inner sum
\((i,j) \mapsto (j,i)\) and using commutativity of \(\F\) shows
\(\sum_{i+j=k} a_i b_j = \sum_{i+j=k} b_i a_j\). Hence \(f \cdot g = g \cdot f\).

\emph{Associativity of multiplication.} For \(f, g, h\) with coefficients
\(a_i, b_j, c_\ell\), the coefficient of \(X^m\) in \((fg)h\) is
\[
  \sum_{k + \ell = m}\Bigl(\sum_{i+j=k} a_i b_j\Bigr) c_\ell
  = \sum_{i + j + \ell = m} a_i b_j c_\ell,
\]
and the same triple sum arises symmetrically as the coefficient of \(X^m\) in
\(f(gh)\). The two agree, so \((fg)h = f(gh)\).

\emph{Distributivity.} The coefficient of \(X^k\) in \(f(g+h)\) is
\(\sum_{i+j=k} a_i (b_j + c_j) = \sum_{i+j=k} a_i b_j + \sum_{i+j=k} a_i c_j\),
which is the coefficient of \(X^k\) in \(fg + fh\); this uses distributivity in
\(\F\). Hence \(f(g+h) = fg + fh\), and the other distributive law follows by
commutativity.

\emph{Identities.} The sequence \((0,0,\dots)\) is plainly additively neutral.
For \(1 = (1,0,0,\dots)\), the convolution gives \((1 \cdot f)_k =
\sum_{i+j=k} \mathbf{1}_{[i=0]} a_j = a_k\), so \(1 \cdot f = f\).

\emph{The embedding.} The map \(\iota \colon c \mapsto (c,0,0,\dots)\) sends sums
to sums and, since \((c,0,\dots)\cdot(d,0,\dots)\) has \(k\)-th coefficient
\(\sum_{i+j=k}\iota(c)_i \iota(d)_j = cd\,\mathbf{1}_{[k=0]}\), products to
products; it is injective because distinct constants give distinct sequences.
Thus \(\iota\) is an injective ring homomorphism, and \(\F[X]\) is a commutative
ring containing a copy of \(\F\).
\end{proof}

\begin{remark}\label{polynomials:rem:scalar}
Multiplication of a polynomial \(f = \sum_i a_i X^i\) by a constant
\(c \in \F\) (viewed as the constant polynomial) is the
\emph{scalar multiple} \(cf = \sum_i (c\,a_i) X^i\). Together with addition this
makes \(\F[X]\) a vector space over \(\F\) — indeed an infinite-dimensional one,
with basis \(\{1, X, X^2, X^3, \dots\}\). We rely on this linear structure
constantly, for example when counting polynomials of bounded degree.
\end{remark}

\subsection{Degree and leading coefficients}
\label{polynomials:subsec:degree}

\begin{definition}[Degree, leading coefficient, monic]\label{polynomials:def:degree}
Let \(f = \sum_{i} a_i X^i \in \F[X]\) be nonzero. Its \emph{degree}
\(\deg f\) is the largest index \(n\) for which \(a_n \neq 0\). The corresponding
coefficient \(a_n\) is the \emph{leading coefficient} of \(f\), written
\(\operatorname{lc}(f)\), and \(a_n X^n\) is the \emph{leading term}. The
coefficient \(a_0\) is the \emph{constant term}. A nonzero polynomial is
\emph{monic} if its leading coefficient is \(1\). The \emph{constant polynomials}
are the elements of \(\F\) (embedded as in
Theorem~\ref{polynomials:thm:ring}); the nonzero ones have degree \(0\), while
the zero polynomial receives degree \(-\infty\) by the following convention. We set \(\deg 0 = -\infty\), with the arithmetic conventions \(-\infty < n\),
\(-\infty + n = -\infty\), and \(-\infty + (-\infty) = -\infty\) for every
\(n \in \Z\). We call polynomials of degree \(1, 2, 3\) \emph{linear},
\emph{quadratic}, \emph{cubic} respectively.
\end{definition}

The \(-\infty\) convention for the zero polynomial is not pedantry: it makes the
degree formulas below hold without exceptions, which streamlines every later
induction on degree.

\begin{proposition}[Degree of sums and products]\label{polynomials:prop:degree}
For all \(f, g \in \F[X]\),
\begin{align}
  \deg(f + g) &\le \max\{\deg f, \deg g\},
    \label{polynomials:eq:deg-sum}\\
  \deg(f \cdot g) &= \deg f + \deg g.
    \label{polynomials:eq:deg-prod}
\end{align}
Equality holds in \eqref{polynomials:eq:deg-sum} whenever
\(\deg f \neq \deg g\).
\end{proposition}

\begin{proof}
If either \(f\) or \(g\) is zero, both claims reduce to the \(-\infty\)
conventions, so assume \(f, g \neq 0\) with \(\deg f = m\), \(\deg g = n\), and
leading coefficients \(a_m, b_n \neq 0\).

For the sum, the coefficient of \(X^k\) in \(f + g\) is \(a_k + b_k\), which is
\(0\) for every \(k > \max\{m,n\}\); hence \(\deg(f+g) \le \max\{m,n\}\). If
\(m \neq n\), say \(m > n\), then the coefficient of \(X^m\) in \(f+g\) equals
\(a_m + 0 = a_m \neq 0\), giving \(\deg(f+g) = m = \max\{m,n\}\). (When
\(m = n\), cancellation \(a_m + b_m = 0\) can occur, lowering the degree; this is
why \eqref{polynomials:eq:deg-sum} is only an inequality in general.)

For the product, the coefficient of \(X^{m+n}\) is
\(c_{m+n} = \sum_{i+j = m+n} a_i b_j\). Every term with \(i > m\) has
\(a_i = 0\), and every term with \(i < m\) forces \(j = m+n-i > n\), so
\(b_j = 0\); thus the only surviving term is \(i = m, j = n\), giving
\(c_{m+n} = a_m b_n\). For \(k > m+n\) the same argument shows every term
vanishes, so \(c_k = 0\). Since \(\F\) is a field it has no zero divisors, so
\(a_m b_n \neq 0\); hence \(\deg(fg) = m + n\) and the leading coefficient of
\(fg\) is \(a_m b_n\).
\end{proof}

The product formula~\eqref{polynomials:eq:deg-prod} is the first place the field
hypothesis does real work: it required \(a_m b_n \neq 0\), i.e.\ the
absence of zero divisors. This single fact has sweeping consequences.

\begin{corollary}[{\(\F[X]\) is an integral domain}]\label{polynomials:cor:domain}
The ring \(\F[X]\) has no zero divisors: if \(f g = 0\) then \(f = 0\) or
\(g = 0\). Equivalently, \(\F[X]\) is an integral domain.
\end{corollary}

\begin{proof}
If both \(f\) and \(g\) were nonzero, then \(\deg(fg) = \deg f + \deg g \ge 0\)
by \eqref{polynomials:eq:deg-prod}, so \(fg \neq 0\). Contrapositively,
\(fg = 0\) forces \(f = 0\) or \(g = 0\).
\end{proof}

\begin{corollary}[Cancellation]\label{polynomials:cor:cancellation}
If \(f, g, h \in \F[X]\) with \(h \neq 0\) and \(f h = g h\), then \(f = g\).
\end{corollary}

\begin{proof}
From \(fh = gh\) it follows that \((f - g) h = 0\); since \(h \neq 0\),
Corollary~\ref{polynomials:cor:domain} forces \(f - g = 0\).
\end{proof}

\begin{proposition}[{Units of \(\F[X]\)}]\label{polynomials:prop:units}
A polynomial \(f \in \F[X]\) is a \emph{unit} (has a multiplicative inverse in
\(\F[X]\)) if and only if it is a nonzero constant. Thus the units of \(\F[X]\)
are exactly the nonzero elements of \(\F\).
\end{proposition}

\begin{proof}
A nonzero constant \(c \in \F^\times\) has inverse \(c^{-1} \in \F \subseteq
\F[X]\). Conversely, if \(fg = 1\) then by \eqref{polynomials:eq:deg-prod},
\(\deg f + \deg g = \deg 1 = 0\); since degrees of nonzero polynomials are
nonnegative, \(\deg f = \deg g = 0\), so \(f\) is a nonzero constant.
\end{proof}

\begin{definition}[Associates]\label{polynomials:def:associates}
Two polynomials \(f, g \in \F[X]\) are \emph{associates} if \(f = u g\) for some
unit \(u \in \F^\times\), i.e.\ if they differ only by a nonzero scalar factor.
\end{definition}

Every nonzero polynomial \(f\) has exactly one monic associate, obtained by
dividing through by its leading coefficient: \(\operatorname{lc}(f)^{-1} f\).
Choosing the monic representative is the standard way to pin down "the" gcd or
"the" irreducible factor uniquely, and we adopt it below.
\subsection{The division algorithm}
\label{polynomials:subsec:division}

The arithmetic of \(\F[X]\) parallels that of the integers \(\Z\) with remarkable
fidelity, and the source of the parallel is a division-with-remainder statement.
Here the degree plays the role that absolute value plays in \(\Z\).

\begin{theorem}[Division algorithm]\label{polynomials:thm:division}
Let \(f, g \in \F[X]\) with \(g \neq 0\). Then there exist \emph{unique}
polynomials \(q, r \in \F[X]\) such that
\[
  f = q\, g + r,
  \qquad\text{with}\qquad
  \deg r < \deg g.
\]
We call \(q\) the \emph{quotient} and \(r\) the \emph{remainder}; we write
\(r = f \bmod g\) and \(q = f \operatorname{div} g\).
\end{theorem}

\begin{proof}
\emph{Existence.} Fix \(g \neq 0\) and induct on \(\deg f\). If
\(\deg f < \deg g\) (which includes the case \(f = 0\), where
\(\deg f = -\infty\)), take \(q = 0\) and \(r = f\); then \(f = 0 \cdot g + f\)
with \(\deg r = \deg f < \deg g\).

Otherwise \(m := \deg f \ge \deg g =: n\). Write the leading terms as
\(a_m X^m\) and \(b_n X^n\), with \(b_n \neq 0\). Consider
\[
  f_1 = f - \frac{a_m}{b_n}\, X^{m-n}\, g.
\]
The subtracted term has degree \(m\) and leading coefficient
\((a_m / b_n)\, b_n = a_m\), so it cancels the leading term of \(f\); hence
\(\deg f_1 < m\) (or \(f_1 = 0\)). By the induction hypothesis applied to
\(f_1\), there are \(q_1, r\) with \(f_1 = q_1 g + r\) and \(\deg r < \deg g\).
Then
\[
  f = f_1 + \frac{a_m}{b_n} X^{m-n} g
    = \Bigl(q_1 + \frac{a_m}{b_n} X^{m-n}\Bigr) g + r,
\]
so \(q = q_1 + (a_m/b_n) X^{m-n}\) and \(r\) work. The field hypothesis enters
visibly here: division by the leading coefficient \(b_n\) requires \(b_n^{-1}\)
to exist.

\emph{Uniqueness.} Suppose \(f = q g + r = q' g + r'\) with
\(\deg r, \deg r' < \deg g\). Subtracting, \((q - q') g = r' - r\). If
\(q \neq q'\), then \(\deg\bigl((q-q')g\bigr) = \deg(q - q') + \deg g \ge
\deg g\) by \eqref{polynomials:eq:deg-prod}; but
\(\deg(r' - r) \le \max\{\deg r', \deg r\} < \deg g\) by
\eqref{polynomials:eq:deg-sum}. These contradict each other, so \(q = q'\), and
then \(r = r'\).
\end{proof}

\begin{remark}[The underlying algorithm]\label{polynomials:rem:long-division}
The existence proof is constructive: it is precisely \emph{polynomial long
division}. We repeatedly eliminate the current leading term by subtracting an
appropriate monomial multiple of \(g\); each step lowers the degree, so the
process terminates after at most \(\deg f - \deg g + 1\) steps. For example,
dividing \(f = X^3 + 2X + 1\) by \(g = X^2 + 1\) over \(\Q\): first subtract
\(X \cdot g = X^3 + X\) to obtain \(X + 1\), whose degree \(1\) is below
\(\deg g = 2\); hence \(q = X\), \(r = X + 1\), and indeed
\(X^3 + 2X + 1 = X(X^2 + 1) + (X + 1)\).
\end{remark}

\begin{definition}[Divisibility]\label{polynomials:def:divisibility}
We say \(g\) \emph{divides} \(f\), written \(g \mid f\), if there is some
\(q \in \F[X]\) with \(f = q g\); equivalently (for \(g \neq 0\)), if the
remainder \(f \bmod g\) is \(0\). In that case \(g\) is a \emph{divisor} or
\emph{factor} of \(f\).
\end{definition}

\subsection{Roots and the factor theorem}
\label{polynomials:subsec:roots}

We now make precise the substitution of an element of \(\F\) for the
indeterminate. This is the first place where a polynomial gives rise to a
function.

\begin{definition}[Evaluation and roots]\label{polynomials:def:evaluation}
For \(f = \sum_{i=0}^n a_i X^i \in \F[X]\) and \(\alpha \in \F\), the
\emph{evaluation} of \(f\) at \(\alpha\) is the field element
\[
  f(\alpha) = \sum_{i=0}^n a_i\, \alpha^i \in \F,
\]
which substitutes \(\alpha\) for the indeterminate \(X\). An element
\(\alpha \in \F\) is a \emph{root} (or \emph{zero}) of \(f\) if
\(f(\alpha) = 0\).
\end{definition}

\begin{proposition}[Evaluation is a ring homomorphism]\label{polynomials:prop:eval-hom}
Fix \(\alpha \in \F\). The \emph{evaluation map}
\(\operatorname{ev}_\alpha \colon \F[X] \to \F\), \(f \mapsto f(\alpha)\),
satisfies, for all \(f, g \in \F[X]\),
\[
  (f + g)(\alpha) = f(\alpha) + g(\alpha),
  \qquad
  (f \cdot g)(\alpha) = f(\alpha)\, g(\alpha),
  \qquad
  1(\alpha) = 1.
\]
That is, \(\operatorname{ev}_\alpha\) is a ring homomorphism.
\end{proposition}

\begin{proof}
Additivity is immediate from coordinatewise addition:
\(\sum_i (a_i + b_i)\alpha^i = \sum_i a_i \alpha^i + \sum_i b_i \alpha^i\). For
multiplicativity, the product \(fg\) has \(k\)-th coefficient
\(c_k = \sum_{i+j=k} a_i b_j\), so
\[
  (fg)(\alpha) = \sum_k c_k \alpha^k
    = \sum_k \sum_{i+j=k} a_i b_j\, \alpha^{i+j}
    = \Bigl(\sum_i a_i \alpha^i\Bigr)\Bigl(\sum_j b_j \alpha^j\Bigr)
    = f(\alpha)\, g(\alpha),
\]
where the middle step expands the product of the two sums and groups terms by
\(k = i + j\). Finally \(1(\alpha) = 1\). This is the formal reason evaluation
respects the algebra: the convolution defining polynomial multiplication
collapses to ordinary multiplication in \(\F\) because \(\alpha^i \alpha^j =
\alpha^{i+j}\).
\end{proof}

The link between roots and divisibility is the factor theorem, an immediate but
powerful consequence of the division algorithm.

\begin{theorem}[Remainder theorem and factor theorem]\label{polynomials:thm:factor}
Let \(f \in \F[X]\) and \(\alpha \in \F\).
\begin{enumerate}
  \item \emph{(Remainder theorem.)} The remainder of \(f\) upon division by the
        linear polynomial \(X - \alpha\) equals the constant \(f(\alpha)\):
        \[
          f = q\,(X - \alpha) + f(\alpha)
          \quad\text{for some } q \in \F[X].
        \]
  \item \emph{(Factor theorem.)} \(\alpha\) is a root of \(f\) if and only if
        \((X - \alpha) \mid f\).
\end{enumerate}
\end{theorem}

\begin{proof}
(1) Divide \(f\) by \(X - \alpha\). Since \(\deg(X - \alpha) = 1\), the
remainder \(r\) has degree \(< 1\), hence is a constant \(r \in \F\). Write
\(f = q(X - \alpha) + r\) and evaluate at \(\alpha\) using
Proposition~\ref{polynomials:prop:eval-hom}:
\[
  f(\alpha) = q(\alpha)\,(\alpha - \alpha) + r = q(\alpha)\cdot 0 + r = r.
\]
Thus \(r = f(\alpha)\).

(2) By part (1), \((X - \alpha) \mid f\) means the remainder \(f(\alpha)\) is
\(0\), i.e.\ \(\alpha\) is a root.
\end{proof}

This single theorem yields the central counting bound of the section.

\begin{theorem}[At most \(d\) roots]\label{polynomials:thm:root-bound}
A nonzero polynomial \(f \in \F[X]\) of degree \(d\) has at most \(d\) distinct
roots in \(\F\).
\end{theorem}

\begin{proof}
Induct on \(d = \deg f\). If \(d = 0\), then \(f\) is a nonzero constant, which
is never \(0\), so it has \(0 \le d\) roots. Suppose \(d \ge 1\) and the claim
holds for all smaller degrees. If \(f\) has no root in \(\F\), then
\(0 \le d\) and the claim holds. Otherwise let \(\alpha\) be a root. By the
factor theorem, \(f = (X - \alpha)\, q\) for some \(q\) with
\(\deg q = d - 1\) (by \eqref{polynomials:eq:deg-prod}). Now let
\(\beta \neq \alpha\) be any root of \(f\). Then
\[
  0 = f(\beta) = (\beta - \alpha)\, q(\beta),
\]
and since \(\beta - \alpha \neq 0\) and \(\F\) has no zero divisors, it follows
that \(q(\beta) = 0\). Thus every root of \(f\) other than \(\alpha\) is a root
of \(q\). By the induction hypothesis \(q\) has at most \(d - 1\) distinct roots,
so \(f\) has at most \((d - 1) + 1 = d\) distinct roots.
\end{proof}

\begin{remark}\label{polynomials:rem:field-needed}
Both the field hypothesis (no zero divisors) and the degree bookkeeping are
essential. Over a ring with zero divisors the bound fails: in
\(\Z/8\Z\) the quadratic \(X^2 - 1\) has the four roots \(1, 3, 5, 7\). This is
why the polynomial methods underlying Halo~2 require a field — and indeed
a finite field \(\Fp\) or \(\Fq\) — beneath them.
\end{remark}

\begin{corollary}[Few points determine a low-degree polynomial]
\label{polynomials:cor:agree-implies-equal}
Let \(f, g \in \F[X]\) both have degree at most \(d\). If \(f(\alpha) =
g(\alpha)\) for \(d + 1\) distinct values \(\alpha \in \F\), then \(f = g\) as
polynomials.
\end{corollary}

\begin{proof}
The difference \(h := f - g\) has \(\deg h \le \max\{\deg f, \deg g\} \le d\) and
vanishes at \(d + 1\) distinct points. If \(h\) were nonzero it would have degree
at most \(d\) and yet \(d + 1\) distinct roots, contradicting
Theorem~\ref{polynomials:thm:root-bound}. Hence \(h = 0\), i.e.\ \(f = g\).
\end{proof}

Corollary~\ref{polynomials:cor:agree-implies-equal} is the deterministic
skeleton of the Schwartz--Zippel principle and of Lagrange interpolation; we
develop both below.

\subsection{Greatest common divisors and the Euclidean algorithm}
\label{polynomials:subsec:gcd}

Because \(\F[X]\) supports division with remainder, the entire theory of greatest
common divisors from elementary number theory transfers verbatim, with degree
replacing magnitude.

\begin{definition}[Greatest common divisor]\label{polynomials:def:gcd}
Let \(f, g \in \F[X]\), not both zero. A \emph{greatest common divisor} of \(f\)
and \(g\) is a polynomial \(d \in \F[X]\) such that
\begin{enumerate}
  \item \(d \mid f\) and \(d \mid g\) (it is a common divisor), and
  \item every common divisor \(c\) of \(f\) and \(g\) satisfies \(c \mid d\)
        (it is divisible by every other common divisor).
\end{enumerate}
A gcd is determined only up to multiplication by a unit; we single out the unique
\emph{monic} gcd and denote it \(\gcd(f, g)\). When \(\gcd(f, g) = 1\) we say
\(f\) and \(g\) are \emph{coprime} or \emph{relatively prime}.
\end{definition}

The existence of such a \(d\), its uniqueness up to associates, and its
expression as an \(\F[X]\)-linear combination of \(f\) and \(g\) constitute the
content of the next two results. The Euclidean algorithm delivers all of this
algorithmically.

\begin{lemma}[Euclidean step]\label{polynomials:lem:euclid-step}
Let \(f, g \in \F[X]\) with \(g \neq 0\), and write \(f = q g + r\) as in the
division algorithm. Then the pair \((f, g)\) and the pair \((g, r)\) have exactly
the same common divisors; in particular their gcds coincide.
\end{lemma}

\begin{proof}
If \(c \mid f\) and \(c \mid g\), then \(c \mid (f - q g) = r\), so \(c\) is a
common divisor of \(g\) and \(r\). Conversely, if \(c \mid g\) and \(c \mid r\),
then \(c \mid (q g + r) = f\), so \(c\) is a common divisor of \(f\) and \(g\).
The two sets of common divisors are equal, hence so are their maximal elements
under divisibility.
\end{proof}

\begin{theorem}[Euclidean algorithm and B\'ezout's identity]
\label{polynomials:thm:euclid}
Let \(f, g \in \F[X]\), not both zero. Then \(\gcd(f, g)\) exists, is unique as a
monic polynomial, and there exist \(s, t \in \F[X]\) with
\begin{equation}\label{polynomials:eq:bezout}
  \gcd(f, g) = s\, f + t\, g.
\end{equation}
Concretely, define a sequence by \(r_0 = f\), \(r_1 = g\), and, as long as
\(r_{i} \neq 0\),
\[
  r_{i-1} = q_i\, r_{i} + r_{i+1},
  \qquad \deg r_{i+1} < \deg r_{i}.
\]
The degrees strictly decrease, so some \(r_{N+1} = 0\); the last nonzero
remainder \(r_N\), made monic, is \(\gcd(f, g)\).
\end{theorem}

\begin{proof}
\emph{Termination and the gcd property.} If \(g = 0\) the gcd is the monic
associate of \(f\); so assume \(g \neq 0\). The remainders satisfy
\(\deg r_1 > \deg r_2 > \cdots \ge 0\) (a strictly decreasing sequence of
nonnegative integers once we pass \(r_1\)), so the algorithm halts with some
\(r_{N+1} = 0\) and \(r_N \neq 0\). Applying
Lemma~\ref{polynomials:lem:euclid-step} repeatedly,
\[
  \{\text{common divisors of } (f, g)\}
  = \{\text{common divisors of } (r_1, r_2)\}
  = \cdots
  = \{\text{common divisors of } (r_N, 0)\}.
\]
The common divisors of \((r_N, 0)\) are exactly the divisors of \(r_N\) (since
everything divides \(0\)), whose greatest element under divisibility is \(r_N\)
itself. Hence \(r_N\) is a gcd of \(f\) and \(g\): it divides both, and any
common divisor divides it. Making it monic gives the distinguished
\(\gcd(f, g)\).

\emph{Uniqueness.} If \(d, d'\) are both monic gcds, then \(d \mid d'\) and
\(d' \mid d\) by property (2), so \(d' = u d\) for a unit \(u\); comparing
leading coefficients of two monic polynomials forces \(u = 1\), whence
\(d = d'\).

\emph{B\'ezout.} We show by induction on the index \(i\) that each \(r_i\) is an
\(\F[X]\)-linear combination of \(f\) and \(g\). Indeed \(r_0 = f\) and
\(r_1 = g\) are such combinations, and the recurrence
\(r_{i+1} = r_{i-1} - q_i r_i\) expresses \(r_{i+1}\) as a combination once
\(r_{i-1}\) and \(r_i\) are. Hence \(r_N = s_0 f + t_0 g\) for some
\(s_0, t_0\); dividing through by \(\operatorname{lc}(r_N)\) to make it monic
yields \eqref{polynomials:eq:bezout}.
\end{proof}

\begin{remark}[Extended Euclidean algorithm]\label{polynomials:rem:eea}
Carrying the linear-combination coefficients along with the remainders — the
\emph{extended Euclidean algorithm} — produces \(s\) and \(t\) explicitly
alongside the gcd. This is precisely how one inverts elements in a quotient
\(\F[X]/(m)\) by an irreducible \(m\): if \(\gcd(f, m) = 1\) then
\(s f + t m = 1\), so \(s\) is the inverse of \(f\) modulo \(m\). The curves of
Halo~2 and Orchard, by contrast, are defined over \emph{prime} fields \(\Fp\),
where inversion uses the integer extended Euclidean algorithm
(Theorem~\ref{intmod:thm:exteuclid}) or Fermat's little theorem
(\cref{sec:finitefields}); the polynomial version above matters because it
parallels the integer one exactly and because it computes gcds of the
polynomials a proof system manipulates.
\end{remark}

\begin{corollary}[{Euclid's lemma in \(\F[X]\)}]\label{polynomials:cor:euclid-lemma}
Let \(p, f, g \in \F[X]\).
\begin{enumerate}
  \item If \(\gcd(p, f) = 1\) and \(p \mid f g\), then \(p \mid g\).
  \item If \(p\) is irreducible (Definition~\ref{polynomials:def:irreducible})
        and \(p \mid f g\), then \(p \mid f\) or \(p \mid g\).
\end{enumerate}
\end{corollary}

\begin{proof}
(1) By B\'ezout there are \(s, t\) with \(s p + t f = 1\). Multiply by \(g\):
\(s p g + t f g = g\). Now \(p\) divides \(s p g\) trivially and divides
\(t f g\) because \(p \mid f g\); hence \(p\) divides their sum \(g\).

(2) If \(p \nmid f\), then \(\gcd(p, f)\) is a monic divisor of \(p\) that is not
the monic associate of \(p\); as \(p\) is irreducible its only monic divisors are
\(1\) and its own monic associate, so \(\gcd(p, f) = 1\). By part (1),
\(p \mid g\).
\end{proof}
\subsection{Irreducible polynomials and unique factorization}
\label{polynomials:subsec:ufd}

\begin{definition}[Irreducible polynomial]\label{polynomials:def:irreducible}
A polynomial \(p \in \F[X]\) is \emph{irreducible} (over \(\F\)) if
\(\deg p \ge 1\) and \(p\) admits no factorization as a product \(p = g h\) with both
\(\deg g \ge 1\) and \(\deg h \ge 1\). Equivalently, \(p\) is nonconstant and its
only factorizations are the trivial ones \(p = (\text{unit})\cdot(\text{associate
of } p)\). A nonconstant polynomial that is not irreducible is \emph{reducible}.
\end{definition}

\begin{remark}[Irreducibility is relative to the field]\label{polynomials:rem:relative}
Irreducibility depends crucially on \(\F\). The polynomial \(X^2 + 1\) is
irreducible over \(\R\) (it has no real root, and a real quadratic factors iff it
has a root), but over \(\C\) it splits as \((X - i)(X + i)\), and over
\(\F = \mathbb{F}_2\) it factors as \((X+1)^2\). The degree-\(1\) polynomials are
irreducible over every field. For later use, observe that a polynomial of degree
\(2\) or \(3\) is irreducible over \(\F\) if and only if it has no root in \(\F\):
any nontrivial factorization of a quadratic or cubic must include a linear
factor, which by the factor theorem supplies a root. This rootlessness test fails
for degree \(\ge 4\), where a factorization into two irreducible quadratics has
no roots yet is reducible.
\end{remark}

Irreducible polynomials are the ``primes'' of \(\F[X]\). We now show that
\(\F[X]\) admits unique factorization into them, exactly as \(\Z\) factors into
primes.

\begin{theorem}[{\(\F[X]\) is a unique factorization domain}]
\label{polynomials:thm:ufd}
Every nonconstant polynomial \(f \in \F[X]\) factors as
\begin{equation}\label{polynomials:eq:factorization}
  f = c\, p_1\, p_2 \cdots p_k,
\end{equation}
where \(c = \operatorname{lc}(f) \in \F^\times\) is a constant and
\(p_1, \dots, p_k\) are monic irreducible polynomials. This factorization is
unique up to the order of the factors: if also \(f = c'\, q_1 \cdots q_\ell\)
with \(c' \in \F^\times\) and each \(q_j\) monic irreducible, then \(c = c'\),
\(k = \ell\), and the \(q_j\) are a reordering of the \(p_i\).
\end{theorem}

\begin{proof}[Proof (sketch)]
The argument is the template of the fundamental theorem of arithmetic
(Theorem~\ref{intmod:thm:fta}), with degree in place of absolute value.
\emph{Existence}: induct on \(\deg f\); an irreducible \(f\) is
\(\operatorname{lc}(f)\) times its monic associate, and a reducible \(f = gh\)
splits into factors of smaller degree, whose factorizations multiply together.
\emph{Uniqueness}: comparing leading coefficients gives
\(c = c' = \operatorname{lc}(f)\); then Euclid's lemma
(Corollary~\ref{polynomials:cor:euclid-lemma}(2)) lets one match \(p_1\) with
some \(q_j\), cancel (Corollary~\ref{polynomials:cor:cancellation}), and recurse,
exactly as in \(\Z\).
\end{proof}

\begin{remark}\label{polynomials:rem:pid}
The structural reason underlying Theorem~\ref{polynomials:thm:ufd} is that
\(\F[X]\) is a \emph{Euclidean domain} (via the degree function and the division
algorithm), hence a \emph{principal ideal domain}: every ideal of \(\F[X]\) is of
the form \((g) = \{q g : q \in \F[X]\}\) for a single generator \(g\), namely a
gcd of any generating set. B\'ezout's identity~\eqref{polynomials:eq:bezout} is
the concrete shadow of this fact. Euclidean \(\Rightarrow\) PID \(\Rightarrow\)
UFD is the standard chain of implications in ring theory; we have just walked it
explicitly for \(\F[X]\).
\end{remark}

\subsection{Lagrange interpolation}
\label{polynomials:subsec:lagrange}

Corollary~\ref{polynomials:cor:agree-implies-equal} established that a polynomial of
degree at most \(d\) is determined by its values at \(d + 1\) points. Lagrange
interpolation is the constructive converse: \emph{any} prescribed values at
\(d + 1\) distinct points arise from a \emph{unique} such polynomial. This
is the algebraic backbone of secret sharing, of Reed--Solomon codes, and of the
polynomial encodings in Halo~2.

\begin{theorem}[Lagrange interpolation]\label{polynomials:thm:lagrange}
Let \(x_0, x_1, \dots, x_d \in \F\) be \(d + 1\) distinct points, and let
\(y_0, y_1, \dots, y_d \in \F\) be arbitrary target values. Then there exists a
\emph{unique} polynomial \(f \in \F[X]\) with \(\deg f \le d\) and
\(f(x_i) = y_i\) for all \(i = 0, 1, \dots, d\). The explicit
formula
\begin{equation}\label{polynomials:eq:lagrange}
  f(X) = \sum_{i=0}^{d} y_i\, \ell_i(X),
  \qquad\text{where}\qquad
  \ell_i(X) = \prod_{\substack{0 \le j \le d \\ j \neq i}}
              \frac{X - x_j}{x_i - x_j}
\end{equation}
gives it. We call the polynomials \(\ell_i\) the \emph{Lagrange basis polynomials} for
the nodes \(x_0, \dots, x_d\).
\end{theorem}

\begin{proof}
\emph{The basis polynomials are well defined.} Each denominator
\(\prod_{j \neq i}(x_i - x_j)\) is a product of nonzero elements of \(\F\)
(distinctness of the nodes), hence is a nonzero element of \(\F\) and is
invertible. Thus each \(\ell_i\) is a genuine polynomial of degree exactly
\(d\) (a product of \(d\) linear factors).

\emph{The interpolation property.} Evaluate \(\ell_i\) at a node \(x_m\). If
\(m = i\), every factor is \((x_i - x_j)/(x_i - x_j) = 1\), so \(\ell_i(x_i) =
1\). If \(m \neq i\), then the factor with \(j = m\) has numerator
\(x_m - x_m = 0\), so \(\ell_i(x_m) = 0\). In short,
\begin{equation}\label{polynomials:eq:lagrange-delta}
  \ell_i(x_m) = \delta_{im} =
  \begin{cases}
    1 & m = i,\\
    0 & m \neq i.
  \end{cases}
\end{equation}
Therefore
\(f(x_m) = \sum_{i} y_i\, \ell_i(x_m) = \sum_i y_i\, \delta_{im} = y_m\) for
each \(m\), as required, and \(\deg f \le \max_i \deg \ell_i = d\).

\emph{Uniqueness.} If \(f\) and \(\tilde f\) both have degree at most \(d\) and
agree at the \(d + 1\) distinct points \(x_0, \dots, x_d\), then by
Corollary~\ref{polynomials:cor:agree-implies-equal} they are equal as
polynomials. Hence the interpolant is unique.
\end{proof}

\begin{example}[A small interpolation]\label{polynomials:exa:lagrange}
Over \(\Q\), find the polynomial of degree \(\le 2\) with \(f(0) = 1\),
\(f(1) = 3\), \(f(2) = 7\). The basis polynomials for nodes \(0, 1, 2\) are
\begin{align*}
  \ell_0 &= \frac{(X-1)(X-2)}{(0-1)(0-2)} = \frac{(X-1)(X-2)}{2}, &
  \ell_1 &= \frac{(X-0)(X-2)}{(1-0)(1-2)} = -X(X-2), \\
  \ell_2 &= \frac{(X-0)(X-1)}{(2-0)(2-1)} = \frac{X(X-1)}{2}.
\end{align*}
Then
\(f = 1 \cdot \ell_0 + 3 \cdot \ell_1 + 7 \cdot \ell_2 = X^2 + X + 1\), and one
verifies \(f(0) = 1\), \(f(1) = 3\), \(f(2) = 7\).
\end{example}

\begin{remark}[The evaluation isomorphism]\label{polynomials:rem:eval-iso}
Fix \(d + 1\) distinct nodes \(x_0, \dots, x_d\) and let \(V = \{f \in \F[X] :
\deg f \le d\}\), an \((d{+}1)\)-dimensional \(\F\)-vector space (basis
\(1, X, \dots, X^d\)). The \emph{evaluation map}
\[
  \operatorname{ev} \colon V \to \F^{d+1},
  \qquad f \mapsto \bigl(f(x_0), \dots, f(x_d)\bigr),
\]
is \(\F\)-linear. Surjectivity is exactly the existence half of Lagrange
interpolation, and injectivity is exactly the uniqueness half; together they say
\(\operatorname{ev}\) is a \emph{linear isomorphism}
\(V \cong \F^{d+1}\). The Lagrange basis \(\{\ell_i\}\) is precisely the basis of
\(V\) dual to the standard basis of \(\F^{d+1}\), by
\eqref{polynomials:eq:lagrange-delta}. This isomorphism — switching between a
polynomial's \emph{coefficient representation} and its \emph{evaluation
representation} — is the conceptual core of the fast (FFT-based) polynomial
arithmetic and the evaluation-domain encodings used throughout Halo~2.
\end{remark}

\begin{remark}[The vanishing polynomial]\label{polynomials:rem:vanishing}
The product of all the linear factors over the node set,
\[
  Z(X) = \prod_{i=0}^{d} (X - x_i),
\]
is the unique monic polynomial of degree \(d + 1\) that vanishes at every node;
we call it the \emph{vanishing polynomial} of the set
\(\{x_0, \dots, x_d\}\). A polynomial \(g\) vanishes on the whole node set if and
only if \(Z \mid g\) — one direction is the factor theorem applied repeatedly
(using that the distinct \(X - x_i\) are pairwise coprime), the other is
immediate. The pair (Lagrange basis, vanishing polynomial) is exactly the data a
proof system manipulates when it argues that a committed polynomial agrees with
prescribed values on an evaluation domain.
\end{remark}

\subsection{The formal derivative and multiplicities}
\label{polynomials:subsec:derivative}

Roots can occur ``with multiplicity,'' and to detect repeated roots without any
appeal to limits we introduce a purely algebraic derivative.

\begin{definition}[Multiplicity of a root]\label{polynomials:def:multiplicity}
Let \(f \in \F[X]\) be nonzero and \(\alpha \in \F\) a root. The
\emph{multiplicity} of \(\alpha\) (as a root of \(f\)) is the largest integer
\(m \ge 1\) such that \((X - \alpha)^m \mid f\). We write \(m = \operatorname{mult}_\alpha(f)\).
A root of multiplicity \(1\) is \emph{simple}; a root of multiplicity \(\ge 2\)
is \emph{repeated} (or \emph{multiple}).
\end{definition}

This is well defined: the powers \((X - \alpha)^m\) have strictly increasing
degree \(m\), so only finitely many can divide the fixed polynomial \(f\); the
largest such \(m\) exists. Equivalently, \(m\) is determined by writing
\(f = (X - \alpha)^m\, g\) with \(g(\alpha) \neq 0\).

\begin{proposition}[Counting roots with multiplicity]
\label{polynomials:prop:roots-with-mult}
If \(f \in \F[X]\) is nonzero of degree \(d\) with distinct roots
\(\alpha_1, \dots, \alpha_r \in \F\) of multiplicities \(m_1, \dots, m_r\), then
\[
  f = (X - \alpha_1)^{m_1} \cdots (X - \alpha_r)^{m_r}\, g
\]
for some \(g \in \F[X]\) with no roots in \(\F\), and
\(\sum_{i=1}^r m_i \le d\). In particular the number of roots \emph{counted with
multiplicity} is at most \(d\), with equality iff \(f\) splits into linear
factors over \(\F\).
\end{proposition}

\begin{proof}
The distinct linear polynomials \(X - \alpha_i\) are pairwise coprime (their gcd
is \(1\), as a common divisor would be a nonzero constant by degree). Hence if
\((X - \alpha_i)^{m_i} \mid f\) for each \(i\), their product divides \(f\): this
follows by induction using that if coprime \(a, b\) both divide \(f\) then
\(ab \mid f\) (write \(f = a u = b v\); since \(\gcd(a,b)=1\), B\'ezout gives
\(sa + tb = 1\), so \(f = saf + tbf = sa(bv) + tb(au) = ab(sv + tu)\)). Thus
\(f = \prod_i (X - \alpha_i)^{m_i}\, g\); by maximality of each \(m_i\),
\(g(\alpha_i) \neq 0\), and \(g\) has no other roots either (any root of \(g\)
would be a root of \(f\), hence some \(\alpha_i\)). Taking degrees,
\(\sum_i m_i + \deg g = d\), so \(\sum_i m_i \le d\), with equality iff
\(\deg g = 0\), i.e.\ \(g\) is a nonzero constant and \(f\) is a product of
linear factors.
\end{proof}

\begin{definition}[Formal derivative]\label{polynomials:def:derivative}
The \emph{formal derivative} is the map \(D \colon \F[X] \to \F[X]\) defined on
\(f = \sum_{i=0}^n a_i X^i\) by
\[
  f'(X) = D f = \sum_{i=1}^{n} i\, a_i\, X^{i-1}
        = a_1 + 2 a_2 X + 3 a_3 X^2 + \cdots + n\, a_n X^{n-1},
\]
where \(i\, a_i\) means \(a_i\) added to itself \(i\) times in \(\F\) (the image
of the integer \(i\) under \(\Z \to \F\)). No limits are involved; this is a
formal, algebraic definition that makes sense over any field.
\end{definition}

\begin{proposition}[Rules of formal differentiation]\label{polynomials:prop:deriv-rules}
For all \(f, g \in \F[X]\) and \(c \in \F\):
\begin{enumerate}
  \item \emph{(Linearity.)} \((f + g)' = f' + g'\) and \((c f)' = c f'\).
  \item \emph{(Product rule / Leibniz.)} \((f g)' = f' g + f g'\).
  \item \emph{(Power rule.)} \(\bigl((X - \alpha)^m\bigr)' = m\,(X - \alpha)^{m-1}\)
        for every \(\alpha \in \F\) and \(m \ge 1\).
\end{enumerate}
\end{proposition}

\begin{proof}
(1) is immediate from the coordinatewise, linear definition of \(D\).

(2) By linearity it suffices to prove the product rule on monomials \(f = X^a\),
\(g = X^b\), since both sides are bilinear in \((f, g)\) and every polynomial is
a linear combination of monomials. For monomials, \((X^a X^b)' = (X^{a+b})' =
(a+b) X^{a+b-1}\), while \(f' g + f g' = a X^{a-1} X^b + X^a\, b X^{b-1} =
(a + b) X^{a+b-1}\). The two agree. Extending by bilinearity: writing
\(f = \sum_a u_a X^a\), \(g = \sum_b v_b X^b\), we have \((fg)' =
\sum_{a,b} u_a v_b (X^a X^b)' = \sum_{a,b} u_a v_b\,(X^a)'X^b +
\sum_{a,b} u_a v_b\, X^a (X^b)' = f' g + f g'\).

(3) Induct on \(m\) using the product rule. For \(m = 1\),
\((X - \alpha)' = 1 = 1\cdot(X-\alpha)^0\). Assuming the claim for \(m\),
\[
  \bigl((X-\alpha)^{m+1}\bigr)' = \bigl((X-\alpha)^m (X - \alpha)\bigr)'
   = m(X-\alpha)^{m-1}(X-\alpha) + (X-\alpha)^m \cdot 1
   = (m+1)(X-\alpha)^m. \qedhere
\]
\end{proof}

\begin{theorem}[Derivative test for multiple roots]\label{polynomials:thm:multiple-root}
Let \(f \in \F[X]\) be nonzero and \(\alpha \in \F\). Then \(\alpha\) is a
repeated root of \(f\) (multiplicity \(\ge 2\)) if and only if
\(f(\alpha) = 0\) and \(f'(\alpha) = 0\). Consequently \(f\) has a repeated root
in \(\F\) only if \(\gcd(f, f') \neq 1\); and \(f\) has no repeated roots in any
field containing \(\F\) (it is \emph{separable}) if and only if
\(\gcd(f, f') = 1\) in \(\F[X]\).
\end{theorem}

\begin{proof}
Write \(\operatorname{mult}_\alpha(f) = m \ge 1\), so \(f = (X - \alpha)^m g\)
with \(g(\alpha) \neq 0\). By the product and power rules,
\[
  f' = m (X - \alpha)^{m-1} g + (X - \alpha)^m g'
     = (X - \alpha)^{m-1}\bigl[m\, g + (X - \alpha) g'\bigr].
\]
Evaluate the bracket at \(\alpha\): it equals \(m\, g(\alpha) + 0 = m\, g(\alpha)\).

If \(m = 1\): then \(f' = g + (X-\alpha)g'\), so \(f'(\alpha) = g(\alpha) \neq 0\).
Thus a simple root is \emph{not} a root of \(f'\).

If \(m \ge 2\): then \((X-\alpha)^{m-1}\) with \(m - 1 \ge 1\) divides \(f'\),
so \(f'(\alpha) = 0\); and of course \(f(\alpha) = 0\).

Hence \(f(\alpha) = f'(\alpha) = 0\) exactly when \(m \ge 2\), proving the first
claim. For the gcd statements: a common root \(\alpha \in \F\) of \(f\) and
\(f'\) forces \((X - \alpha) \mid \gcd(f, f')\), so a repeated root makes the gcd
nontrivial. The full separability statement requires care over extensions: one
shows \(\gcd(f, f')\), computed in \(\F[X]\), is unchanged under field extension
(the Euclidean algorithm uses only field operations in \(\F\)), and that over a
splitting field a repeated factor \((X-\alpha)^m\), \(m\ge 2\), contributes a
common factor \((X-\alpha)\) to \(f\) and \(f'\) by the computation above, while
distinct simple linear factors contribute none. Thus \(\gcd(f, f') = 1\) in
\(\F[X]\) if and only if \(f\) has no repeated root in any extension.
\end{proof}

\begin{remark}[The pitfall of positive characteristic]\label{polynomials:rem:char-p}
The integer coefficient \(i\) in \(f' = \sum_i i a_i X^{i-1}\) is interpreted
\emph{in \(\F\)}, and may vanish there. Over a field of characteristic \(p\)
(meaning \(p = 0\) in \(\F\), with \(p\) prime — the setting of \(\Fp\) and
\(\Fq\)), the derivative of \(X^p\) is \(p X^{p-1} = 0\). Thus \(f = X^p - 1\) over
\(\Fp\) has \(f' = 0\), and \(\gcd(f, f') = f \neq 1\); indeed
\(X^p - 1 = (X - 1)^p\) in \(\Fp[X]\), a single root of multiplicity \(p\). This
characteristic-\(p\) behavior is not a pathology to avoid but a structural
feature one must respect when working over the finite fields of Orchard.
\end{remark}
\subsection{Polynomials as formal objects versus functions}
\label{polynomials:subsec:formal-vs-function}

This section sharpens the distinction noted at the outset, and extracts from it
the probabilistic principle on which the remainder of this book rests.

\begin{definition}[Polynomial function]\label{polynomials:def:poly-function}
Every \(f \in \F[X]\) induces a \emph{polynomial function}
\(\widetilde f \colon \F \to \F\), \(\alpha \mapsto f(\alpha)\). The assignment
\(f \mapsto \widetilde f\) is a ring homomorphism from \(\F[X]\) to the ring
\(\F^\F\) of all functions \(\F \to \F\) (with pointwise operations), by
Proposition~\ref{polynomials:prop:eval-hom} applied at every point.
\end{definition}

The central question is whether two \emph{distinct} polynomials can induce the
\emph{same} function. Over an infinite field they cannot; over a finite field
they can.

\begin{proposition}[When formal equals functional]
\label{polynomials:prop:formal-vs-fn}
Let \(f, g \in \F[X]\).
\begin{enumerate}
  \item If \(\F\) is infinite, then \(\widetilde f = \widetilde g\) (equal as
        functions) implies \(f = g\) (equal as polynomials). Equivalently, the
        map \(f \mapsto \widetilde f\) is injective, and one may identify
        polynomials with polynomial functions.
  \item If \(\F\) is finite, the map \(f \mapsto \widetilde f\) is \emph{not}
        injective. For instance, over \(\Fp\) the nonzero polynomial
        \(X^p - X\) induces the zero function, since by Fermat's little theorem
        (Corollary~\ref{groups:cor-euler-fermat})
        \(\alpha^p = \alpha\) for every \(\alpha \in \Fp\).
\end{enumerate}
\end{proposition}

\begin{proof}
(1) Suppose \(\widetilde f = \widetilde g\). Then \(h := f - g\) satisfies
\(h(\alpha) = 0\) for every \(\alpha \in \F\). If \(h\) were nonzero it would
have some finite degree \(d\) and thus at most \(d\) roots by
Theorem~\ref{polynomials:thm:root-bound}; but \(\F\) is infinite, so \(h\) has
infinitely many roots — a contradiction. Hence \(h = 0\) and \(f = g\).

(2) The polynomial \(X^p - X \in \Fp[X]\) is nonzero (its coefficients are
\(-1\) at \(X^1\) and \(1\) at \(X^p\)), yet by Fermat's little theorem
\(\alpha^p = \alpha\), so \((X^p - X)(\alpha) = 0\) for all \(\alpha \in \Fp\).
Thus \(X^p - X\) and the zero polynomial induce the same (zero) function while
being distinct formal objects, so the map is not injective.
\end{proof}

\begin{remark}[Why the formal viewpoint is the correct one]
\label{polynomials:rem:formal-right}
Proposition~\ref{polynomials:prop:formal-vs-fn}(2) is precisely the reason
we define polynomials as coefficient sequences rather than as functions. In
cryptography one works over finite fields \(\Fp, \Fq\), where the functional
viewpoint loses information; the \emph{degree} of a polynomial — a property of
the formal object, invisible to the induced function once the field is finite —
is what the security of polynomial commitment schemes rests on. A prover commits
to a low-degree formal polynomial; the verifier tests it by evaluation. The gap
between agreement as a function at the tested points and identity to the claimed
formal polynomial is precisely the gap that the next principle quantifies.
\end{remark}

\begin{theorem}[Univariate Schwartz--Zippel]\label{polynomials:thm:schwartz-zippel}
Let \(f \in \F[X]\) be a nonzero polynomial of degree at most \(d\), and let
\(S \subseteq \F\) be a finite, nonempty set. If \(\alpha\) is chosen uniformly
at random from \(S\), then
\begin{equation}\label{polynomials:eq:schwartz-zippel}
  \Pr_{\alpha \sim S}\bigl[\, f(\alpha) = 0 \,\bigr] \le \frac{d}{|S|}.
\end{equation}
More generally, if \(f, g \in \F[X]\) are distinct polynomials each of degree at
most \(d\), then \(\Pr_{\alpha \sim S}[\, f(\alpha) = g(\alpha)\,] \le d / |S|\).
\end{theorem}

\begin{proof}
By Theorem~\ref{polynomials:thm:root-bound}, the nonzero polynomial \(f\) of
degree at most \(d\) has at most \(d\) roots in \(\F\), hence at most \(d\) roots
inside \(S\). The number of \(\alpha \in S\) with \(f(\alpha) = 0\) is therefore
at most \(d\), and since \(\alpha\) is uniform on \(S\),
\[
  \Pr_{\alpha \sim S}[f(\alpha) = 0]
    = \frac{|\{\alpha \in S : f(\alpha) = 0\}|}{|S|}
    \le \frac{d}{|S|}.
\]
For the general statement, apply this to the nonzero polynomial \(h = f - g\),
which has degree at most \(d\); \(f(\alpha) = g(\alpha)\) iff \(h(\alpha) = 0\).
\end{proof}

\begin{remark}[The principle stated plainly, and its application]
\label{polynomials:rem:sz-use}
Inequality~\eqref{polynomials:eq:schwartz-zippel} formalizes the statement that
\emph{two low-degree polynomials that agree at many points must be equal.} The
deterministic version is Corollary~\ref{polynomials:cor:agree-implies-equal}:
agreement at \(d + 1\) points forces equality. Schwartz--Zippel is its
probabilistic refinement: even a \emph{single} random evaluation point detects
unequal degree-\(\le d\) polynomials except with probability at most
\(d / |S|\). If \(|S|\) is exponentially large (as for \(S = \Fq\) with
\(q \approx 2^{255}\), the field sizes in Orchard) and \(d\) is polynomially
bounded, this failure probability is negligible. This constitutes the foundational
soundness argument for polynomial identity testing, and through it for the
polynomial commitment schemes and arithmetization checks at the heart of
Halo~2: to verify a purported polynomial identity, one evaluates both sides at a
random challenge and compares. We invoke this principle repeatedly in the
analysis of those protocols.
\end{remark}

\begin{remark}[Looking ahead: multivariate generalization]\label{polynomials:rem:multivariate}
The full Schwartz--Zippel lemma extends \eqref{polynomials:eq:schwartz-zippel} to
multivariate polynomials \(f \in \F[X_1, \dots, X_n]\): a nonzero polynomial of
total degree \(d\) vanishes at a uniformly random point of \(S^n\) with
probability at most \(d / |S|\). The univariate case proved here is both the base
case of its inductive proof and the version most directly used in the analysis of
Halo~2; the later volumes develop the multivariate ring \(\F[X_1, \dots, X_n]\)
and this generalization where they require it.
\end{remark}

\clearpage

\section{Finite fields}\label{sec:finitefields}

A \emph{field} is a set in which one can add, subtract, multiply, and divide by anything nonzero, with all the usual algebraic laws holding. The familiar examples \(\Q\), \(\R\), and \(\C\) are all infinite. This section develops the theory of \emph{finite} fields: fields with only finitely many elements. These objects are the arithmetic substrate of essentially all modern public-key cryptography, and in particular of the elliptic-curve constructions underlying Halo~2 and Zcash Orchard. There the coordinates of curve points live in a prime field \(\Fp\), while the scalars that multiply those points live in a second prime field \(\Fq\); understanding both, and the relationship between them, begins with understanding finite fields in general.

The ring theory of \cref{sec:ringsfields} and the polynomial theory of \cref{sec:polynomials} are assumed; the emphasis here is on the structure theory specific to the finite case.

\subsection{The prime field \texorpdfstring{\(\Fp\)}{Fp} and modular arithmetic}
\label{finitefields:subsec:primefield}

The basic example is already in hand: by Corollary~\ref{intmod:cor:fpfield}, for a prime \(p\) the ring \(\Z/p\Z\) is a field, the \emph{prime field} \(\Fp\), because every \(a \in \{1,\dots,p-1\}\) satisfies \(\gcd(a,p)=1\) and is therefore invertible (Theorem~\ref{intmod:thm:unitiff}). Its arithmetic is integer arithmetic followed by reduction modulo \(p\). Division is the one operation whose implementation deserves comment, and two methods coexist in practice.

\paragraph{Inversion via the extended Euclidean algorithm.}
The extended Euclidean algorithm (Theorem~\ref{intmod:thm:exteuclid}) on input \((a, p)\) with \(0 < a < p\) returns \(u\) with \(ua \equiv 1 \pmod p\), so \(\overline a^{-1} = \overline u\), in \(O(\log p)\) division steps; Example~\ref{intmod:exa:modinverse} carried this out. This is the standard general-purpose method.

\paragraph{Inversion via Fermat's little theorem.}
By Fermat's little theorem (Corollary~\ref{groups:cor-euler-fermat}), every nonzero \(\overline a \in \Fp\) satisfies \(\overline a^{\,p-1} = \overline 1\), so \(\overline a \cdot \overline a^{\,p-2} = \overline 1\) and
\[
  \overline a^{-1} = \overline a^{\,p-2}.
\]
The exponentiation is computed by \emph{square-and-multiply}: to raise to an exponent \(e\) with binary expansion \(e = \sum_i e_i 2^i\), square repeatedly to obtain \(a^{2^i}\) and multiply together those factors with \(e_i = 1\), for a total of \(O(\log e)\) multiplications rather than \(e - 1\). (Double-and-add, \cref{ellipticcurves:def-dbladd}, is the same algorithm written additively for elliptic-curve scalar multiplication.) Inversion thus costs \(O(\log p)\) modular multiplications. The method is constant-time-friendly --- the sequence of operations does not branch on the secret \(a\) --- which is why implementations that must resist side-channel attacks often prefer this formula to the extended Euclidean algorithm.

\subsection{Characteristic, prime subfield, and order}
\label{finitefields:subsec:char}

Let \(F\) be any field and recall from \cref{ringsfields:subsec-char} the ring homomorphism \(\iota\colon \Z \to F\), \(\iota(n) = n\cdot 1_F\), and the \emph{characteristic} \(\operatorname{char} F\), the nonnegative generator of \(\ker\iota\) (Definition~\ref{ringsfields:def-char}). If \(F\) is finite, \(\iota\) cannot be injective (its domain is infinite), so \(\operatorname{char} F = p > 0\), and \(p\) is prime by Theorem~\ref{ringsfields:thm-char-prime}. For \(F = \Fp\) itself, \(\operatorname{char} \Fp = p\).

\begin{definition}[Prime subfield]
\label{finitefields:def:primesubfield}
The \emph{prime subfield} of \(F\) is the smallest subfield of \(F\), namely the intersection of all subfields, equivalently the subfield generated by \(1_F\).
\end{definition}

\begin{proposition}
\label{finitefields:prop:primesub}
If \(\operatorname{char} F = p > 0\), the prime subfield of \(F\) is isomorphic to \(\Fp\). If \(\operatorname{char} F = 0\), it is isomorphic to \(\Q\).
\end{proposition}

\begin{proof}
Consider \(\iota\colon \Z \to F\) as above. When \(\operatorname{char} F = p\), the kernel is \(p\Z\), so \(\iota\) factors through an injective ring homomorphism \(\overline\iota\colon \Z/p\Z \to F\), whose image is a subfield isomorphic to \(\Fp\). Any subfield contains \(1_F\) and hence this image, so it is the prime subfield. When \(\operatorname{char} F = 0\), \(\iota\) is injective and extends to an embedding of \(\Q\) by sending \(a/b \mapsto (a\cdot 1_F)(b\cdot 1_F)^{-1}\); its image is the prime subfield.
\end{proof}

Thus every finite field \(F\) is a vector space over its prime subfield \(\Fp\). Since \(F\) is finite it is finite-dimensional; if \(\dim_{\Fp} F = k\), then choosing a basis \(e_1, \dots, e_k\) shows that every element of \(F\) is uniquely \(c_1 e_1 + \cdots + c_k e_k\) with \(c_i \in \Fp\), so \(F\) has exactly \(p^k\) elements (basis and dimension are developed formally in \cref{sec:linalg}; the argument needs only basis existence and uniqueness of coordinates).

\begin{theorem}[Order of a finite field is a prime power]
\label{finitefields:thm:primepower}
Every finite field has cardinality \(p^k\) for some prime \(p\) and integer \(k \ge 1\), where \(p = \operatorname{char} F\).
\end{theorem}

\begin{proof}
As shown above, \(p = \operatorname{char} F\) is prime, and by Proposition~\ref{finitefields:prop:primesub} the prime subfield is \(\Fp\). As just observed, \(F\) is a finite-dimensional \(\Fp\)-vector space, say of dimension \(k \ge 1\), and counting coordinate vectors gives \(|F| = p^k\).
\end{proof}

Henceforth \(q = p^k\) denotes a prime power and \(\Fq\) a field with \(q\) elements. A field of every prime-power order exists, and any two fields of the same order are isomorphic (one constructs \(\Fq\) as a quotient \(\Fp[x]/(f)\) by an irreducible polynomial of degree \(k\), cf.\ Example~\ref{ringsfields:exa-quotient-Zn}); we do not prove these facts, since the fields of Halo~2 and Orchard are the prime fields \(q = p\), where \cref{intmod:cor:fpfield} already settles everything. We state results for general \(\Fq\) when the proofs are identical.

\subsection{The Frobenius map and the freshman's dream}
\label{finitefields:subsec:frobenius}

A recurring phenomenon of prime characteristic is that the \(p\)-th power map respects addition.

\begin{lemma}[Freshman's dream]
\label{finitefields:lem:freshman}
Let \(R\) be a commutative ring of prime characteristic \(p\). Then for all \(a, b \in R\),
\[
  (a + b)^p = a^p + b^p,
  \qquad\text{and more generally}\qquad
  (a + b)^{p^n} = a^{p^n} + b^{p^n}.
\]
\end{lemma}

\begin{proof}
The binomial theorem gives \((a+b)^p = \sum_{i=0}^{p} \binom{p}{i} a^i b^{p-i}\). For \(0 < i < p\), the integer \(\binom{p}{i} = \frac{p!}{i!\,(p-i)!}\) is divisible by \(p\): the numerator contains the factor \(p\), while the denominator, being a product of integers all smaller than \(p\), is coprime to \(p\). Since \(p \cdot 1_R = 0\), every middle term vanishes in \(R\), leaving \((a+b)^p = a^p + b^p\). The general case follows by induction on \(n\), applying the case \(n=1\) to \(a^{p^{n-1}}\) and \(b^{p^{n-1}}\).
\end{proof}

\begin{definition}[Frobenius endomorphism]
\label{finitefields:def:frobenius}
Let \(F\) be a field of characteristic \(p\). The \emph{Frobenius map} is
\[
  \operatorname{Frob}\colon F \to F, \qquad \operatorname{Frob}(x) = x^p.
\]
\end{definition}

\begin{proposition}
\label{finitefields:prop:frobhom}
The Frobenius map of a field \(F\) of characteristic \(p\) is a ring homomorphism fixing the prime subfield \(\Fp\) pointwise. If \(F\) is finite it is an automorphism of \(F\).
\end{proposition}

\begin{proof}
It preserves products since \((xy)^p = x^p y^p\) by commutativity, and sums by Lemma~\ref{finitefields:lem:freshman}; it sends \(1 \mapsto 1\). Thus it is a ring homomorphism, and any ring homomorphism of fields is injective: if \(x \neq 0\) lay in the kernel, then \(1 = x^{-1}x\) would map to \(0\), contradicting \(1 \mapsto 1\). On \(\Fp\), Fermat's little theorem (Corollary~\ref{groups:cor-euler-fermat}) gives \(x^p = x\) for all \(x\), so \(\operatorname{Frob}\) fixes \(\Fp\). When \(F\) is finite, an injective self-map of a finite set is bijective, so \(\operatorname{Frob}\) is an automorphism.
\end{proof}
\subsection{The multiplicative group of a finite field is cyclic}
\label{finitefields:subsec:cyclic}

For a field \(F\) write \(F^\times = F \setminus \{0\}\) for its multiplicative group. When \(F\) is finite, \(F^\times\) has a strikingly simple structure: it is cyclic, that is, generated by a single element. This fact underlies discrete-logarithm cryptography and the construction of finite fields with prescribed roots of unity used in Halo~2. We prove it in the natural generality of an arbitrary finite subgroup of \(F^\times\); the proof rests on the root bound for polynomials (Theorem~\ref{polynomials:thm:root-bound}), the order theory of \cref{sec:groups}, and one counting identity for Euler's totient \(\varphi\) (the order of \((\Z/n\Z)^\times\), Proposition~\ref{intmod:prop:unitgroup}).

\begin{lemma}[Gauss's divisor-sum identity]
\label{finitefields:lem:totient-sum}
For every \(n \ge 1\), \(\sum_{d \mid n} \varphi(d) = n\).
\end{lemma}

\begin{proof}
Partition \(\{1,2,\dots,n\}\) according to the value \(g = \gcd(a,n)\). For a fixed divisor \(d \mid n\), the integers \(a\) with \(\gcd(a,n)=d\) are exactly \(a = d b\) with \(1 \le b \le n/d\) and \(\gcd(b, n/d)=1\); there are \(\varphi(n/d)\) of these. Hence \(n = \sum_{d \mid n}\varphi(n/d)\), and as \(d\) ranges over divisors of \(n\) so does \(n/d\), giving \(\sum_{d \mid n}\varphi(d)=n\).
\end{proof}

\begin{theorem}[Finite subgroups of \(F^\times\) are cyclic]
\label{finitefields:thm:cyclic}
Let \(F\) be any field and let \(G \le F^\times\) be a finite subgroup of order \(N\). Then \(G\) is cyclic. In particular, for a finite field \(F\) with \(q\) elements, \(F^\times\) is cyclic of order \(q-1\).
\end{theorem}

\begin{proof}
For each divisor \(d \mid N\), let \(\psi(d)\) be the number of elements of \(G\) of order exactly \(d\). Every element's order divides \(N\) (Corollary~\ref{groups:cor-order-elt-divides}), so
\[
\sum_{d \mid N} \psi(d) = N.
\]
We claim that for each \(d \mid N\), either \(\psi(d) = 0\) or \(\psi(d) = \varphi(d)\). Suppose \(\psi(d) > 0\), so some \(a \in G\) has order \(d\). The \(d\) distinct powers \(a^0, a^1, \dots, a^{d-1}\) are all roots of \(X^d - 1\) in \(F\). But \(X^d - 1\) has at most \(d\) roots (Theorem~\ref{polynomials:thm:root-bound}), so these powers are \emph{all} the roots, and every element of \(G\) of order \(d\) is among them. Among the powers \(a^j\), those of order exactly \(d\) are the \(a^j\) with \(\gcd(j,d)=1\) by Corollary~\ref{groups:cor-order-power}, and there are \(\varphi(d)\) of them. Hence \(\psi(d) = \varphi(d)\) whenever \(\psi(d)>0\).

Therefore \(\psi(d) \le \varphi(d)\) for all \(d \mid N\). Summing and using Lemma~\ref{finitefields:lem:totient-sum},
\[
N = \sum_{d \mid N}\psi(d) \;\le\; \sum_{d \mid N}\varphi(d) = N.
\]
Equality throughout forces \(\psi(d) = \varphi(d)\) for every \(d \mid N\). In particular \(\psi(N) = \varphi(N) \ge 1\), so \(G\) contains an element of order \(N = |G|\), which generates it. For a finite field, take \(G = F^\times\), of order \(q-1\).
\end{proof}

The count \(\psi(d) = \varphi(d)\) established along the way is itself useful: a cyclic group of order \(N\) contains exactly \(\varphi(d)\) elements of order \(d\) for each \(d \mid N\), and in particular \(\varphi(N)\) generators.

\begin{definition}[Primitive root]
\label{finitefields:def:primitive}
We call a generator of \(\Fq^\times\) a \emph{primitive root} (or \emph{primitive element}) of \(\Fq\). Equivalently, \(g\) is primitive iff \(\ord(g) = q-1\).
\end{definition}

\begin{example}
\label{finitefields:ex:primitiveroot}
In \(\F_7^\times\), of order \(6\), test \(g = 3\): the powers are \(3^1 = 3,\ 3^2 = 2,\ 3^3 = 6,\ 3^4 = 4,\ 3^5 = 5,\ 3^6 = 1\) (mod \(7\)), running through all of \(\{1,\dots,6\}\). So \(3\) is a primitive root. By contrast \(2\) has order \(3\) (\(2,4,1\)) and is not primitive. By Corollary~\ref{groups:cor-order-power} the remaining primitive roots are the powers \(3^j\) with \(\gcd(j,6)=1\), namely \(3^1 = 3\) and \(3^5 = 5\); indeed there are \(\varphi(6) = 2\) of them.
\end{example}

\begin{corollary}[Roots of unity]
\label{finitefields:cor:rootsofunity}
Let \(F\) be a finite field with \(q\) elements. For every \(n \ge 1\), the solutions of \(x^n = 1\) in \(F^\times\) number exactly \(\gcd(n, q-1)\), and they form the cyclic subgroup of that order. In particular, when \(n \mid q-1\) there are exactly \(n\) such solutions.
\end{corollary}

\begin{proof}
Let \(g\) be a primitive root, so \(F^\times = \langle g\rangle\) is cyclic of order \(q-1\). Writing \(x = g^j\), the equation \(x^n = 1\) becomes \(g^{jn} = 1\), i.e.\ \((q-1) \mid jn\), i.e.\ \(\frac{q-1}{d} \mid j\) where \(d = \gcd(n, q-1)\). The values \(j \in \{0, 1, \dots, q-2\}\) satisfying this are \(j = \frac{q-1}{d}\cdot t\) for \(t = 0, \dots, d-1\), exactly \(d\) of them, and they form the cyclic subgroup generated by \(g^{(q-1)/d}\). When \(n \mid q-1\) we have \(d = n\).
\end{proof}

This corollary is precisely what renders a curve such as Pallas or Vesta suitable for Halo~2: one chooses \(q\) so that \(q - 1\) is divisible by a large power of \(2\), guaranteeing a multiplicative subgroup of \(2\)-power order on which the number-theoretic transform, and hence efficient polynomial arithmetic, operates.

\subsection{Prime fields in elliptic-curve cryptography}
\label{finitefields:subsec:ecc}

This section closes by indicating how the prime field \(\Fp\) — the simplest finite field — enters elliptic-curve cryptography, the setting of Halo~2 and Zcash Orchard. We develop the geometry of elliptic curves in full later; here we only locate the relevant fields.

An elliptic curve over \(\Fp\) (with \(p > 3\)) typically takes the short Weierstrass form, the set of solutions
\[
  E(\Fp) = \{\, (x, y) \in \Fp \times \Fp : y^2 = x^3 + ax + b \,\} \cup \{\mathcal O\},
\]
together with a distinguished point \(\mathcal O\) at infinity, where \(a, b \in \Fp\) and \(4a^3 + 27b^2 \neq 0\). All the arithmetic defining the group law — the computation of slopes \(\lambda = (y_2 - y_1)(x_2 - x_1)^{-1}\), squaring, and the inversions required to add and double points — proceeds in \(\Fp\) using exactly the operations of Subsection~\ref{finitefields:subsec:primefield}. We call this \(\Fp\) the \emph{base field} (or coordinate field) of the curve.

The set \(E(\Fp)\) forms a finite abelian group under the chord-and-tangent addition law, of order \(N = |E(\Fp)|\) close to \(p\) (the Hasse bound, Theorem~\ref{ellipticcurves:thm-hasse}). For cryptography one selects a curve whose group has a large prime factor \(r\); the subgroup of order \(r\) is where keys and signatures reside. The scalars acting on this subgroup live modulo \(r\), and because \(r\) is prime, these scalars themselves form a field.

\begin{definition}[Base field and scalar field]
\label{finitefields:def:basescalar}
For an elliptic curve used in cryptography with a prime-order subgroup of order \(r\):
\begin{itemize}
\item the \emph{base field} \(\Fp\) is the field over which the curve is defined and in which point coordinates lie;
\item the \emph{scalar field} \(\Fr = \Z/r\Z\) is the field of scalars by which one multiplies points, i.e.\ the ring of integers modulo the subgroup order \(r\).
\end{itemize}
Both are prime fields, but in general \(p \neq r\).
\end{definition}

\begin{remark}[Two fields, and the Pasta cycle]
\label{finitefields:rem:pasta}
Scalar multiplication \(P \mapsto sP\) on a point \(P\) of order \(r\) depends only on \(s \bmod r\), because \(rP = \mathcal O\); the exponents thus lie in \(\Fr\), while the coordinates of \(sP\) lie in \(\Fp\). A typical elliptic-curve operation therefore manipulates two distinct prime fields simultaneously. Halo~2's Pallas and Vesta curves (the ``Pasta'' cycle) exploit this deliberately: they form a \emph{cycle of curves}, meaning the base field of one is the scalar field of the other and vice versa,
\[
  |E_{\text{Pallas}}| = p_{\text{Vesta}}, \qquad |E_{\text{Vesta}}| = p_{\text{Pallas}},
\]
so that the scalar field of Pallas equals the base field of Vesta. This permits a proof system to express the verifier of one curve's arithmetic efficiently inside the native field of the other, enabling the recursive proof composition at the heart of Halo~2. One chooses both \(p_{\text{Pallas}}\) and \(p_{\text{Vesta}}\) to be \(\equiv 1 \pmod{2^{32}}\), so that by Corollary~\ref{finitefields:cor:rootsofunity} each field has a multiplicative subgroup of order \(2^{32}\), furnishing the high-order roots of unity needed for fast (number-theoretic-transform-based) polynomial multiplication.
\end{remark}

\begin{remark}[The role of cyclicity and Frobenius in cryptography]
\label{finitefields:rem:cryptosummary}
Two structural facts from this section recur throughout the cryptographic development. First, Theorem~\ref{finitefields:thm:cyclic} (cyclicity of \(\Fq^\times\)) underlies the discrete-logarithm problem and guarantees the existence of the roots of unity used in polynomial commitment schemes. Second, the Frobenius endomorphism (Definition~\ref{finitefields:def:frobenius}) is not merely an abstraction: on elliptic curves over \(\Fp\) it induces an endomorphism whose characteristic polynomial \(T^2 - tT + p\) encodes the curve order \(N = p + 1 - t\) via the Hasse trace \(t\), tying the count \(|E(\Fp)|\) directly to the field-theoretic Frobenius studied here. We take up these connections in \cref{sec:ellipticcurves}.
\end{remark}

\clearpage

\section{Number theory for cryptography}\label{sec:numbertheory}

Modern public-key cryptography resides almost entirely inside finite cyclic groups and finite fields, and the security of the schemes used in Halo~2 and Zcash Orchard rests on number-theoretic facts about such structures. This section assembles the classical number theory of the integers modulo \(n\) that the later volumes use: the Chinese Remainder Theorem and the computation of Euler's totient, Fermat's little theorem and Euler's generalisation in their role as workhorses of modular exponentiation, primitive roots, the theory of quadratic residues with the Legendre symbol and the Euler criterion, algorithms for extracting square roots modulo a prime, and finally the discrete logarithm problem and the reason the \emph{prime order} of a group is so important for cryptographic hardness.

Throughout, \(n\) and \(m\) denote positive integers and \(p\), \(q\) denote primes. The section builds directly on the arithmetic of \cref{sec:intmod}, the group theory of \cref{sec:groups}, and the cyclicity theorem of \cref{sec:finitefields}; several statements proved there reappear here, restated in arithmetic language, with proofs reduced to citations.

\subsection{The integers modulo \texorpdfstring{\(n\)}{n} and units}\label{numbertheory:subsec:units}

Recall from \cref{sec:intmod} that \(a \in \Z/n\Z\) is a unit if and only if \(\gcd(a,n) = 1\) (Theorem~\ref{intmod:thm:unitiff}), and that the units form a finite abelian group \((\Z/n\Z)^\times\) under multiplication (Proposition~\ref{intmod:prop:unitgroup}). In particular, when \(p\) is prime every nonzero residue is coprime to \(p\), so \(\Fp^\times = \{1,2,\dots,p-1\}\) is a group of order \(p-1\) and \(\Fp\) is a field (Corollary~\ref{intmod:cor:fpfield}).

\subsection{Euler's totient function}\label{numbertheory:subsec:totient}

\begin{definition}[Euler totient]\label{numbertheory:def:totient}
\emph{Euler's totient function} \(\varphi(n)\) is the number of integers \(a\) with \(1 \le a \le n\) and \(\gcd(a,n)=1\). Equivalently, \(\varphi(n) = \lvert (\Z/n\Z)^\times \rvert\).
\end{definition}

\begin{example}\label{numbertheory:ex:totient-small}
We have \(\varphi(1)=1\), \(\varphi(p)=p-1\) for prime \(p\), and \(\varphi(12)=\lvert\{1,5,7,11\}\rvert = 4\).
\end{example}

Computing \(\varphi\) requires two facts: its value on prime powers, and its behaviour on coprime factors.

\begin{proposition}[Totient of a prime power]\label{numbertheory:prop:totient-prime-power}
For a prime \(p\) and integer \(k \ge 1\), \(\varphi(p^k) = p^k - p^{k-1} = p^k\left(1 - \tfrac1p\right)\).
\end{proposition}

\begin{proof}
Among \(1,\dots,p^k\) the integers \emph{not} coprime to \(p^k\) are exactly the multiples of \(p\), namely \(p, 2p, \dots, p^{k-1}\cdot p\), of which there are \(p^{k-1}\). Subtracting gives \(p^k - p^{k-1}\).
\end{proof}

\begin{theorem}[Chinese Remainder Theorem]\label{numbertheory:thm:crt}
If \(\gcd(m,n)=1\), then the map
\[
\Z/mn\Z \longrightarrow \Z/m\Z \times \Z/n\Z, \qquad a \bmod mn \longmapsto (a \bmod m,\; a \bmod n)
\]
is an isomorphism of rings. It restricts to a group isomorphism \((\Z/mn\Z)^\times \cong (\Z/m\Z)^\times \times (\Z/n\Z)^\times\).
\end{theorem}

\begin{proof}
The map is a well-defined ring homomorphism since reduction modulo \(m\) and modulo \(n\) are. It is injective: if \(a \equiv 0 \pmod m\) and \(a \equiv 0 \pmod n\) with \(\gcd(m,n)=1\), then \(mn \mid a\), so \(a \equiv 0 \pmod{mn}\); thus the kernel is trivial. Both sides have \(mn\) elements, so injectivity forces bijectivity. A ring isomorphism carries units to units, and \((u,v)\) is a unit of the product ring iff \(u,v\) are units in their factors; this gives the multiplicative statement.
\end{proof}

\begin{corollary}[Multiplicativity and the product formula]\label{numbertheory:cor:totient-mult}
If \(\gcd(m,n)=1\) then \(\varphi(mn) = \varphi(m)\varphi(n)\). Hence if \(n = \prod_i p_i^{k_i}\) is the prime factorisation of \(n\), then
\[
\varphi(n) = \prod_i \left(p_i^{k_i} - p_i^{k_i-1}\right) = n \prod_{p \mid n}\left(1 - \frac1p\right).
\]
\end{corollary}

\begin{proof}
The cardinality of a product of groups is the product of cardinalities, so Theorem~\ref{numbertheory:thm:crt} gives \(\varphi(mn)=\varphi(m)\varphi(n)\). Applying this across the distinct prime powers in the factorisation and using Proposition~\ref{numbertheory:prop:totient-prime-power} yields the product formula.
\end{proof}

Gauss's divisor-sum identity \(\sum_{d \mid n} \varphi(d) = n\), proved as Lemma~\ref{finitefields:lem:totient-sum}, completes the totient's basic theory.

\subsection{Fermat's little theorem and Euler's theorem}\label{numbertheory:subsec:fermat-euler}

The statements of this subsection were proved as Corollary~\ref{groups:cor-euler-fermat}, as instances of Lagrange's theorem; we restate them here in the arithmetic form in which cryptography uses them.

\begin{theorem}[Euler's theorem]\label{numbertheory:thm:euler}
If \(\gcd(a,n)=1\), then
\[
a^{\varphi(n)} \equiv 1 \pmod{n}.
\]
\end{theorem}

\begin{proof}
The hypothesis says \(a \bmod n \in (\Z/n\Z)^\times\), a finite group of order \(\varphi(n)\); now apply Corollary~\ref{groups:cor-aG-e} (every element of a finite group of order \(N\) satisfies \(g^N = e\)), a consequence of Lagrange's theorem.
\end{proof}

\begin{corollary}[Fermat's little theorem]\label{numbertheory:cor:fermat}
Let \(p\) be prime. If \(p \nmid a\), then \(a^{p-1} \equiv 1 \pmod p\). For \emph{every} integer \(a\), \(a^{p} \equiv a \pmod p\).
\end{corollary}

\begin{proof}
Take \(n=p\) in Theorem~\ref{numbertheory:thm:euler}, using \(\varphi(p)=p-1\), to get the first statement. Multiplying \(a^{p-1}\equiv 1\) by \(a\) yields \(a^p \equiv a\) when \(p \nmid a\); and when \(p \mid a\) both sides are \(\equiv 0\), so \(a^p \equiv a \pmod p\) in all cases.
\end{proof}

\begin{remark}\label{numbertheory:rem:fermat-rsa}
Euler's theorem underlies RSA: if \(ed \equiv 1 \pmod{\varphi(n)}\) then \((a^e)^d = a^{1+k\varphi(n)} \equiv a \pmod n\) whenever \(\gcd(a,n)=1\). For the elliptic-curve cryptography of Orchard one instead exploits Fermat's little theorem in the field \(\Fp\): it shows every nonzero element satisfies \(x^{p-1}=1\), so one can compute inversion as \(x^{-1} = x^{p-2}\), and it underlies the Euler criterion for square roots developed below.
\end{remark}

\subsection{Multiplicative order and primitive roots}\label{numbertheory:subsec:order}

\begin{definition}[Multiplicative order]\label{numbertheory:def:ord}
Let \(\gcd(a,n)=1\). The \emph{multiplicative order} of \(a\) modulo \(n\), written \(\ord_n(a)\), is the least positive integer \(k\) with \(a^k \equiv 1 \pmod n\). (Such a \(k\) exists because \(a^{\varphi(n)}\equiv 1\) by Theorem~\ref{numbertheory:thm:euler}.)
\end{definition}

This is the order of \(a\) as an element of the group \((\Z/n\Z)^\times\), so the order theory of \cref{sec:groups} applies verbatim: \(a^m \equiv 1 \pmod n\) if and only if \(\ord_n(a) \mid m\) (Proposition~\ref{groups:prop-order-divides}) --- in particular \(\ord_n(a) \mid \varphi(n)\) by Euler's theorem --- and \(\ord_n(a^j) = \ord_n(a)/\gcd(j, \ord_n(a))\) (Corollary~\ref{groups:cor-order-power}).

\begin{definition}[Primitive root]\label{numbertheory:def:primitive-root}
A \emph{primitive root} modulo \(n\) is an element \(g\) with \(\ord_n(g) = \varphi(n)\); equivalently, \(g\) generates \((\Z/n\Z)^\times\) as a cyclic group, so that \(\{g^0, g^1, \dots, g^{\varphi(n)-1}\}\) lists all the units.
\end{definition}

Primitive roots need not exist for every \(n\) (for instance \((\Z/8\Z)^\times \cong \Z/2\Z \times \Z/2\Z\) has no element of order \(4\)). The decisive case for cryptography is \(n=p\) prime, where the multiplicative group is always cyclic.

\begin{corollary}[\(\Fp^\times\) is cyclic]\label{numbertheory:thm:fp-cyclic}
For every prime \(p\), the group \(\Fp^\times = (\Z/p\Z)^\times\) is cyclic of order \(p-1\). In particular primitive roots modulo \(p\) exist, and there are \(\varphi(p-1)\) of them; Example~\ref{finitefields:ex:primitiveroot} exhibits them for \(p = 7\).
\end{corollary}

\begin{proof}
\(\Fp^\times\) is a finite subgroup (of order \(p-1\)) of the multiplicative group of the field \(\Fp\), hence cyclic by Theorem~\ref{finitefields:thm:cyclic}; the generator count is \(\varphi(p-1)\) by Corollary~\ref{groups:cor-order-power}.
\end{proof}

\begin{remark}\label{numbertheory:rem:cyclic-subgroup}
Cyclicity is why the prime-order groups used in elliptic-curve cryptography behave like \(\Z/q\Z\) under addition once one fixes a generator: a cyclic group of order \(q\) is isomorphic to \((\Z/q\Z, +)\) via \(g^k \leftrightarrow k\) (Theorem~\ref{groups:thm-cyclic-classification}). The full classification, which we do not need, states that primitive roots modulo \(n\) exist exactly when \(n \in \{1,2,4,p^k,2p^k\}\) for an odd prime \(p\).
\end{remark}

\subsection{Quadratic residues and the Euler criterion}\label{numbertheory:subsec:qr}

We now turn to squares modulo a prime, which control whether a curve point with a given \(x\)-coordinate exists, and which we need to compute square roots in \(\Fp\).

\begin{definition}[Quadratic residue]\label{numbertheory:def:qr}
Let \(p\) be an odd prime and \(a\) an integer with \(p \nmid a\). We say \(a\) is a \emph{quadratic residue} modulo \(p\) (a QR) if the congruence \(x^2 \equiv a \pmod p\) has a solution, and a \emph{quadratic non-residue} (a QNR) otherwise.
\end{definition}

\begin{proposition}[Exactly half are residues]\label{numbertheory:prop:half-qr}
For an odd prime \(p\), exactly \(\tfrac{p-1}{2}\) of the nonzero residues modulo \(p\) are quadratic residues and \(\tfrac{p-1}{2}\) are non-residues. The squaring map \(x \mapsto x^2\) is exactly two-to-one on \(\Fp^\times\).
\end{proposition}

\begin{proof}
Consider \(\sigma : \Fp^\times \to \Fp^\times\), \(\sigma(x)=x^2\). Its image is exactly the set of QRs. For \(b\) in the image, \(\sigma(x)=\sigma(y)\) iff \(x^2=y^2\) iff \((x-y)(x+y)=0\) iff \(y = \pm x\); since \(p\) is odd, \(x \ne -x\), so each fibre has exactly two elements. Hence \(\lvert \img \sigma\rvert = \tfrac{p-1}{2}\), and the complement has the same size.
\end{proof}

\begin{theorem}[Euler's criterion]\label{numbertheory:thm:euler-criterion}
Let \(p\) be an odd prime and \(p \nmid a\). Then
\[
a^{\frac{p-1}{2}} \equiv \begin{cases} \phantom{-}1 \pmod p & \text{if } a \text{ is a quadratic residue,}\\[2pt] -1 \pmod p & \text{if } a \text{ is a quadratic non-residue.}\end{cases}
\]
\end{theorem}

\begin{proof}
By Fermat, \(a^{p-1}\equiv 1\), so \(\bigl(a^{(p-1)/2}\bigr)^2 \equiv 1\), and since the only square roots of \(1\) modulo a prime are \(\pm 1\), we have \(a^{(p-1)/2}\equiv \pm 1 \pmod p\).

If \(a\) is a QR, say \(a \equiv x^2\), then \(a^{(p-1)/2}\equiv x^{p-1}\equiv 1\) by Fermat. Thus all \(\tfrac{p-1}{2}\) quadratic residues are roots of the polynomial \(t^{(p-1)/2}-1\) over \(\Fp\). By Theorem~\ref{polynomials:thm:root-bound} this polynomial has at most \(\tfrac{p-1}{2}\) roots, so the QRs are \emph{exactly} its roots. Any QNR \(a\) is therefore not a root, yet still satisfies \(a^{(p-1)/2}=\pm 1\); hence \(a^{(p-1)/2}\equiv -1\).
\end{proof}

The standard notation for quadratic-residue status is the Legendre symbol.

\begin{definition}[Legendre symbol]\label{numbertheory:def:legendre}
For an odd prime \(p\) and an integer \(a\), the \emph{Legendre symbol} is
\[
\legendre{a}{p} = \begin{cases} \phantom{-}0 & \text{if } p \mid a,\\ \phantom{-}1 & \text{if } a \text{ is a quadratic residue modulo } p,\\ -1 & \text{if } a \text{ is a quadratic non-residue modulo } p.\end{cases}
\]
Euler's criterion (Theorem~\ref{numbertheory:thm:euler-criterion}) then reads, uniformly,
\[
\legendre{a}{p} \equiv a^{\frac{p-1}{2}} \pmod p,
\]
so the symbol is computable by one modular exponentiation. The criterion also shows the symbol is \emph{completely multiplicative} in \(a\): \(\legendre{ab}{p} = \legendre{a}{p}\legendre{b}{p}\), since \((ab)^{(p-1)/2} = a^{(p-1)/2} b^{(p-1)/2}\).
\end{definition}

\begin{example}[\(13\) is a non-square in both Pasta fields]\label{numbertheory:ex:pasta-nonsquare}
The Pasta hash-to-curve construction used by Orchard fixes \(Z = -13\) as a verifiably non-square constant in both base fields. Since \(p \equiv q \equiv 1 \pmod 4\), the element \(-1\) is a square in both fields, so \(-13\) is a non-square precisely when \(13\) is. We verify the latter here. Let \(p, q\) be the Pasta primes of \cref{sec:ellipticcurves}. Both satisfy \(p \equiv q \equiv 1 \pmod 4\), and \(13 \equiv 1 \pmod 4\), so the law of quadratic reciprocity --- which for two odd primes, at least one \(\equiv 1 \pmod 4\), states that each is a square modulo the other or neither is --- gives \(\legendre{13}{p} = \legendre{p}{13}\) and \(\legendre{13}{q} = \legendre{q}{13}\). Reducing the primes modulo \(13\) and comparing against the squares modulo \(13\), namely \(\{1, 3, 4, 9, 10, 12\}\):
\[
p \equiv 5 \pmod{13}, \qquad q \equiv 8 \pmod{13}, \qquad 5, 8 \notin \{1,3,4,9,10,12\},
\]
whence
\[
\legendre{13}{p} = \legendre{5}{13} = -1, \qquad \legendre{13}{q} = \legendre{8}{13} = -1 .
\]
Thus \(13\) is a quadratic non-residue modulo both \(p\) and \(q\). (One can avoid reciprocity entirely by evaluating \(13^{(p-1)/2} \bmod p\) and \(13^{(q-1)/2} \bmod q\) directly via Euler's criterion; both equal \(-1\).)
\end{example}

\subsection{Square roots modulo a prime}\label{numbertheory:subsec:sqrt}

Deciding that \(a\) is a residue is one matter; \emph{producing} an \(x\) with \(x^2 \equiv a\) is another, and we need it, for instance, to recover a curve point's \(y\)-coordinate from its \(x\)-coordinate (point decompression). The difficulty depends on \(p \bmod 4\).

\begin{proposition}[The \(p \equiv 3 \pmod 4\) shortcut]\label{numbertheory:prop:sqrt-3mod4}
Let \(p \equiv 3 \pmod 4\) be prime and let \(a\) be a quadratic residue modulo \(p\) with \(p \nmid a\). Then
\[
x \equiv a^{\frac{p+1}{4}} \pmod p
\]
is a square root of \(a\); the two square roots are \(\pm x\).
\end{proposition}

\begin{proof}
Since \(p \equiv 3 \pmod 4\), \(\tfrac{p+1}{4}\) is a positive integer. Compute
\[
x^2 = a^{\frac{p+1}{2}} = a^{\frac{p-1}{2}}\cdot a \equiv \legendre{a}{p}\cdot a \equiv 1 \cdot a = a \pmod p,
\]
using Euler's criterion and the assumption that \(a\) is a residue, so \(\legendre{a}{p}=1\). The two roots are \(\pm x\) by Proposition~\ref{numbertheory:prop:half-qr}.
\end{proof}

\begin{remark}\label{numbertheory:rem:pasta-3mod4}
The single-exponentiation shortcut requires \(p \equiv 3 \pmod 4\), which is one reason cryptographers often favour such primes. For \(p \equiv 5 \pmod 8\) a closed-form variant (Atkin's formula) still exists. In general the cost of square-root extraction grows with the \(2\)-adic valuation \(s\) of \(p - 1\): writing \(p - 1 = 2^s d\) with \(d\) odd, the Tonelli--Shanks algorithm below performs up to \(s\) corrective iterations. The Pasta primes have \(s = 32\) (chosen deliberately for FFT-friendliness, \cref{rootsofunity:rem-twoadic}), so their implementations use Tonelli--Shanks.
\end{remark}

When \(p \equiv 1 \pmod 4\) one writes \(p - 1 = 2^s d\) with \(d\) odd and \(s \ge 2\). The Tonelli--Shanks algorithm operates inside the cyclic \(2\)-power-torsion subgroup of \(\Fp^\times\).

\begin{theorem}[Tonelli--Shanks]\label{numbertheory:thm:tonelli-shanks}
Let \(p\) be an odd prime, \(p-1 = 2^s d\) with \(d\) odd, and let \(a\) be a quadratic residue with \(p \nmid a\). One can compute a square root of \(a\) modulo \(p\) using \(O(s^2 + \log p)\) multiplications plus a few modular exponentiations, given any quadratic non-residue \(z\) modulo \(p\).
\end{theorem}

\begin{proof}[Sketch]
Fix a non-residue \(z\) (found by trying random elements; by Proposition~\ref{numbertheory:prop:half-qr} a random element is a non-residue with probability \(\tfrac12\), and one checks \(\legendre{z}{p}\) in \(O((\log p)^2)\) time). Set \(c = z^{d}\); this generates the unique cyclic subgroup \(H\) of order \(2^s\) in \(\Fp^\times\). Indeed, by Euler's criterion \(z^{(p-1)/2} \equiv -1\), so \((z^d)^{2^{s-1}} = z^{(p-1)/2} = -1 \neq 1\) while \((z^d)^{2^s} = z^{p-1} = 1\); hence \(\ord_p(z^d) = 2^s\). Initialise
\[
x = a^{\frac{d+1}{2}}, \qquad t = a^{d}, \qquad m = s.
\]
Then \(x^2 = a\cdot t\), so \(x\) is a square root of \(a\) \emph{up to} the factor \(t \in H\); moreover \(t^{2^{s-1}}=a^{(p-1)/2}=1\), so \(t\) lies in \(H\) and has \(2\)-power order. The algorithm repeatedly drives \(t\) to \(1\): at each step find the least \(i\) with \(t^{2^{i}} = 1\); then \(b = c^{\,2^{m-i-1}}\) satisfies \(b^2 \in H\) and replacing
\[
x \leftarrow x b, \qquad t \leftarrow t b^2, \qquad c \leftarrow b^2, \qquad m \leftarrow i
\]
strictly decreases the order of \(t\) while preserving the invariant \(x^2 = a\,t\). Since the order of \(t\) is a strictly decreasing power of \(2\), after at most \(s\) iterations \(t = 1\) and \(x^2 = a\). Correctness rests on the cyclic structure of \(H\) (Corollary~\ref{numbertheory:thm:fp-cyclic}); the running-time bound counts the at most \(s\) iterations, each costing \(O(s)\) squarings to locate \(i\).
\end{proof}

\begin{remark}\label{numbertheory:rem:sqrt-uniqueness}
In every case the square roots come in a pair \(\pm x\) (Proposition~\ref{numbertheory:prop:half-qr}), so to pin down a canonical root one fixes a convention, e.g.\ ``choose the root whose least nonnegative representative is even'', as elliptic-curve point encodings do. If \(a\) is \emph{not} a residue, all these algorithms detect failure (e.g.\ the candidate fails the final check \(x^2 \equiv a\)).
\end{remark}

\subsection{The discrete logarithm problem}\label{numbertheory:subsec:dlog}

The hardness underpinning elliptic-curve cryptography is the difficulty of \emph{inverting} exponentiation in a finite cyclic group. We phrase it here for a generic cyclic group, since in Orchard the group is the group of points of an elliptic curve, written additively, of large prime order.

\begin{definition}[Discrete logarithm problem]\label{numbertheory:def:dlp}
Let \(G = \langle g\rangle\) be a finite cyclic group of order \(N\), written multiplicatively, with generator \(g\). Given \(h \in G\), the \emph{discrete logarithm} of \(h\) to base \(g\), written \(\log_g h\), is the unique residue \(x \in \Z/N\Z\) with \(g^x = h\). The \emph{discrete logarithm problem} (DLP) is to compute \(x\) from \((G,g,h)\).
\end{definition}

The map \(x \mapsto g^x\) is a group isomorphism \((\Z/N\Z, +) \xrightarrow{\sim} (G,\cdot)\) (it is a homomorphism with trivial kernel because \(g\) has order \(N\)), so \(\log_g\) is genuinely a well-defined inverse. Computing \(g^x\) from \(x\) is fast (repeated squaring, \(O(\log N)\) operations); the DLP asks for the reverse, for which no efficient algorithm is known in well-chosen groups.

\begin{remark}[Generic complexity and group choice]\label{numbertheory:rem:dlp-generic}
Two ``generic'' algorithms that operate in any group of order \(N\), using only the group operation, solve the DLP in \(O(\sqrt N)\) time: \emph{baby-step giant-step} (deterministic, \(O(\sqrt N)\) time and space) and \emph{Pollard's rho} (randomised, \(O(\sqrt N)\) time and \(O(1)\) space). Shoup's theorem shows \(\Omega(\sqrt N)\) is optimal for generic algorithms when \(N\) is prime. Thus to reach a \(\lambda\)-bit security level one requires \(N \approx 2^{2\lambda}\); this is why the Pallas/Vesta scalar fields have \(\approx 255\)-bit prime order, targeting \(\approx 128\)-bit security. For the integers modulo \(p\) under multiplication, sub-exponential index-calculus attacks exist, which is why elliptic-curve groups (where no such attack is known) permit much smaller parameters.
\end{remark}

The next result makes the phrase ``well-chosen'' precise: if the group order \(N\) is \emph{composite} with small prime factors, the DLP reduces to several DLPs in small subgroups, and is therefore easy.

\begin{theorem}[Pohlig--Hellman reduction]\label{numbertheory:thm:pohlig-hellman}
Let \(G = \langle g\rangle\) be cyclic of order \(N = \prod_{i=1}^{r} q_i^{e_i}\). One can compute the discrete logarithm in \(G\) by solving the discrete logarithm in each subgroup of prime order \(q_i\) (a bounded number of times) and combining the answers with the Chinese Remainder Theorem. If the largest prime factor of \(N\) is \(q_{\max}\), the total cost is
\[
O\!\left(\sum_i e_i\bigl(\log N + \sqrt{q_i}\bigr)\right)
\]
group operations, dominated by \(O(\sqrt{q_{\max}})\).
\end{theorem}

\begin{proof}[Sketch]
The reduction proceeds in two stages. \emph{(a) CRT across coprime factors.} By the cyclic structure of \(G\) and the Chinese Remainder Theorem, \(G\) decomposes as the internal direct product of its subgroups \(G_i\) of prime-power order \(q_i^{e_i}\) (the unique subgroups of those orders, Theorem~\ref{groups:thm-cyclic-subgroups}). Knowing \(x = \log_g h\) modulo each \(q_i^{e_i}\) determines \(x\) modulo \(N = \prod q_i^{e_i}\) uniquely, by the Chinese Remainder Theorem (Theorem~\ref{numbertheory:thm:crt}, applied to the integers \(\Z/N\Z\)). To find \(x \bmod q_i^{e_i}\) one projects: set \(g_i = g^{N/q_i^{e_i}}\) and \(h_i = h^{N/q_i^{e_i}}\); then \(g_i\) has order \(q_i^{e_i}\) and \(\log_{g_i} h_i \equiv x \pmod{q_i^{e_i}}\).

\emph{(b) Lifting a prime power digit by digit.} Within a group of order \(q^{e}\), write the unknown exponent in base \(q\) as \(x \equiv x_0 + x_1 q + \dots + x_{e-1}q^{e-1} \pmod{q^e}\) with \(0 \le x_j < q\). Raising the relevant elements to the power \(q^{e-1}\) lands inside the order-\(q\) subgroup and isolates \(x_0\) as a single discrete log in a group of order \(q\); subtracting off its contribution and repeating isolates \(x_1\), and so on. Each digit costs one DLP in the order-\(q\) subgroup, solvable in \(O(\sqrt q)\) by baby-step giant-step or Pollard rho, plus \(O(\log N)\) for the exponentiations. Summing over the \(e_i\) digits of each of the \(r\) primes yields the stated bound.
\end{proof}

\begin{corollary}[Why prime order matters]\label{numbertheory:cor:prime-order}
If \(N\) has a small prime factor \(q\), then by Theorem~\ref{numbertheory:thm:pohlig-hellman} an attacker can recover the discrete logarithm modulo that factor in \(O(\sqrt q)\) time, leaking part of the secret \(x\). Security therefore requires \(N\) to have a \emph{large} prime factor; the safest and simplest choice is to make \(N\) itself a large prime, so that \(G\) has \emph{no} proper nontrivial subgroups and Pohlig--Hellman offers no reduction at all. This is precisely why the elliptic-curve groups used in Halo~2 and Orchard have large prime order.
\end{corollary}

\begin{remark}[Cofactors and subgroup attacks]\label{numbertheory:rem:cofactor}
In practice an elliptic curve over \(\Fp\) has order \(h \cdot q\) where \(q\) is a large prime and \(h\) is a small \emph{cofactor}; cryptographic operations stay confined to the prime-order-\(q\) subgroup. If protocols neglect to check that incoming points lie in this subgroup (or fail to clear the cofactor), an adversary can supply a low-order point and use Pohlig--Hellman in the small factor to extract information modulo \(h\) -- a real-world ``small-subgroup attack''. The Pasta curves of Halo~2 have cofactor \(h = 1\), i.e.\ \emph{prime} order, which eliminates this class of attacks at the group level. This is the cryptographic payoff of all the number theory developed above: the multiplicative and additive structure of \(\Z/q\Z\) for a large prime \(q\) is what gives the scheme both its functionality (a field of scalars) and its security (no Pohlig--Hellman reduction).
\end{remark}

\clearpage

\section{Linear algebra over a field}\label{sec:linalg}

Linear algebra is the study of \emph{vector spaces} — collections of objects that can be added together and scaled by elements of a field — together with the structure-preserving maps between them. It is the computational backbone of the cryptography this monograph develops: commitments to vectors, the linear constraints that encode arithmetic circuits, the polynomial-evaluation arguments at the heart of the inner-product argument, and the Pedersen-style multiscalar multiplications that Halo 2 and Orchard use are all, at bottom, statements about linear maps and (bi)linear forms over a finite field. This section develops the theory from the ground up, assuming only that the reader knows what a field is (a set \(\F\) with addition and multiplication in which every nonzero element is invertible; \cref{sec:ringsfields}). The development then specialises to the two settings that matter here: the coordinate space \(\F^n\) over a finite field \(\F = \Fq\), and the ``mixed'' pairing between scalar vectors and vectors of group elements.

Throughout, \(\F\) denotes an arbitrary field. We write its additive identity \(0\), its multiplicative identity \(1\), and call elements of \(\F\) \emph{scalars}. We use the field axioms freely: addition and multiplication are commutative and associative, multiplication distributes over addition, and every \(a \neq 0\) has an inverse \(a^{-1}\).

\subsection{Vector spaces and subspaces}\label{linalg:subsec-vs}

\begin{definition}[Vector space]\label{linalg:def-vs}
Let \(\F\) be a field. A \emph{vector space over \(\F\)} (or an \emph{\(\F\)-vector space}) is a set \(V\) together with two operations,
\[
+ \colon V \times V \to V \qquad\text{(vector addition)}, \qquad
\cdot \colon \F \times V \to V \qquad\text{(scalar multiplication)},
\]
satisfying the following axioms for all \(u, v, w \in V\) and all \(a, b \in \F\). We call the elements of \(V\) \emph{vectors}.
\begin{enumerate}
\item \emph{(Additive abelian group.)} \(V\) is an abelian group under \(+\): addition is associative and commutative, there is a zero vector \(\mathbf{0} \in V\) with \(v + \mathbf{0} = v\), and every \(v\) has an additive inverse \(-v\) with \(v + (-v) = \mathbf{0}\).
\item \emph{(Compatibility.)} \(a \cdot (b \cdot v) = (ab) \cdot v\).
\item \emph{(Identity.)} \(1 \cdot v = v\).
\item \emph{(Distributivity over vector addition.)} \(a \cdot (u + v) = a \cdot u + a \cdot v\).
\item \emph{(Distributivity over scalar addition.)} \((a + b) \cdot v = a \cdot v + b \cdot v\).
\end{enumerate}
We usually write the scalar product \(a \cdot v\) as \(av\), and we drop the distinction between the scalar \(0 \in \F\) and the vector \(\mathbf{0} \in V\) when context makes it clear.
\end{definition}

A few immediate consequences of the axioms recur so frequently that we record them once and for all.

\begin{proposition}\label{linalg:prop-vs-basic}
Let \(V\) be an \(\F\)-vector space, \(v \in V\), and \(a \in \F\). Then:
\begin{enumerate}
\item \(0 \cdot v = \mathbf{0}\);
\item \(a \cdot \mathbf{0} = \mathbf{0}\);
\item \((-1) \cdot v = -v\);
\item if \(a v = \mathbf{0}\) then \(a = 0\) or \(v = \mathbf{0}\).
\end{enumerate}
\end{proposition}

\begin{proof}
(1) By distributivity over scalar addition, \(0 \cdot v = (0 + 0) \cdot v = 0 \cdot v + 0 \cdot v\). Adding \(-(0 \cdot v)\) to both sides gives \(\mathbf{0} = 0 \cdot v\).

(2) Similarly \(a \cdot \mathbf{0} = a \cdot (\mathbf{0} + \mathbf{0}) = a \cdot \mathbf{0} + a \cdot \mathbf{0}\), and cancellation gives \(a \cdot \mathbf{0} = \mathbf{0}\).

(3) Using (1) and distributivity, \(v + (-1)\cdot v = 1\cdot v + (-1)\cdot v = (1 + (-1))\cdot v = 0 \cdot v = \mathbf{0}\), so \((-1)\cdot v\) is the additive inverse of \(v\), i.e.\ \((-1)\cdot v = -v\).

(4) Suppose \(av = \mathbf{0}\) and \(a \neq 0\). Since \(\F\) is a field, \(a^{-1}\) exists, and by the compatibility and identity axioms \(v = 1\cdot v = (a^{-1}a)\cdot v = a^{-1}\cdot(a v) = a^{-1}\cdot \mathbf{0} = \mathbf{0}\) using (2).
\end{proof}

\begin{example}[The coordinate space \(\F^n\)]\label{linalg:exa-Fn}
The single most important example is the \emph{coordinate space}
\[
\F^n = \{\, (x_1, \dots, x_n) : x_i \in \F \,\},
\]
the set of \(n\)-tuples of scalars, with componentwise operations
\[
(x_1, \dots, x_n) + (y_1, \dots, y_n) = (x_1 + y_1, \dots, x_n + y_n), \qquad
a \cdot (x_1, \dots, x_n) = (a x_1, \dots, a x_n).
\]
The zero vector is \((0, \dots, 0)\). All the axioms follow from the field axioms applied coordinatewise. When \(n = 0\) the result is the \emph{trivial} (or \emph{zero}) vector space \(\{\mathbf{0}\}\), consisting of the empty tuple alone. We typically write vectors of \(\F^n\) as columns and in boldface, \(\mathbf{x} = (x_1, \dots, x_n)^{\mathsf T}\).
\end{example}

\begin{example}[Further examples]\label{linalg:exa-more}
\begin{enumerate}
\item \emph{Polynomials.} The set \(\F[X]\) of all polynomials in one variable with coefficients in \(\F\) is an \(\F\)-vector space under ordinary addition and scalar multiplication. So is the subset \(\F[X]_{<n}\) of polynomials of degree less than \(n\); this is exactly the setting in which polynomial commitment schemes operate.
\item \emph{Functions.} For any set \(S\), the set \(\F^S\) of all functions \(S \to \F\) is an \(\F\)-vector space under pointwise operations. Example~\ref{linalg:exa-Fn} is the special case \(S = \{1, \dots, n\}\).
\item \emph{Matrices.} The set \(\F^{m \times n}\) of \(m \times n\) matrices with entries in \(\F\) is an \(\F\)-vector space under entrywise addition and scaling.
\item \emph{Field extensions.} If \(\F \subseteq \mathbb{K}\) is a subfield of a larger field \(\mathbb{K}\), then \(\mathbb{K}\) is an \(\F\)-vector space. For instance \(\C\) is a \(2\)-dimensional \(\R\)-vector space, and the finite field \(\Fq\) with \(q = p^k\) is a \(k\)-dimensional \(\Fp\)-vector space (Theorem~\ref{finitefields:thm:primepower} rested on exactly this).
\end{enumerate}
\end{example}

A vector space often contains smaller vector spaces sitting inside it. These are the subspaces.

\begin{definition}[Subspace]\label{linalg:def-subspace}
A subset \(W \subseteq V\) of an \(\F\)-vector space \(V\) is a \emph{(linear) subspace} if it is itself a vector space under the operations inherited from \(V\). Equivalently, \(W\) is a subspace iff it is nonempty and closed under the vector-space operations.
\end{definition}

\begin{proposition}[Subspace criterion]\label{linalg:prop-subspace}
A subset \(W \subseteq V\) is a subspace if and only if all three of the following hold:
\begin{enumerate}
\item \(\mathbf{0} \in W\);
\item \(W\) is closed under addition: \(u, v \in W \implies u + v \in W\);
\item \(W\) is closed under scalar multiplication: \(a \in \F,\ v \in W \implies a v \in W\).
\end{enumerate}
Equivalently, \(W \neq \varnothing\) and for all \(a, b \in \F\) and \(u, v \in W\) we have \(a u + b v \in W\).
\end{proposition}

\begin{proof}
If \(W\) is a subspace, the three conditions hold because they are part of (or follow from) the vector-space axioms restricted to \(W\). Conversely, suppose (1)–(3) hold. Closure under \(+\) and \(\cdot\) makes the operations well defined on \(W\). The additive inverse of \(v \in W\) is \((-1)v \in W\) by (3) and Proposition~\ref{linalg:prop-vs-basic}(3). The space \(V\) supplies the remaining axioms (associativity, commutativity, distributivity, etc.) to \(W\), because they are universally quantified identities that hold for all elements of \(V\), in particular for elements of \(W\). The ``equivalently'' clause follows by taking \((a,b) = (1,1)\) and \((a,b)=(a,0)\) to recover (2) and (3), and noting \(\mathbf 0 = 0\cdot v\) for any \(v\) in the nonempty \(W\).
\end{proof}

\begin{example}\label{linalg:exa-subspace}
In \(\F^3\), the set \(\{(x_1, x_2, x_3) : x_1 + x_2 + x_3 = 0\}\) is a subspace (the solution set of a homogeneous linear equation), whereas \(\{(x_1,x_2,x_3) : x_1 + x_2 + x_3 = 1\}\) is not, since it does not contain \(\mathbf 0\). More generally, the solution set of any homogeneous linear system \(A\mathbf x = \mathbf 0\) is a subspace of \(\F^n\); we revisit this in the study of kernels.
\end{example}

\begin{proposition}[Intersections and sums]\label{linalg:prop-int-sum}
Let \(\{W_i\}_{i \in I}\) be a family of subspaces of \(V\).
\begin{enumerate}
\item The intersection \(\bigcap_{i \in I} W_i\) is a subspace.
\item The \emph{sum} \(W_1 + \dots + W_k = \{ w_1 + \dots + w_k : w_j \in W_j \}\) of finitely many subspaces is a subspace, and it is the smallest subspace containing all the \(W_j\).
\end{enumerate}
In general the \emph{union} of two subspaces is \emph{not} a subspace.
\end{proposition}

\begin{proof}
(1) Each \(W_i\) contains \(\mathbf 0\), so the intersection does too. If \(u, v\) lie in every \(W_i\) and \(a, b \in \F\), then \(au + bv \in W_i\) for every \(i\) by Proposition~\ref{linalg:prop-subspace}, hence in the intersection.

(2) Closure: given \(\sum_j w_j\) and \(\sum_j w_j'\) in the sum and scalars \(a,b\), we have \(a\sum_j w_j + b\sum_j w_j' = \sum_j (a w_j + b w_j')\) with \(a w_j + b w_j' \in W_j\). It contains \(\mathbf 0 = \sum_j \mathbf 0\) and each \(W_j\) (take all other summands zero). Any subspace containing every \(W_j\) must, by closure under addition, contain all such sums; hence the sum is the smallest such subspace. An example shows the union claim: in \(\F^2\) with \(|\F| > 1\), the two axes \(\{(x,0)\}\) and \(\{(0,y)\}\) are subspaces whose union omits \((1,1) = (1,0)+(0,1)\).
\end{proof}
\subsection{Linear combinations, span, and linear independence}\label{linalg:subsec-span}

\begin{definition}[Linear combination and span]\label{linalg:def-span}
Let \(S \subseteq V\). A \emph{linear combination} of vectors in \(S\) is any finite sum
\[
a_1 v_1 + a_2 v_2 + \dots + a_k v_k, \qquad a_i \in \F,\ v_i \in S.
\]
The \emph{span} of \(S\), written \(\operatorname{span}(S)\) or \(\langle S\rangle\), is the set of all linear combinations of finitely many vectors in \(S\). By convention, \(\operatorname{span}(\varnothing) = \{\mathbf 0\}\). \(S\) \emph{spans} (or \emph{generates}) \(V\) if \(\operatorname{span}(S) = V\), and \(V\) is \emph{finitely generated} (or \emph{finite-dimensional}) if some finite set spans it.
\end{definition}

\begin{proposition}\label{linalg:prop-span-subspace}
For any \(S \subseteq V\), \(\operatorname{span}(S)\) is the smallest subspace of \(V\) containing \(S\); that is, it is a subspace, it contains \(S\), and it is contained in every subspace that contains \(S\).
\end{proposition}

\begin{proof}
A sum of two linear combinations of vectors in \(S\) is again such a linear combination, as is a scalar multiple of one, and \(\mathbf 0\) is the empty combination; hence \(\operatorname{span}(S)\) is a subspace by Proposition~\ref{linalg:prop-subspace}, and \(S \subseteq \operatorname{span}(S)\). If \(W\) is any subspace containing \(S\), then \(W\) contains every linear combination of elements of \(S\) (by closure under addition and scaling), so \(\operatorname{span}(S) \subseteq W\).
\end{proof}

\begin{definition}[Linear independence]\label{linalg:def-indep}
A finite list of vectors \(v_1, \dots, v_k \in V\) is \emph{linearly independent} if the only expression of the zero vector as a linear combination of them is the trivial one:
\[
a_1 v_1 + \dots + a_k v_k = \mathbf 0 \quad\Longrightarrow\quad a_1 = \dots = a_k = 0.
\]
Otherwise the list is \emph{linearly dependent}; we call a nontrivial relation \(\sum_i a_i v_i = \mathbf 0\) with some \(a_i \neq 0\) a \emph{dependence relation}. An arbitrary (possibly infinite) set \(S\) is linearly independent if every finite sublist of distinct vectors from \(S\) is independent.
\end{definition}

\begin{remark}\label{linalg:rem-indep}
Linear independence captures non-redundancy. A single vector \(v\) is independent iff \(v \neq \mathbf 0\). Two vectors are dependent iff one is a scalar multiple of the other. In general, a list is dependent iff some vector in it lies in the span of the others, as the next lemma makes precise.
\end{remark}

\begin{lemma}[Dependence and redundancy]\label{linalg:lem-dependence}
A list \(v_1, \dots, v_k\) with \(k \geq 1\) is linearly dependent if and only if some \(v_j\) is a linear combination of the others. Moreover, if \(v_1, \dots, v_k\) are independent but \(v_1, \dots, v_k, w\) are dependent, then \(w \in \operatorname{span}(v_1, \dots, v_k)\), and the coefficients expressing \(w\) are unique.
\end{lemma}

\begin{proof}
If \(v_j = \sum_{i \neq j} a_i v_i\), then \(\sum_{i\neq j} a_i v_i + (-1) v_j = \mathbf 0\) is a nontrivial relation. Conversely, given a dependence relation \(\sum_i a_i v_i = \mathbf 0\) with \(a_j \neq 0\), one solves \(v_j = -a_j^{-1}\sum_{i \neq j} a_i v_i\).

For the second statement, a dependence relation \(\sum_i a_i v_i + b w = \mathbf 0\) must have \(b \neq 0\) (otherwise it would be a nontrivial relation among the independent \(v_i\)), so \(w = -b^{-1}\sum_i a_i v_i \in \operatorname{span}(v_1,\dots,v_k)\). Uniqueness: if \(w = \sum_i c_i v_i = \sum_i c_i' v_i\), then \(\sum_i (c_i - c_i') v_i = \mathbf 0\), forcing \(c_i = c_i'\) by independence.
\end{proof}

\subsection{Bases, dimension, and linear maps}\label{linalg:subsec-basis}

\begin{definition}[Basis]\label{linalg:def-basis}
A \emph{basis} of \(V\) is a linearly independent list \(b_1, \dots, b_n\) that spans \(V\). By Lemma~\ref{linalg:lem-dependence}, every \(v \in V\) is then a \emph{unique} linear combination \(v = \sum_i c_i b_i\); the scalars \(c_i\) are the \emph{coordinates} of \(v\) in the basis, written \([v]_B = (c_1, \dots, c_n)\).
\end{definition}

The standard basis \(e_1, \dots, e_n\) of \(\F^n\) (where \(e_i\) has a \(1\) in position \(i\) and \(0\) elsewhere) and the monomial basis \(1, X, \dots, X^{n-1}\) of \(\F[X]_{<n}\) are the two examples in constant use. That all bases of a given space have the same size rests on the following exchange principle.

\begin{lemma}[Steinitz exchange]\label{linalg:lem-exchange}
If \(v_1, \dots, v_m\) are linearly independent in \(V\) and \(w_1, \dots, w_n\) span \(V\), then \(m \le n\).
\end{lemma}

\begin{proof}
We show by induction on \(k = 0, 1, \dots, m\) that, after renumbering the \(w_j\), the list \(v_1, \dots, v_k, w_{k+1}, \dots, w_n\) spans \(V\); taking \(k = m\) forces \(m \le n\), since the process consumes one \(w\) for each \(v\). The case \(k=0\) is the hypothesis. For the step, the spanning property at stage \(k\) writes \(v_{k+1} = \sum_{i \le k} a_i v_i + \sum_{j > k} c_j w_j\). Some \(c_j \neq 0\) --- otherwise \(v_{k+1}\) would be a combination of \(v_1, \dots, v_k\), contradicting independence --- say \(c_{k+1} \neq 0\) after renumbering. Solving for \(w_{k+1}\) expresses it through \(v_1, \dots, v_{k+1}, w_{k+2}, \dots, w_n\), so replacing \(w_{k+1}\) by \(v_{k+1}\) preserves the span.
\end{proof}

\begin{theorem}[Dimension]\label{linalg:thm-dimension}
Let \(V\) be a finitely generated \(\F\)-vector space. Then \(V\) has a basis, and any two bases have the same number of elements. We call this number the \emph{dimension} \(\dim_\F V\). In particular \(\dim_\F \F^n = n\).
\end{theorem}

\begin{proof}
A minimal spanning list (obtained from a finite spanning set by discarding redundant vectors, Lemma~\ref{linalg:lem-dependence}) is independent, hence a basis. If \(b_1, \dots, b_m\) and \(b_1', \dots, b_n'\) are two bases, Lemma~\ref{linalg:lem-exchange} applied in both directions gives \(m \le n\) and \(n \le m\). For \(\F^n\), the standard basis has \(n\) elements.
\end{proof}

\begin{definition}[Linear map, kernel, image]\label{linalg:def-linear-map}
A function \(T \colon V \to W\) between \(\F\)-vector spaces is \emph{linear} if \(T(au + bv) = a\,T(u) + b\,T(v)\) for all \(a, b \in \F\) and \(u, v \in V\). Its \emph{kernel} and \emph{image} are
\[
\ker T = \{v \in V : T(v) = \mathbf 0\} \subseteq V, \qquad \img T = \{T(v) : v \in V\} \subseteq W ;
\]
both are subspaces (immediate from the subspace criterion and linearity). \(T\) is injective iff \(\ker T = \{\mathbf 0\}\): if \(T(u) = T(v)\) then \(u - v \in \ker T\), and conversely. A bijective linear map is an \emph{isomorphism}.
\end{definition}

\begin{theorem}[Rank--nullity]\label{linalg:thm-rank-nullity}
Let \(T \colon V \to W\) be linear with \(V\) finite-dimensional. Then
\[
\dim V = \dim \ker T + \dim \img T .
\]
\end{theorem}

\begin{proof}
Choose a basis \(u_1, \dots, u_k\) of \(\ker T\) and extend it to a basis \(u_1, \dots, u_k, v_1, \dots, v_r\) of \(V\) (adjoin vectors outside the current span until the list spans; each adjunction preserves independence by Lemma~\ref{linalg:lem-dependence}, and the process stops by Lemma~\ref{linalg:lem-exchange}). We claim \(T(v_1), \dots, T(v_r)\) is a basis of \(\img T\), which yields \(\dim V = k + r = \dim\ker T + \dim\img T\). \emph{Spanning}: any \(T(v)\) with \(v = \sum_i a_i u_i + \sum_j c_j v_j\) equals \(\sum_j c_j T(v_j)\), since \(T(u_i) = \mathbf 0\). \emph{Independence}: if \(\sum_j c_j T(v_j) = \mathbf 0\), then \(\sum_j c_j v_j \in \ker T\), so \(\sum_j c_j v_j = \sum_i a_i u_i\) for some scalars \(a_i\); independence of the combined basis of \(V\) forces all \(c_j = 0\).
\end{proof}

Rank--nullity is the dimension-counting engine of the section: it shows, for instance, that a linear map from a higher-dimensional space into a lower-dimensional one must have a nontrivial kernel --- the fact that makes vector commitments compressing and only \emph{computationally} binding (Remark~\ref{linalg:rem-msm-meaning}).

\subsection{Bilinear forms and inner products on \(\F^n\)}\label{linalg:subsec-bilinear}

The maps considered so far are linear in a single argument. Commitment and proof systems are pervaded by expressions that are linear in \emph{each} of two arguments — most importantly the inner product. We treat these next. Because the fields of interest are finite, there is no notion of ``positive length'' or ``angle''; the appropriate level of generality is the \emph{bilinear form}.

\begin{definition}[Bilinear form]\label{linalg:def-bilinear}
A \emph{bilinear form} on an \(\F\)-vector space \(V\) is a function \(\beta \colon V \times V \to \F\) that is linear in each argument separately:
\[
\beta(au + a'u', v) = a\,\beta(u,v) + a'\,\beta(u',v), \qquad
\beta(u, bv + b'v') = b\,\beta(u,v) + b'\,\beta(u,v').
\]
It is \emph{symmetric} if \(\beta(u,v) = \beta(v,u)\) for all \(u,v\), and \emph{nondegenerate} if \(\beta(u,v) = 0\) for all \(v\) implies \(u = \mathbf 0\).
\end{definition}

Relative to a fixed basis \(b_1, \dots, b_n\) of \(V\), every bilinear form \(\beta\) is represented by the unique matrix \(M \in \F^{n \times n}\) with \(M_{ij} = \beta(b_i, b_j)\), via \(\beta(u,v) = [u]_B^{\mathsf T} M\, [v]_B\); we do not use this representation below and record only the special case that matters.

\begin{definition}[Standard inner product on \(\F^n\)]\label{linalg:def-stdip}
The \emph{standard inner product} (or \emph{dot product}) on \(\F^n\) is the symmetric bilinear form
\[
\ip{\mathbf a}{\mathbf b} = \sum_{i=1}^n a_i b_i = \mathbf a^{\mathsf T}\mathbf b,
\qquad \mathbf a = (a_1,\dots,a_n),\ \mathbf b = (b_1,\dots,b_n).
\]
Its matrix in the standard basis is the identity \(I_n\); it is symmetric and nondegenerate.
\end{definition}

\begin{remark}[No positivity over finite fields]\label{linalg:rem-no-positivity}
Over \(\R\) the standard inner product is \emph{positive definite}: \(\ip{\mathbf a}{\mathbf a} = \sum a_i^2 \ge 0\), vanishing only at \(\mathbf 0\), which yields the Euclidean norm \(\|\mathbf a\| = \sqrt{\ip{\mathbf a}{\mathbf a}}\) and the Cauchy--Schwarz inequality. Over a finite field \(\Fq\) there is no ordering, so ``positivity'' is meaningless: one may have \(\mathbf a \neq \mathbf 0\) with \(\ip{\mathbf a}{\mathbf a} = 0\) (an \emph{isotropic} vector). For example over \(\F_5\), \(\mathbf a = (1,2)\) gives \(1 + 4 = 5 \equiv 0\). The algebraically robust properties that survive are bilinearity, symmetry, and nondegeneracy; these are all that the arguments in proof systems rely on.
\end{remark}

\begin{remark}[Functionals via the inner product]\label{linalg:rem-riesz}
Nondegeneracy yields a clean dictionary between vectors and linear functionals. Each \(\mathbf a \in \F^n\) gives a functional \(\mathbf b \mapsto \ip{\mathbf a}{\mathbf b}\), and the map \(\mathbf a \mapsto \ip{\mathbf a}{\cdot}\) is an isomorphism from \(\F^n\) onto its \emph{dual space} \(\Hom(\F^n, \F)\) (both have dimension \(n\), and the map is injective by nondegeneracy). Thus ``a linear functional on \(\F^n\)'' and ``a vector in \(\F^n\)'' are interchangeable. Evaluating a committed polynomial \(p\) of coefficient vector \(\mathbf c\) at a point \(z\) is the inner product \(\ip{\mathbf c}{(1, z, z^2, \dots, z^{n-1})}\); this is precisely the shape an inner-product argument proves.
\end{remark}

\subsection{The mixed inner product with group elements}\label{linalg:subsec-mixed}

The commitments used in Halo 2 and Orchard pair a vector of \emph{scalars} with a vector of \emph{group elements}. We isolate the algebraic structure here; the section on elliptic curves introduces the relevant group fully, but only its abelian-group structure matters at present.

Let \((\G, +)\) be a finite abelian group, written additively, of prime order \(r\) with identity \(\mathcal O\). For an integer \(a\) and \(G \in \G\) write \([a]G\) for the \(a\)-fold sum \(G + \dots + G\) (and \([0]G = \mathcal O\), \([-a]G = -([a]G)\); this is the \emph{scalar multiplication} of the group). Because \(\G\) has order \(r\) with \(r\) prime, \([a]G\) depends only on \(a \bmod r\), so the operation descends to scalars from the field \(\Fr = \Z/r\Z\): for \(a \in \Fr\) and \(G \in \G\), \([a]G\) is well defined. With this action \(\G\) becomes a one-dimensional vector space over \(\Fr\) whenever \(\G\) is nontrivial — a viewpoint we exploit repeatedly.

\begin{proposition}[\(\G\) as an \(\Fr\)-vector space]\label{linalg:prop-G-vs}
Let \(\G\) be a group of prime order \(r\). Then the scalar multiplication \((a, G) \mapsto [a]G\) makes \(\G\) into a vector space over \(\Fr\) of dimension \(1\); any nonidentity element \(G\) is a basis, and every element is uniquely \([a]G\) for some \(a \in \Fr\).
\end{proposition}

\begin{proof}
The vector-space axioms are the standard laws of group scalar multiplication: \([a]([b]G) = [ab]G\), \([1]G = G\), \([a](G+H) = [a]G + [a]H\), and \([a+b]G = [a]G + [b]G\), all of which hold in any abelian group and respect reduction mod \(r\). For a nonidentity \(G\), Lagrange's theorem in a group of prime order \(r\) forces \(G\) to have order \(r\), so the \(r\) multiples \([0]G, [1]G, \dots, [r-1]G\) are distinct and exhaust \(\G\) (which has \(r\) elements). Hence \(\{G\}\) spans and is independent (\([a]G = \mathcal O\) iff \(a \equiv 0\)), a basis; uniqueness of \(a\) is unique representation in this basis.
\end{proof}

\begin{definition}[Mixed inner product / multiscalar multiplication]\label{linalg:def-msm}
Let \(\mathbf G = (G_1, \dots, G_n) \in \G^n\) be a fixed tuple of group elements and \(\mathbf a = (a_1, \dots, a_n) \in \Fr^n\) a vector of scalars. Their \emph{mixed inner product} (also called a \emph{multiscalar multiplication}, MSM, or \emph{Pedersen-style inner product}) is the group element
\[
\ip{\mathbf a}{\mathbf G} \;=\; \sum_{i=1}^n [a_i]\,G_i \;\in\; \G .
\]
We use the same angle-bracket notation as for the scalar inner product; whether the second argument is a scalar vector or a vector of group elements determines which one we mean.
\end{definition}

\begin{proposition}[Bilinearity of the mixed product]\label{linalg:prop-msm-bilinear}
Fix \(\mathbf G \in \G^n\). The map \(\mathbf a \mapsto \ip{\mathbf a}{\mathbf G}\) is \(\Fr\)-linear from \(\Fr^n\) to \(\G\):
\[
\ip{c\,\mathbf a + c'\,\mathbf a'}{\mathbf G} = [c]\ip{\mathbf a}{\mathbf G} + [c']\ip{\mathbf a'}{\mathbf G}.
\]
Symmetrically, for fixed \(\mathbf a\) the map \(\mathbf G \mapsto \ip{\mathbf a}{\mathbf G}\) is a group homomorphism \(\G^n \to \G\) that is \(\Fr\)-linear when one views \(\G^n\) as an \(\Fr\)-vector space. Thus \(\ip{\cdot}{\cdot} \colon \Fr^n \times \G^n \to \G\) is bilinear over \(\Fr\).
\end{proposition}

\begin{proof}
Using the vector-space laws on \(\G\) (Proposition~\ref{linalg:prop-G-vs}) componentwise,
\[
\ip{c\mathbf a + c'\mathbf a'}{\mathbf G} = \sum_i [\,c a_i + c' a_i'\,]G_i = \sum_i\big([c][a_i]G_i + [c'][a_i']G_i\big) = [c]\sum_i [a_i]G_i + [c']\sum_i [a_i']G_i,
\]
which is the first claim. Linearity in \(\mathbf G\) is identical by symmetry of the defining sum and the law \([a](G + H) = [a]G + [a]H\).
\end{proof}

\begin{remark}[Why this is the right abstraction]\label{linalg:rem-msm-meaning}
The mixed inner product is the computational heart of Pedersen commitments and of the Halo 2 inner-product argument. Reading \(\ip{\mathbf a}{\mathbf G}\) as a \emph{commitment} to the scalar vector \(\mathbf a\) under the public ``basis'' \(\mathbf G\), bilinearity is exactly what endows the commitment with its \emph{homomorphic} property: a commitment to \(c\mathbf a + c'\mathbf a'\) is the same linear combination \([c]\)- and \([c']\)-scaled of the individual commitments. Provided one chooses the \(G_i\) so that the prover knows no nontrivial relation \(\sum_i [a_i]G_i = \mathcal O\) (the discrete-logarithm hardness assumption, treated later), the commitment is binding: the map \(\mathbf a \mapsto \ip{\mathbf a}{\mathbf G}\) is computationally injective. This map \emph{cannot} be injective as a set map once \(n \ge 2\), since by rank--nullity its kernel (a subspace of \(\Fr^n\) mapping into the one-dimensional \(\G\)) has dimension at least \(n - 1\); binding is therefore a \emph{computational}, not information-theoretic, statement. A typical Pedersen commitment to a single value \(a\) with randomness \(\rho\) is the special case \([a]G + [\rho]H = \ip{(a,\rho)}{(G,H)}\), which is moreover \emph{perfectly hiding} because the \([\rho]H\) term is uniform in \(\G\).
\end{remark}

\subsection{How this machinery appears in commitment and proof systems}\label{linalg:subsec-crypto}

We now draw the threads together; each point below is developed rigorously in the later volumes of this series, but the linear-algebraic content is already in hand.

\begin{remark}[Witnesses, constraints, and assignments]\label{linalg:rem-constraints}
An arithmetic circuit over \(\Fq\) is, after \emph{arithmetisation}, a system of polynomial constraints in the wire values. Matrices over \(\Fq\) acting on the assignment vector \(\mathbf z \in \Fq^N\) capture the \emph{linear} part of such a system — the wiring/copy constraints of PLONK. Satisfying the linear constraints means \(\mathbf z\) lies in an affine subspace; the remaining (multiplicative) constraints are handled separately. Rank--nullity (Theorem~\ref{linalg:thm-rank-nullity}) quantifies the residual degrees of freedom, i.e.\ the dimension of the witness space.
\end{remark}

\begin{remark}[Commitments as linear maps]\label{linalg:rem-commit-linear}
A Pedersen vector commitment \(\mathbf a \mapsto \ip{\mathbf a}{\mathbf G} = \sum_i [a_i]G_i\) is a linear map \(\Fr^n \to \G\) (Proposition~\ref{linalg:prop-msm-bilinear}). Its linearity is the source of every algebraic manipulation a verifier performs: it can take known linear combinations of commitments, ``fold'' two commitment keys \(\mathbf G, \mathbf G'\) into \(\mathbf G + [x]\mathbf G'\) under a challenge \(x\), and check claimed openings by re-evaluating the same linear map. A polynomial commitment is the same map applied to a coefficient vector, and an \emph{evaluation} of that polynomial at \(z\) is the scalar inner product \(\ip{\mathbf c}{(1,z,\dots,z^{n-1})}\) (Remark~\ref{linalg:rem-riesz}).
\end{remark}

\begin{remark}[The inner-product relation]\label{linalg:rem-ipa}
The core statement an inner-product argument proves has the shape: ``I know vectors \(\mathbf a, \mathbf b \in \Fr^n\) such that \(P = \ip{\mathbf a}{\mathbf G} + \ip{\mathbf b}{\mathbf H} + [\,\ip{\mathbf a}{\mathbf b}\,]U\) for public \(\mathbf G, \mathbf H \in \G^n\) and \(U \in \G\).'' Here three of the constructions meet at once: two mixed inner products \(\ip{\mathbf a}{\mathbf G}, \ip{\mathbf b}{\mathbf H}\) (commitments to the witness vectors), one scalar inner product \(\ip{\mathbf a}{\mathbf b}\) (the value being argued), and the bilinearity of all three (which permits the prover and verifier to fold length-\(n\) instances into length-\(n/2\) ones, using the bilinearity of these inner products). The argument's soundness ultimately reduces to the nondegeneracy of these forms together with the hardness of finding kernel elements of the commitment map — that is, to exactly the linear-algebraic and computational facts assembled in this section.
\end{remark}

\begin{remark}[Summary]\label{linalg:rem-summary}
The objects constructed above — vector spaces, bases and dimension, linear maps with their kernels and images governed by rank--nullity, bilinear forms, and the two flavours of inner product (scalar-with-scalar and scalar-with-group) — constitute the linear-algebraic vocabulary needed for the cryptography ahead. Every commitment is a linear map; every opening is an inner product; every soundness argument rests on (non)degeneracy and dimension counting. The next section applies this machinery to the polynomial spaces \(\F[X]_{<n}\) over evaluation domains, and the elliptic-curve section realises the group \(\G\).
\end{remark}

\clearpage

\section{Roots of unity and evaluation domains}\label{sec:rootsofunity}

The polynomial-based proof systems underlying Halo~2 and Zcash Orchard operate over a single, highly particular kind of object: a finite multiplicative subgroup of a field, generated by a \emph{root of unity}. Nearly every operation such a system performs---committing to a polynomial, checking that two polynomials agree, enforcing that a constraint vanishes everywhere it should, multiplying or interpolating efficiently---is ultimately a computation indexed by such a subgroup. This section develops the theory of these subgroups from first principles. We assume only the basic theory of fields, polynomials over a field, and finite cyclic groups, and we recall these as needed.

Throughout, \(\F\) denotes a field. When \(\F\) is to be finite the notation \(\Fq\) denotes the field with \(q\) elements (so \(q\) is a prime power); when we intend a prime field specifically we use the notation \(\Fp\). We write the multiplicative group of nonzero elements as \(\F^\times = \F \setminus \{0\}\). Set \(\N = \{1,2,3,\dots\}\), and reserve the letter \(n\) for the order of the domain we shall construct.

\subsection{Roots of unity}

The starting point is the equation \(X^n = 1\). Its solutions in a field are the roots of unity of order dividing \(n\).

\begin{definition}\label{rootsofunity:def-root}
Let \(\F\) be a field and \(n \in \N\). An element \(\omega \in \F\) is an \emph{\(n\)-th root of unity} if \(\omega^n = 1\). We write
\[
  \mu_n(\F) = \{\,\omega \in \F : \omega^n = 1\,\}
\]
for the set of all \(n\)-th roots of unity in \(\F\).
\end{definition}

Note that \(0 \notin \mu_n(\F)\), since \(0^n = 0 \ne 1\). Thus every \(n\)-th root of unity is a unit, i.e.\ lies in \(\F^\times\). The set \(\mu_n(\F)\) is not merely a set; it is a group.

\begin{proposition}\label{rootsofunity:prp-group}
\(\mu_n(\F)\) is a finite subgroup of \(\F^\times\) of order at most \(n\).
\end{proposition}

\begin{proof}
The constant \(1\) satisfies \(1^n = 1\), so \(1 \in \mu_n(\F)\). If \(\omega, \zeta \in \mu_n(\F)\) then \((\omega\zeta)^n = \omega^n \zeta^n = 1\cdot 1 = 1\) because \(\F^\times\) is abelian, so \(\omega\zeta \in \mu_n(\F)\). If \(\omega \in \mu_n(\F)\) then \((\omega^{-1})^n = (\omega^n)^{-1} = 1^{-1} = 1\), so \(\omega^{-1} \in \mu_n(\F)\). Hence \(\mu_n(\F)\) is a subgroup of \(\F^\times\). Every element of \(\mu_n(\F)\) is a root of the polynomial \(X^n - 1 \in \F[X]\), which has degree \(n\); a nonzero polynomial of degree \(n\) over a field has at most \(n\) roots, so \(|\mu_n(\F)| \le n\).
\end{proof}

The bound \(|\mu_n(\F)| \le n\) need not be attained. Whether it is attained is precisely the question of the existence of a \emph{primitive} \(n\)-th root of unity.

\begin{definition}\label{rootsofunity:def-primitive}
An element \(\omega \in \F\) is a \emph{primitive \(n\)-th root of unity} if \(\omega^n = 1\) and \(\omega^k \ne 1\) for every integer \(k\) with \(1 \le k < n\). Equivalently, \(\omega\) is primitive of order \(n\) if its multiplicative order \(\ord(\omega)\) equals \(n\).
\end{definition}

Here \(\ord(\omega)\), the \emph{order} of \(\omega\), is the least positive integer \(d\) with \(\omega^d = 1\); it exists for any \(\omega \in \mu_n(\F)\) because \(n\) itself is such an exponent and the positive integers are well-ordered. The two formulations in \cref{rootsofunity:def-primitive} agree: \(\omega^k \ne 1\) for \(1 \le k < n\) together with \(\omega^n = 1\) states exactly that the least positive exponent killing \(\omega\) is \(n\).

\begin{lemma}\label{rootsofunity:lem-orderdivides}
Let \(\omega \in \F^\times\) have order \(d = \ord(\omega)\). Then for any integer \(m\) we have \(\omega^m = 1\) if and only if \(d \mid m\). In particular \(\omega\) is an \(n\)-th root of unity if and only if \(d \mid n\), and \(\omega\) is a \emph{primitive} \(n\)-th root of unity if and only if \(d = n\).
\end{lemma}

\begin{proof}
If \(d \mid m\), write \(m = dk\); then \(\omega^m = (\omega^d)^k = 1^k = 1\). Conversely suppose \(\omega^m = 1\). Divide with remainder: \(m = qd + r\) with \(0 \le r < d\). Then \(1 = \omega^m = (\omega^d)^q \omega^r = \omega^r\). Minimality of \(d\) forces \(r = 0\), so \(d \mid m\). The remaining claims are immediate: \(\omega^n = 1 \iff d \mid n\), and primitivity is the case \(d = n\).
\end{proof}

A primitive \(n\)-th root of unity, when one exists, generates all the others.

\begin{proposition}\label{rootsofunity:prp-cyclic}
If \(\omega \in \F\) is a primitive \(n\)-th root of unity, then
\[
  \mu_n(\F) = \langle \omega \rangle = \{\,1, \omega, \omega^2, \dots, \omega^{n-1}\,\}
\]
is a cyclic group of order exactly \(n\), and these \(n\) powers are pairwise distinct.
\end{proposition}

\begin{proof}
The powers \(1, \omega, \dots, \omega^{n-1}\) are pairwise distinct: if \(\omega^i = \omega^j\) with \(0 \le i \le j \le n-1\), then \(\omega^{j-i} = 1\) with \(0 \le j - i \le n-1 < n\), so by primitivity (\cref{rootsofunity:lem-orderdivides}, with \(d = n\)) we must have \(j - i = 0\). Thus \(\langle \omega \rangle\) has exactly \(n\) elements, all of which lie in \(\mu_n(\F)\) since \((\omega^k)^n = (\omega^n)^k = 1\). But \(|\mu_n(\F)| \le n\) by \cref{rootsofunity:prp-group}, so the inclusion \(\langle\omega\rangle \subseteq \mu_n(\F)\) is an equality of \(n\)-element sets.
\end{proof}

Over the complex numbers a primitive \(n\)-th root of unity always exists, namely \(\omega = e^{2\pi i / n}\); this is the classical picture of \(n\) equally spaced points on the unit circle. Over a finite field the situation is more delicate, and it is exactly the finite case that matters for cryptography.

\subsection{Existence over finite fields: the condition \(n \mid q-1\)}

The decisive structural fact about finite fields is that their multiplicative group is cyclic; the existence theorem for roots of unity rests on it.

\begin{theorem}[Multiplicative group of a finite field is cyclic]\label{rootsofunity:thm-cyclic}
Let \(\Fq\) be a finite field with \(q\) elements. Then \(\Fq^\times\) is a cyclic group of order \(q-1\).
\end{theorem}

\begin{proof}
This is Theorem~\ref{finitefields:thm:cyclic} applied to the finite subgroup \(\Fq^\times\) of itself.
\end{proof}

A generator of \(\Fq^\times\) is classically called a \emph{primitive element} or \emph{primitive root} of \(\Fq\); it is a primitive \((q-1)\)-th root of unity in the sense of \cref{rootsofunity:def-primitive}.

\begin{theorem}[Existence and count of primitive \(n\)-th roots of unity]\label{rootsofunity:thm-existence}
Let \(\Fq\) be a finite field and \(n \in \N\) with \(\gcd(n, q) = 1\) (automatic when \(n < p\), where \(p = \operatorname{char}\Fq\); equivalently, in characteristic \(p\), exactly the condition \(p \nmid n\)). Then:
\begin{enumerate}
\item \(\Fq\) contains a primitive \(n\)-th root of unity if and only if \(n \mid q - 1\).
\item When \(n \mid q-1\), the group \(\mu_n(\Fq)\) is cyclic of order exactly \(n\), and the number of primitive \(n\)-th roots of unity in \(\Fq\) is \(\varphi(n)\), where \(\varphi\) is Euler's totient function.
\end{enumerate}
\end{theorem}

\begin{proof}
Let \(g\) be a generator of \(\Fq^\times\), which is cyclic of order \(q-1\) by \cref{rootsofunity:thm-cyclic}.

\smallskip
\noindent\emph{(\(\Rightarrow\))} Suppose \(\omega\) is a primitive \(n\)-th root of unity in \(\Fq\). Then \(\omega \in \Fq^\times\) has order \(n\) by \cref{rootsofunity:def-primitive}, and by Lagrange's theorem the order of any element divides the group order \(q-1\). Hence \(n \mid q-1\).

\smallskip
\noindent\emph{(\(\Leftarrow\))} Suppose \(n \mid q-1\), say \(q - 1 = n m\). Set \(\omega = g^{m}\). Then \(\omega^n = g^{mn} = g^{q-1} = 1\), so \(\omega \in \mu_n(\Fq)\). Its order is
\[
  \ord(\omega) = \ord(g^m) = \frac{\ord(g)}{\gcd(m, \ord(g))} = \frac{q-1}{\gcd(m, q-1)} = \frac{q-1}{m} = n,
\]
using the standard formula for the order of a power in a cyclic group together with \(m \mid q-1\) (so \(\gcd(m,q-1)=m\)). Thus \(\omega\) is a primitive \(n\)-th root of unity. By \cref{rootsofunity:prp-cyclic}, \(\mu_n(\Fq) = \langle \omega \rangle\) is cyclic of order exactly \(n\).

\smallskip
\noindent\emph{Count.} In the cyclic group \(\mu_n(\Fq) = \langle \omega \rangle\) of order \(n\), the element \(\omega^k\) has order \(n / \gcd(k,n)\). This equals \(n\) precisely when \(\gcd(k,n) = 1\). The number of \(k \in \{0,1,\dots,n-1\}\) with \(\gcd(k,n)=1\) is by definition \(\varphi(n)\). Hence there are exactly \(\varphi(n)\) primitive \(n\)-th roots of unity. (The coprimality hypothesis \(\gcd(n,q)=1\) is what guarantees \(X^n-1\) is separable and that \(\mu_n\) has the full \(n\) elements; if \(p = \operatorname{char}\Fq\) divided \(n\), then \(X^n - 1\) would have repeated roots and \(\mu_n(\Fq)\) would be strictly smaller. But \(n \mid q-1\) already forces \(\gcd(n,q)=1\), since \(\gcd(q-1,q)=1\), so the hypothesis is in fact subsumed.)
\end{proof}

\begin{remark}\label{rootsofunity:rem-subsumed}
The proof shows that for a finite field the clean statement is simply: \emph{\(\Fq\) has a primitive \(n\)-th root of unity \(\iff n \mid q-1\)}, with no separate coprimality side-condition, because \(n \mid q-1\) automatically makes \(n\) coprime to \(q\). We stated the condition \(\gcd(n,q)=1\) explicitly only to flag the role of the characteristic.
\end{remark}

\begin{example}\label{rootsofunity:exp-f7}
Take \(\Fq = \Fp\) with \(p = 7\). Here \(q - 1 = 6\), and \(3\) is a generator of \(\Fp^\times\) (its powers are \(3,2,6,4,5,1\), all six nonzero residues). For \(n = 3 \mid 6\), set \(m = 6/3 = 2\) and \(\omega = 3^2 = 2\). Indeed \(2^1 = 2\), \(2^2 = 4\), \(2^3 = 8 \equiv 1\), so \(\omega = 2\) has order \(3\) and \(\mu_3(\F_7) = \{1, 2, 4\}\). There are \(\varphi(3) = 2\) primitive cube roots, namely \(2\) and \(4\). By contrast \(n = 4 \nmid 6\), so \(\F_7\) has \emph{no} primitive \(4\)-th root of unity: \(X^4 - 1 = (X^2-1)(X^2+1)\) and \(X^2 + 1\) is irreducible over \(\F_7\) (it has no root since \(-1\) is not a square mod \(7\)).
\end{example}

\begin{remark}[Two-adicity and FFT-friendly primes]\label{rootsofunity:rem-twoadic}
For the fast algorithms below it is desirable that \(n\) be a power of two (or have many factors of two and other small primes). \cref{rootsofunity:thm-existence} shows this is possible exactly when \(2^k \mid q-1\). The largest \(k\) with \(2^k \mid q-1\) is called the \emph{two-adicity} of the field. Designers deliberately choose primes \(p\) with high two-adicity---so that \(p - 1 = 2^k \cdot c\) for large \(k\)---as scalar fields in proof systems precisely so that large smooth evaluation domains exist. Both Pasta primes satisfy \(2^{32} \mid p - 1\) (two-adicity exactly \(32\)), so each Pasta field contains primitive \(2^k\)-th roots of unity for every \(k \le 32\), furnishing power-of-two evaluation domains of every size up to \(2^{32}\); this is what makes the entire machinery of this section available to Halo~2 and Orchard.
\end{remark}

\subsection{The evaluation domain \(H\) and its vanishing polynomial}

We now fix the central object of this section.

\begin{definition}\label{rootsofunity:def-domain}
Let \(\F\) be a field containing a primitive \(n\)-th root of unity \(\omega\). The \emph{evaluation domain} of order \(n\) is the multiplicative subgroup
\[
  H = \langle \omega \rangle = \{\,1, \omega, \omega^2, \dots, \omega^{n-1}\,\} = \mu_n(\F) \subseteq \F^\times,
\]
a cyclic group of order \(n\) by \cref{rootsofunity:prp-cyclic}. We frequently index its elements as \(\omega^0, \omega^1, \dots, \omega^{n-1}\), and abbreviate \(\omega^i\) by \(h_i\) when convenient.
\end{definition}

Because \(H\) is a multiplicative \emph{group}, it is closed under products and inverses, and \(1 \in H\). These elementary facts are exactly what make \(H\) substantially more useful than an arbitrary set of \(n\) interpolation points: cosets, products, and shifts of \(H\) interact algebraically. The single most important polynomial attached to \(H\) is the one that vanishes precisely on it.

\begin{definition}\label{rootsofunity:def-vanishing}
The \emph{vanishing polynomial} (or \emph{zerofier}) of \(H\) is
\[
  Z_H(X) = \prod_{h \in H}(X - h) \in \F[X].
\]
\end{definition}

\begin{theorem}[Closed form and factorization of the vanishing polynomial]\label{rootsofunity:thm-vanishing}
With \(H = \mu_n(\F)\) of order \(n\) generated by a primitive \(n\)-th root of unity \(\omega\),
\[
  Z_H(X) \;=\; X^n - 1 \;=\; \prod_{i=0}^{n-1}\bigl(X - \omega^i\bigr).
\]
Moreover \(Z_H\) is separable: it has \(n\) distinct roots and \(\gcd(Z_H, Z_H') = 1\).
\end{theorem}

\begin{proof}
Every \(h \in H\) satisfies \(h^n = 1\), hence is a root of \(X^n - 1\). The \(n\) elements of \(H\) are distinct (\cref{rootsofunity:prp-cyclic}), so \(X^n - 1\) has \(n\) distinct roots, namely all of \(H\). A degree-\(n\) polynomial with \(n\) distinct roots \(h\) is \(c \prod_{h\in H}(X-h)\) for some constant \(c\); comparing leading coefficients (both \(X^n - 1\) and the product are monic) gives \(c = 1\). Hence \(X^n - 1 = \prod_{h\in H}(X-h) = Z_H(X)\). Separability: the derivative is \(Z_H'(X) = n X^{n-1}\). The existence of the \(n\) distinct roots already forces \(\operatorname{char}\F \nmid n\): if \(\operatorname{char}\F = \ell\) divided \(n\), say \(n = \ell m\), the freshman's dream (\cref{finitefields:lem:freshman}) would give \(X^n - 1 = (X^m - 1)^\ell\), a polynomial with at most \(m < n\) distinct roots. Hence \(n \ne 0\) in \(\F\), so \(n X^{n-1}\) is nonzero and its only root is \(0\), which is not a root of \(X^n - 1\). Thus \(Z_H\) and \(Z_H'\) share no root, so \(\gcd(Z_H, Z_H') = 1\) and \(Z_H\) is separable.
\end{proof}

The closed form \(Z_H(X) = X^n - 1\) is the workhorse of polynomial proof systems: testing whether a polynomial \(f\) vanishes on all of \(H\) reduces to checking divisibility by \(X^n - 1\), which one can verify by exhibiting a quotient \(q(X)\) with \(f(X) = q(X)(X^n-1)\). Moreover, evaluating \(Z_H\) at any point costs a single exponentiation \(X^n\), not \(n\) multiplications. We record several structural consequences below.

\begin{proposition}[Cosets and shifts]\label{rootsofunity:prp-coset}
Let \(\gamma \in \F^\times\) and let \(\gamma H = \{\gamma h : h \in H\}\) be the coset. Then the polynomial vanishing on \(\gamma H\) is
\[
  Z_{\gamma H}(X) = \prod_{h \in H}(X - \gamma h) = X^n - \gamma^n.
\]
In particular, distinct cosets of \(H\) in \(\F^\times\) have vanishing polynomials \(X^n - c\) for distinct constants \(c = \gamma^n\), and the cosets of \(H\) partition \(\F^\times\) into translates of \(H\), on each of which \(X^n\) is constant.
\end{proposition}

\begin{proof}
Substitute \(Y = X/\gamma\): \(\prod_{h}(X - \gamma h) = \gamma^n \prod_h (Y - h) = \gamma^n (Y^n - 1) = \gamma^n\bigl((X/\gamma)^n - 1\bigr) = X^n - \gamma^n\), using \cref{rootsofunity:thm-vanishing}. The constant is \(c = \gamma^n\), and \(\gamma_1 H = \gamma_2 H \iff \gamma_1/\gamma_2 \in H \iff (\gamma_1/\gamma_2)^n = 1 \iff \gamma_1^n = \gamma_2^n\), so distinct cosets give distinct \(c\).
\end{proof}

\begin{proposition}[Sum of roots of unity]\label{rootsofunity:prp-sumroots}
If \(n > 1\) then \(\sum_{i=0}^{n-1} \omega^i = 0\), and more generally for any integer \(k\),
\[
  \sum_{i=0}^{n-1} \omega^{ki} = \begin{cases} n & \text{if } n \mid k,\\[2pt] 0 & \text{if } n \nmid k.\end{cases}
\]
\end{proposition}

\begin{proof}
If \(n \mid k\) then \(\omega^{ki} = (\omega^n)^{ki/n} = 1\) for each \(i\), so the sum is \(n\). If \(n \nmid k\) then \(\zeta := \omega^k \ne 1\) (since \(\omega\) has order \(n\), so \(\omega^k = 1 \iff n \mid k\)), and the finite geometric series gives
\[
  \sum_{i=0}^{n-1} \zeta^i = \frac{\zeta^n - 1}{\zeta - 1} = \frac{(\omega^n)^k - 1}{\zeta - 1} = \frac{1 - 1}{\zeta - 1} = 0,
\]
the denominator being nonzero because \(\zeta \ne 1\). The first claim is the case \(k = 1\), valid for \(n > 1\) since then \(n \nmid 1\).
\end{proof}

\subsection{The Lagrange basis on \(H\)}

Attention now passes from the group \(H\) to the polynomials we study on it. The space of polynomials of degree less than \(n\) has dimension \(n\) over \(\F\); since \(H\) has exactly \(n\) points, their values on \(H\) determine such polynomials. The basis adapted to this situation is the Lagrange basis.

\begin{definition}\label{rootsofunity:def-poly}
Let \(\F[X]_{<n}\) denote the \(\F\)-vector space of polynomials of degree strictly less than \(n\), together with the zero polynomial. It has the monomial basis \(\{1, X, \dots, X^{n-1}\}\) and so \(\dim_\F \F[X]_{<n} = n\).
\end{definition}

\begin{definition}\label{rootsofunity:def-lagrange}
For each \(i \in \{0,1,\dots,n-1\}\), the \(i\)-th \emph{Lagrange basis polynomial} for \(H = \{\omega^0,\dots,\omega^{n-1}\}\) is
\[
  L_i(X) \;=\; \prod_{\substack{j=0 \\ j \ne i}}^{n-1} \frac{X - \omega^j}{\omega^i - \omega^j} \;\in\; \F[X]_{<n}.
\]
\end{definition}

\begin{proposition}[Interpolation property]\label{rootsofunity:prp-lagrangeprop}
The polynomials \(L_0, \dots, L_{n-1}\) satisfy the Kronecker-delta conditions
\[
  L_i(\omega^j) = \delta_{ij} = \begin{cases} 1 & i = j,\\ 0 & i \ne j,\end{cases}
\]
they form a basis of \(\F[X]_{<n}\), and they give the unique interpolant: for any data \((v_0,\dots,v_{n-1}) \in \F^n\) the polynomial
\[
  f(X) = \sum_{i=0}^{n-1} v_i\, L_i(X)
\]
is the unique element of \(\F[X]_{<n}\) with \(f(\omega^i) = v_i\) for all \(i\).
\end{proposition}

\begin{proof}
Each \(L_i\) has degree \(n-1\) (a product of \(n-1\) linear factors) and is well-defined because the \(\omega^j\) are distinct, so the denominators \(\omega^i - \omega^j\) are nonzero for \(j \ne i\). Evaluation at \(X = \omega^j\) with \(j \ne i\) makes the factor \((X-\omega^j)\) vanish, giving \(L_i(\omega^j)=0\); at \(X = \omega^i\) every factor is \((\omega^i - \omega^j)/(\omega^i-\omega^j)=1\), giving \(L_i(\omega^i)=1\). This is \(L_i(\omega^j)=\delta_{ij}\).

For the interpolation claim, \(f = \sum_i v_i L_i\) lies in \(\F[X]_{<n}\) and \(f(\omega^j) = \sum_i v_i \delta_{ij} = v_j\). Uniqueness follows because, if \(f, \tilde f \in \F[X]_{<n}\) agree on all \(n\) points of \(H\), then \(f - \tilde f \in \F[X]_{<n}\) has \(n\) distinct roots but degree \(< n\), hence is the zero polynomial. Finally the \(L_i\) are linearly independent---if \(\sum_i c_i L_i = 0\), evaluation at \(\omega^j\) yields \(c_j = 0\)---and there are \(n\) of them in the \(n\)-dimensional space \(\F[X]_{<n}\), so they form a basis.
\end{proof}

On a generic point set the Lagrange polynomials admit no closed form simpler than the defining product. On the \emph{group} \(H\) they collapse to a compact closed form involving the vanishing polynomial.

\begin{theorem}[Closed form of the Lagrange basis on \(H\)]\label{rootsofunity:thm-lagrangeclosed}
For \(H = \mu_n(\F)\) with primitive root \(\omega\), and for each \(i\),
\[
  L_i(X) \;=\; \frac{\omega^i}{n}\cdot\frac{X^n - 1}{X - \omega^i}
  \;=\; \frac{\omega^i\,(X^n-1)}{n\,(X-\omega^i)}.
\]
In particular, dividing out one linear factor and scaling yields all \(n\) Lagrange polynomials from the \emph{single} polynomial \(Z_H(X) = X^n - 1\).
\end{theorem}

\begin{proof}
Fix \(i\). From \cref{rootsofunity:thm-vanishing}, \(Z_H(X) = \prod_{j=0}^{n-1}(X-\omega^j)\), so
\[
  \frac{Z_H(X)}{X-\omega^i} = \prod_{\substack{j=0\\ j\ne i}}^{n-1}(X - \omega^j),
\]
which is exactly the numerator of \(L_i\) in \cref{rootsofunity:def-lagrange}. It remains to evaluate the denominator \(\prod_{j\ne i}(\omega^i-\omega^j)\). This equals \(\lim_{X\to\omega^i} Z_H(X)/(X-\omega^i) = Z_H'(\omega^i)\), the value of the derivative; concretely \(Z_H'(X) = nX^{n-1}\), so
\[
  Z_H'(\omega^i) = n(\omega^i)^{n-1} = n\,\omega^{i n}\,\omega^{-i} = n\cdot 1\cdot \omega^{-i} = \frac{n}{\omega^i},
\]
using \(\omega^{in} = (\omega^n)^i = 1\). Therefore
\[
  L_i(X) = \frac{Z_H(X)/(X-\omega^i)}{Z_H'(\omega^i)} = \frac{(X^n-1)/(X-\omega^i)}{n/\omega^i} = \frac{\omega^i(X^n-1)}{n(X-\omega^i)},
\]
as claimed. (The cancellation \(\prod_{j\ne i}(\omega^i-\omega^j) = Z_H'(\omega^i)\) is the standard fact that, for a polynomial split into distinct linear factors, the value at a root of the product with that factor removed equals the derivative at the root.)
\end{proof}

\begin{remark}[Barycentric form]\label{rootsofunity:rem-barycentric}
The closed form renders the \emph{barycentric} interpolation formula on \(H\) especially clean. Writing \(f(X) = \sum_i v_i L_i(X)\) and substituting \cref{rootsofunity:thm-lagrangeclosed},
\[
  f(X) = \frac{X^n-1}{n}\sum_{i=0}^{n-1}\frac{\omega^i\,v_i}{X - \omega^i}.
\]
Evaluating \(f\) at a point \(z \notin H\) thus costs one evaluation of \(z^n-1\) and \(n\) terms of the sum, with no per-point polynomial reconstruction. This is precisely the formula a verifier uses to evaluate a committed polynomial's Lagrange representation at a random challenge \(z\).
\end{remark}

\begin{proposition}[Lagrange polynomials sum to one]\label{rootsofunity:prp-partition}
\(\sum_{i=0}^{n-1} L_i(X) = 1\) identically.
\end{proposition}

\begin{proof}
The polynomial \(g(X) = \sum_i L_i(X) - 1\) lies in \(\F[X]_{<n}\) (degree \(< n\)) and vanishes at every \(\omega^j\) because \(\sum_i L_i(\omega^j) = \sum_i \delta_{ij} = 1\). A polynomial of degree \(< n\) with \(n\) distinct roots is zero, so \(g \equiv 0\).
\end{proof}

\begin{example}[The all-ones value vector]\label{rootsofunity:exp-constants}
The constant polynomial \(1\) interpolates the data \(v_i = 1\) for all \(i\); \cref{rootsofunity:prp-partition} is the statement \(\sum_i 1\cdot L_i = 1\). Likewise the values of \(X\) on \(H\) are \(v_i = \omega^i\), and \cref{rootsofunity:prp-lagrangeprop} gives \(\sum_i \omega^i L_i(X) = X\) (both sides agree on \(H\) and have degree \(< n\)). We invoke these identities constantly when rewriting monomials in the Lagrange basis.
\end{example}

\subsection{Evaluation and interpolation as mutually inverse linear maps}

An element of \(\F[X]_{<n}\) admits two natural descriptions: by its \(n\) coefficients, or by its \(n\) values on \(H\). Both are coordinate systems on the same \(n\)-dimensional vector space, and passage between them is a linear isomorphism.

\begin{definition}[Evaluation map]\label{rootsofunity:def-evalmap}
The \emph{evaluation map} on \(H\) is the \(\F\)-linear map
\[
  \mathrm{ev}_H \colon \F[X]_{<n} \longrightarrow \F^{\,n}, \qquad
  \mathrm{ev}_H(f) = \bigl(f(\omega^0), f(\omega^1), \dots, f(\omega^{n-1})\bigr).
\]
\end{definition}

Linearity is immediate from \((af+bg)(\omega^i) = a f(\omega^i) + b g(\omega^i)\).

\begin{theorem}[Evaluation and interpolation are mutually inverse]\label{rootsofunity:thm-evaliso}
The evaluation map \(\mathrm{ev}_H\) is a vector-space isomorphism. Its inverse is the \emph{interpolation map}
\[
  \mathrm{int}_H \colon \F^{\,n} \longrightarrow \F[X]_{<n}, \qquad
  \mathrm{int}_H(v_0,\dots,v_{n-1}) = \sum_{i=0}^{n-1} v_i L_i(X).
\]
That is, \(\mathrm{int}_H \circ \mathrm{ev}_H = \mathrm{id}_{\F[X]_{<n}}\) and \(\mathrm{ev}_H \circ \mathrm{int}_H = \mathrm{id}_{\F^n}\).
\end{theorem}

\begin{proof}
Both maps are linear and act between spaces of the same finite dimension \(n\), so it suffices to verify that they compose to the identity in one order; we check both for clarity. For \(v \in \F^n\), \(\mathrm{int}_H(v) = \sum_i v_i L_i\) has value \(\sum_i v_i L_i(\omega^j) = v_j\) at \(\omega^j\) by \cref{rootsofunity:prp-lagrangeprop}, so \(\mathrm{ev}_H(\mathrm{int}_H(v)) = v\). For \(f \in \F[X]_{<n}\), \(\mathrm{int}_H(\mathrm{ev}_H(f)) = \sum_i f(\omega^i) L_i\) is, by the uniqueness clause of \cref{rootsofunity:prp-lagrangeprop}, the unique degree-\(<n\) polynomial taking the values \(f(\omega^i)\) on \(H\)---which is \(f\) itself. Hence the two maps are inverse, and in particular \(\mathrm{ev}_H\) is bijective and linear, i.e.\ an isomorphism.
\end{proof}

In coordinates relative to the monomial basis, \(\mathrm{ev}_H\) is multiplication by a Vandermonde matrix, which we name here because it is the bridge to the Fourier transform.

\begin{definition}[Vandermonde / DFT matrix]\label{rootsofunity:def-vander}
Let \(f(X) = \sum_{k=0}^{n-1} c_k X^k\) have coefficient vector \(c = (c_0,\dots,c_{n-1})\). Then \(\mathrm{ev}_H(f)\) has \(i\)-th entry \(\sum_{k} c_k \omega^{ik}\). Thus, relative to the monomial basis on the domain and the standard basis on \(\F^n\), the \emph{Vandermonde matrix}
\[
  V = \bigl(\omega^{ik}\bigr)_{0\le i,k \le n-1} =
  \begin{pmatrix}
  1 & 1 & 1 & \cdots & 1\\
  1 & \omega & \omega^2 & \cdots & \omega^{n-1}\\
  1 & \omega^2 & \omega^4 & \cdots & \omega^{2(n-1)}\\
  \vdots & & & & \vdots\\
  1 & \omega^{n-1} & \omega^{2(n-1)} & \cdots & \omega^{(n-1)^2}
  \end{pmatrix},
\]
represents \(\mathrm{ev}_H\), so that \(\mathrm{ev}_H(f) = V c\). We call \(V\) the (discrete Fourier) transform matrix on \(H\).
\end{definition}

\begin{theorem}[Invertibility and the inverse transform]\label{rootsofunity:thm-vanderinv}
The matrix \(V\) of \cref{rootsofunity:def-vander} is invertible, with
\[
  (V^{-1})_{k,i} = \frac{1}{n}\,\omega^{-ik}.
\]
Equivalently, \(V^{-1} = \tfrac1n \overline V\) where \(\overline V = (\omega^{-ik})_{i,k}\) is the transform matrix built from \(\omega^{-1}\) (which is itself a primitive \(n\)-th root of unity). Consequently the inversion formula relates the coefficient and value representations by
\[
  c_k = \frac1n \sum_{i=0}^{n-1} f(\omega^i)\,\omega^{-ik}, \qquad
  f(\omega^i) = \sum_{k=0}^{n-1} c_k\,\omega^{ik}.
\]
\end{theorem}

\begin{proof}
It suffices to show \(\overline V V = n I\). The \((k,k')\) entry of \(\overline V V\) is
\[
  \sum_{i=0}^{n-1} \omega^{-ik}\,\omega^{ik'} = \sum_{i=0}^{n-1} \omega^{i(k'-k)}.
\]
By \cref{rootsofunity:prp-sumroots} this sum is \(n\) if \(n \mid (k'-k)\) and \(0\) otherwise. For \(k,k' \in \{0,\dots,n-1\}\) one has \(n \mid (k'-k) \iff k = k'\). Hence \(\overline V V = nI\), giving \(V^{-1} = \tfrac1n \overline V\) (we may divide by \(n\) since \(n \ne 0\) in \(\F\)). The component formulas write out the matrix identities \(\mathrm{ev}_H(f) = Vc\) and \(c = V^{-1}\mathrm{ev}_H(f)\).
\end{proof}

\begin{remark}\label{rootsofunity:rem-isodict}
\cref{rootsofunity:thm-evaliso} and \cref{rootsofunity:thm-vanderinv} are two faces of one statement. The abstract face: evaluation on \(H\) is an isomorphism whose inverse is Lagrange interpolation. The concrete face: that isomorphism is the matrix \(V\), and its inverse is, up to the scalar \(1/n\), the same matrix with \(\omega\) replaced by \(\omega^{-1}\). The near-symmetry between a transform and its inverse---differing only by conjugating \(\omega \mapsto \omega^{-1}\) and a scale \(1/n\)---is the algebraic reason the fast inverse transform costs exactly as much as the forward one.
\end{remark}

\subsection{The number-theoretic Fourier transform}

The map \(\mathrm{ev}_H\), read as the matrix \(V\), is the discrete Fourier transform, but over a finite field rather than over \(\C\). Because it involves no analytic structure (no complex exponentials, no convergence), we call it the \emph{number-theoretic transform}.

\begin{definition}[Number-theoretic transform]\label{rootsofunity:def-ntt}
Let \(\omega\) be a primitive \(n\)-th root of unity in \(\F\). The \emph{number-theoretic transform} (NTT) of a vector \(a = (a_0,\dots,a_{n-1}) \in \F^n\) is the vector \(\hat a = (\hat a_0,\dots,\hat a_{n-1})\) with
\[
  \hat a_j = \sum_{k=0}^{n-1} a_k\,\omega^{jk}, \qquad j = 0,\dots,n-1.
\]
The \emph{inverse number-theoretic transform} is
\[
  a_k = \frac1n \sum_{j=0}^{n-1} \hat a_j\,\omega^{-jk}.
\]
In the language above: if \(a\) is the coefficient vector of \(f \in \F[X]_{<n}\), then \(\hat a = \mathrm{ev}_H(f)\) is its value vector on \(H\); the NTT is evaluation, and the inverse NTT is interpolation.
\end{definition}

That the inverse formula indeed inverts the forward one is exactly \cref{rootsofunity:thm-vanderinv}. The NTT enjoys the same convolution property as the classical DFT, which is the mechanism underlying fast polynomial multiplication.

\begin{theorem}[Convolution theorem]\label{rootsofunity:thm-convolution}
For \(f, g \in \F[X]\) with \(\deg f + \deg g < n\), let \(c = f g\) be their product (degree \(< n\)). Then for every \(h \in H\),
\[
  c(h) = f(h)\,g(h),
\]
i.e.\ pointwise multiplication of value vectors corresponds to polynomial multiplication. Equivalently, writing \(a, b\) for the coefficient vectors of \(f,g\) and \(a * b\) for the (cyclic) convolution defining the coefficients of the product modulo \(X^n-1\),
\[
  \widehat{a * b}_j = \hat a_j\,\hat b_j \qquad (j = 0,\dots,n-1),
\]
so the NTT turns convolution into entrywise multiplication.
\end{theorem}

\begin{proof}
The first identity is the statement that evaluation is a ring homomorphism \(\F[X] \to \F\), \(f \mapsto f(h)\): \((fg)(h) = f(h)g(h)\) for any \(h\). Since \(\deg(fg) < n\), interpolation recovers the product \(c = fg\) from its values on the \(n\)-point set \(H\) (\cref{rootsofunity:prp-lagrangeprop}), so the value vectors multiply entrywise and uniquely determine \(c\). For the convolution statement: multiplication in the quotient ring \(\F[X]/(X^n-1)\) corresponds to cyclic convolution of coefficient vectors, and since \(H = \mu_n\) is exactly the set of roots of \(X^n-1\), the map \(p \mapsto \mathrm{ev}_H(p)\) is a ring isomorphism \(\F[X]/(X^n-1) \xrightarrow{\sim} \F^n\) (with \(\F^n\) under entrywise operations). It sends the product (convolution) of \(a, b\) to the entrywise product of \(\hat a, \hat b\). This is the displayed identity. When \(\deg f + \deg g < n\) the cyclic and ordinary products coincide, so the same statement computes the genuine product \(fg\).
\end{proof}

\begin{remark}[Computational significance]\label{rootsofunity:rem-whyntt}
\cref{rootsofunity:thm-convolution} reduces multiplying two degree-\(<n/2\) polynomials---a quadratic-cost operation by the schoolbook method---to: transform both ($\mathrm{NTT}$), multiply the value vectors entrywise ($n$ multiplications), and transform back (inverse $\mathrm{NTT}$). If the NTT runs in \(O(n\log n)\) operations, the whole multiplication costs \(O(n\log n)\). The fast Fourier transform supplies that conditional.
\end{remark}

\subsection{The fast Fourier transform}

Computing \(\hat a = Va\) directly from \cref{rootsofunity:def-vander} costs \(\Theta(n^2)\) field multiplications: each of the \(n\) outputs is a sum of \(n\) products. The fast Fourier transform (FFT)---here a number-theoretic FFT---computes the same vector in \(O(n\log n)\) operations by a divide-and-conquer recursion, provided \(n\) is a power of two (or more generally smooth). This is where the requirement that \(H\) be a \emph{smooth} subgroup, \cref{rootsofunity:rem-twoadic}, becomes indispensable.

\begin{theorem}[Cooley--Tukey radix-2 step]\label{rootsofunity:thm-fft}
Suppose \(n = 2m\) and \(\omega\) is a primitive \(n\)-th root of unity in \(\F\). Given \(f(X) = \sum_{k=0}^{n-1} a_k X^k\), split its coefficients by parity into
\[
  f_{\mathrm{even}}(Y) = \sum_{k=0}^{m-1} a_{2k}\,Y^k, \qquad
  f_{\mathrm{odd}}(Y) = \sum_{k=0}^{m-1} a_{2k+1}\,Y^k,
\]
so that \(f(X) = f_{\mathrm{even}}(X^2) + X\, f_{\mathrm{odd}}(X^2)\). Then \(\eta := \omega^2\) is a primitive \(m\)-th root of unity, and for \(j = 0,\dots,m-1\),
\begin{align}
  f(\omega^{j}) &= f_{\mathrm{even}}(\eta^{j}) + \omega^{j}\, f_{\mathrm{odd}}(\eta^{j}), \label{rootsofunity:eq-fftbutter1}\\
  f(\omega^{j+m}) &= f_{\mathrm{even}}(\eta^{j}) - \omega^{j}\, f_{\mathrm{odd}}(\eta^{j}). \label{rootsofunity:eq-fftbutter2}
\end{align}
Thus the length-\(n\) NTT of \(a\) reduces to two length-\(m\) NTTs (of the even and odd subsequences, over the domain \(\langle\eta\rangle\)) plus \(O(n)\) extra operations.
\end{theorem}

\begin{proof}
The identity \(f(X) = f_{\mathrm{even}}(X^2) + X f_{\mathrm{odd}}(X^2)\) regroups the sum \(\sum_k a_k X^k\) into even-index and odd-index terms; in the odd terms we pull out one factor of \(X\), leaving \(X^{2k}\). That \(\eta = \omega^2\) is a primitive \(m\)-th root of unity follows from \(\ord(\omega^2) = n/\gcd(2,n) = n/2 = m\) (\cref{rootsofunity:lem-orderdivides} and the order-of-a-power formula), using \(2 \mid n\). It remains to evaluate. For \eqref{rootsofunity:eq-fftbutter1}, put \(X = \omega^j\); then \(X^2 = \omega^{2j} = \eta^{j}\), so \(f(\omega^j) = f_{\mathrm{even}}(\eta^j) + \omega^j f_{\mathrm{odd}}(\eta^j)\). For \eqref{rootsofunity:eq-fftbutter2}, put \(X = \omega^{j+m}\); then \(X^2 = \omega^{2j+2m} = \omega^{2j}\omega^{n} = \eta^{j}\cdot 1 = \eta^j\) (since \(2m = n\) and \(\omega^n=1\)), while the linear factor is \(\omega^{j+m} = \omega^j \omega^m = -\omega^j\). Here \(\omega^m = -1\) because \(\omega^m\) has order \(n/\gcd(m,n) = n/m = 2\), and \(-1\) is the unique element of order \(2\) in a field (it satisfies \(x^2=1\), \(x\ne 1\), so \(x=-1\); note \(-1 \ne 1\) since \(\operatorname{char}\F \nmid n\) and \(n\) is even forces \(\operatorname{char}\F \ne 2\)). Substituting, \(f(\omega^{j+m}) = f_{\mathrm{even}}(\eta^j) - \omega^j f_{\mathrm{odd}}(\eta^j)\).
\end{proof}

\begin{remark}[The butterfly]\label{rootsofunity:rem-butterfly}
Equations \eqref{rootsofunity:eq-fftbutter1}--\eqref{rootsofunity:eq-fftbutter2} constitute the \emph{butterfly}: from the two subtransform outputs \(E_j = f_{\mathrm{even}}(\eta^j)\) and \(O_j = f_{\mathrm{odd}}(\eta^j)\) one forms the \emph{twiddle} \(t_j = \omega^j O_j\) and reads off the two outputs \(E_j + t_j\) and \(E_j - t_j\). Pictorially, \(E_j\) and \(O_j\) sit as two nodes on the left; two crossing edges (forming an \(\times\), the shape that names the operation) carry them to the two outputs on the right, one edge weighted by \(+\omega^j\) and the other by \(-\omega^j\). Each butterfly costs one multiplication (the twiddle) and two additions, and there are \(n/2\) of them per level.
\end{remark}

\begin{theorem}[\(O(n\log n)\) cost of the FFT]\label{rootsofunity:thm-fftcost}
Let \(n = 2^t\) and let \(T(n)\) be the number of field additions and multiplications the radix-2 recursion of \cref{rootsofunity:thm-fft} uses to compute a length-\(n\) NTT. Then
\[
  T(n) = 2\,T(n/2) + cn \quad (n>1), \qquad T(1) = O(1),
\]
for a constant \(c\), and therefore \(T(n) = O(n\log n)\). The inverse NTT has the same cost, since it is the forward NTT with \(\omega\) replaced by \(\omega^{-1}\) followed by scaling each entry by \(1/n\) (\cref{rootsofunity:thm-vanderinv}).
\end{theorem}

\begin{proof}
The recursion is exactly \cref{rootsofunity:thm-fft}: a length-\(n\) transform calls two length-\(n/2\) transforms (on the even and odd coefficient subsequences) and then combines them with \(n/2\) butterflies, each of \(O(1)\) operations, for the \(cn\) overhead. Unrolling, with \(n = 2^t\),
\[
  T(2^t) = 2T(2^{t-1}) + c\,2^t.
\]
Dividing by \(2^t\), and letting \(S(t) = T(2^t)/2^t\), gives \(S(t) = S(t-1) + c\), so \(S(t) = S(0) + ct = O(1) + ct\). Hence \(T(2^t) = 2^t S(t) = O(2^t) + ct\,2^t = O(t\,2^t)\). With \(n = 2^t\), \(t = \log_2 n\), this is \(T(n) = O(n\log n)\). (Equivalently: the recursion tree has \(\log_2 n\) levels, each doing \(\Theta(n)\) total work across all subproblems at that level.) The inverse-transform claim is \cref{rootsofunity:thm-vanderinv}: the same Cooley--Tukey recursion with \(\omega^{-1}\) in place of \(\omega\), and a final pass of \(n\) divisions by \(n\), all within \(O(n\log n)\).
\end{proof}

\begin{remark}[Mixed radix and smoothness]\label{rootsofunity:rem-mixedradix}
The hypothesis that enables fast transforms is that \(n\) be \emph{smooth}---a product of small primes---and, crucially, that the field actually \emph{contain} a primitive \(n\)-th root of unity, i.e.\ \(n \mid q-1\) (\cref{rootsofunity:thm-existence}). A power-of-two \(n\) dividing \(q-1\) is the optimal case: pure radix-2 with the cleanest butterfly.
\end{remark}

\begin{example}[A length-\(4\) transform]\label{rootsofunity:exp-fft4}
Take \(n = 4\) with primitive \(4\)-th root \(\omega\) (so \(\omega^2 = -1\), \(\omega^4 = 1\), \(\eta = \omega^2 = -1\) the primitive \(2\)nd root). For \(f = a_0 + a_1 X + a_2 X^2 + a_3 X^3\), set \(E = (a_0 + a_2,\; a_0 - a_2)\) and \(O = (a_1 + a_3,\; a_1 - a_3)\), the two length-\(2\) transforms over \(\{1,-1\}\). Then the butterflies give
\[
  \hat a_0 = E_0 + \omega^0 O_0,\quad \hat a_2 = E_0 - \omega^0 O_0,\quad
  \hat a_1 = E_1 + \omega^1 O_1,\quad \hat a_3 = E_1 - \omega^1 O_1.
\]
This uses \(2\) sub-transforms of \(2\) additions each plus \(4\) butterfly operations---far fewer than the \(16\) products of the direct \(4\times4\) matrix multiply, and the gap widens as \(n\) grows.
\end{example}

\subsection{The role of smooth multiplicative domains in polynomial proof systems}

This section closes by assembling the pieces into the reason these domains form the substrate of Halo~2 and Orchard. A proof system of the polynomial type encodes a computation as a system of polynomial identities and convinces a verifier that those identities hold. The evaluation domain \(H = \mu_n\) is where the encoding lives, and every property proved above is load-bearing.

\begin{remark}[The roles played by \(H\), \(Z_H\), the Lagrange basis, and the FFT]\label{rootsofunity:rem-whyproofsystems}
\leavevmode
\begin{itemize}
\item \textbf{Witnesses as evaluations on \(H\).} A prover lays out the trace of a computation as the values of one or more polynomials on the \(n\) points of \(H\). By \cref{rootsofunity:thm-evaliso} this is the same data as a polynomial of degree \(< n\); the prover may freely move between the value (Lagrange) representation, natural for stating constraints row by row, and the coefficient representation, natural for committing and for low-degree testing. The change of representation is the (inverse) NTT.

\item \textbf{Constraints as divisibility by \(Z_H\).} ``Constraint polynomial \(C\) holds on every row'' means \(C\) vanishes on all of \(H\), which by \cref{rootsofunity:thm-vanishing} means exactly that \(Z_H = X^n - 1\) divides \(C\). The prover demonstrates this by exhibiting a quotient \(t\) with \(C(X) = t(X)\,(X^n - 1)\). The inexpensive closed form \(Z_H(X) = X^n-1\)---one exponentiation to evaluate---is what makes this check efficient at a random point. \cref{rootsofunity:prp-coset} extends the same technique to cosets \(\gamma H\), used to separate quotient and remainder domains.

\item \textbf{Lagrange selectors.} The Lagrange polynomials \(L_i\) of \cref{rootsofunity:def-lagrange} are \emph{selectors}: \(L_i\) is \(1\) on row \(i\) and \(0\) elsewhere (\cref{rootsofunity:prp-lagrangeprop}), so boundary/instance constraints (e.g.\ ``the first row equals the public input'') become \(L_0 \cdot(\text{something})\). The closed form and barycentric formula (\cref{rootsofunity:thm-lagrangeclosed}, \cref{rootsofunity:rem-barycentric}) let the verifier evaluate these at a random challenge in \(O(n)\) or even \(O(\log n)\) field operations rather than reconstructing polynomials.

\item \textbf{Group structure for permutations and shifts.} Because \(H\) is a cyclic \emph{group}, multiplication by \(\omega\) is the cyclic ``next-row'' shift \(\omega^i \mapsto \omega^{i+1}\). Constraints relating a row to its neighbour become \(f(\omega X)\) versus \(f(X)\); permutation arguments (the heart of Halo~2's copy constraints) rest on the action of the group on itself and on its cosets. An arbitrary set of \(n\) points offers none of this---it requires the multiplicative \emph{group} \(H\) of \cref{rootsofunity:def-domain}.

\item \textbf{Smoothness for speed.} Provers manipulate polynomials of size \(n\) (often after blowing the domain up by a small constant factor for low-degree testing). Every transform between representations, every multiplication, every coset evaluation is an NTT. By \cref{rootsofunity:thm-fftcost} these cost \(O(n\log n)\) instead of \(O(n^2)\)---the difference between feasible and infeasible at the sizes used in practice---\emph{provided} \(n\) is smooth and \(\Fq\) contains a primitive \(n\)-th root of unity. By \cref{rootsofunity:thm-existence} that requires \(n \mid q-1\); by \cref{rootsofunity:rem-twoadic} this is why we choose high-two-adicity primes for the scalar fields of the Pallas and Vesta curves.
\end{itemize}
In summary: the existence condition \(n \mid q-1\) (\cref{rootsofunity:thm-existence}) supplies the group \(H\); the group supplies the clean vanishing polynomial (\cref{rootsofunity:thm-vanishing}), the closed-form Lagrange basis (\cref{rootsofunity:thm-lagrangeclosed}), and the evaluation/interpolation isomorphism (\cref{rootsofunity:thm-evaliso}); and smoothness of \(n\) supplies the \(O(n\log n)\) FFT (\cref{rootsofunity:thm-fftcost}). Together these are exactly the ingredients a polynomial proof system requires, which is why such systems rest on smooth multiplicative evaluation domains and on the fields that contain them.
\end{remark}

\clearpage

\section{Elliptic curves}\label{sec:ellipticcurves}

Having developed the theory of finite fields, we now turn attention to the geometric objects that supply the cyclic groups underlying the cryptography of Halo~2 and Zcash Orchard: \emph{elliptic curves}. An elliptic curve is, concretely, the solution set of a certain cubic equation in two variables, together with one extra ``point at infinity''. The notable feature of elliptic curves is that this solution set carries the structure of an abelian group, defined by a strikingly geometric rule, and that this group is hard to compute in: solving the discrete logarithm problem in a well-chosen elliptic curve group appears to require exponential time. This combination of a clean algebraic structure with a hard computational problem is exactly what cryptography requires.

Throughout this section we work over a field \(K\) whose characteristic is neither \(2\) nor \(3\). In the applications \(K\) will be a finite prime field \(\Fp\) with \(p\) a large prime, but it is cleaner to develop the basic theory over a general field and specialize when needed. The reader should keep in mind the two \emph{Pasta curves}
\[
  E_p : y^2 = x^3 + 5 \quad\text{over } \Fp, \qquad
  E_q : y^2 = x^3 + 5 \quad\text{over } \Fq,
\]
where \(p\) and \(q\) are the two \(255\)-bit primes
\begin{align*}
  p &= 2^{254} + 0x224698fc094cf91b992d30ed00000001,\\
  q &= 2^{254} + 0x224698fc0994a8dd8c46eb2100000001,
\end{align*}
already encountered in the finite-fields section (Remark~\ref{finitefields:rem:pasta}). These are the curves \emph{Pallas} (\(E_p\), defined over \(\Fp\)) and \emph{Vesta} (\(E_q\), defined over \(\Fq\)); together they form a \emph{cycle}, in that the number of \(\Fp\)-points of Pallas equals \(q\) and the number of \(\Fq\)-points of Vesta equals \(p\). The equation \(y^2 = x^3 + 5\) serves as the running example, with a smaller toy curve developed in parallel for hand computation.

\subsection{Weierstrass equations}

\begin{definition}[short Weierstrass equation]\label{ellipticcurves:def-weierstrass}
Let \(K\) be a field of characteristic \(\neq 2, 3\), and let \(a, b \in K\). The \emph{short Weierstrass equation} with coefficients \(a, b\) is the equation
\begin{equation}\label{ellipticcurves:eq-weierstrass}
  y^2 = x^3 + a x + b .
\end{equation}
The elements \(a\) and \(b\) are the \emph{Weierstrass coefficients}.
\end{definition}

Over a field of characteristic \(\neq 2, 3\), an invertible affine change of variables always reduces the most general cubic curve \(y^2 + a_1 xy + a_3 y = x^3 + a_2 x^2 + a_4 x + a_6\) to the form~\eqref{ellipticcurves:eq-weierstrass}: completing the square in \(y\) (possible since \(2 \neq 0\)) removes the \(xy\) and \(y\) terms, and then a shift \(x \mapsto x - a_2/3\) (possible since \(3 \neq 0\)) removes the \(x^2\) term. This is precisely why we exclude characteristics \(2\) and \(3\); those cases require the longer Weierstrass form and are not relevant to the Pasta curves, which live over large prime fields.

Not every choice of \(a, b\) gives a usable curve. We must exclude the cubics whose graph has a self-intersection or a cusp, because the group law (defined below) breaks down at such a singular point.

\begin{definition}[discriminant and nonsingularity]\label{ellipticcurves:def-discriminant}
The \emph{discriminant} of the short Weierstrass equation~\eqref{ellipticcurves:eq-weierstrass} is
\begin{equation}\label{ellipticcurves:eq-discriminant}
  \Delta = -16\,(4 a^3 + 27 b^2) \in K .
\end{equation}
The equation, and the cubic \(f(x) = x^3 + ax + b\), is \emph{nonsingular} (or \emph{smooth}) if \(\Delta \neq 0\), equivalently \(4 a^3 + 27 b^2 \neq 0\).
\end{definition}

\begin{proposition}\label{ellipticcurves:prop-smooth}
Let \(f(x) = x^3 + ax + b \in K[x]\). The following are equivalent.
\begin{enumerate}
  \item \(4a^3 + 27 b^2 \neq 0\).
  \item \(f\) has no repeated root in the algebraic closure \(\overline{K}\).
  \item \(f\) and its derivative \(f'\) are coprime in \(\overline{K}[x]\).
  \item There is no point \((x_0, y_0)\) in \(\overline{K}^2\) satisfying both \(y^2 = f(x)\) and the two partial-derivative equations \(2y = 0\) and \(f'(x) = 0\).
\end{enumerate}
\end{proposition}

\begin{proof}
\((2)\Leftrightarrow(3)\): a polynomial has a repeated root iff it shares a root with its derivative, i.e. iff \(\gcd(f, f') \neq 1\) in \(\overline{K}[x]\). \((1)\Leftrightarrow(2)\): the discriminant of the cubic \(x^3 + ax + b\) equals, up to the nonzero factor \(-16\), \(4a^3 + 27 b^2\); a univariate polynomial has a repeated root iff its discriminant vanishes. One may verify the identity directly: if \(f(x) = (x-r_1)(x-r_2)(x-r_3)\) then \(\prod_{i<j}(r_i - r_j)^2 = -(4a^3 + 27b^2)\), using \(r_1 + r_2 + r_3 = 0\) (the coefficient of \(x^2\) is zero). \((4)\Leftrightarrow(2)\): a singular point of the affine curve \(y^2 = f(x)\) is a point where both partial derivatives of \(g(x,y) = y^2 - f(x)\) vanish, namely \(\partial g/\partial y = 2y = 0\) and \(\partial g/\partial x = -f'(x) = 0\). Since \(\mathrm{char}(K) \neq 2\), the first forces \(y_0 = 0\), and then the curve equation gives \(f(x_0) = 0\); combined with \(f'(x_0) = 0\) this says \(x_0\) is a repeated root of \(f\). Conversely a repeated root \(x_0\) yields the singular point \((x_0, 0)\).
\end{proof}

\begin{definition}[elliptic curve, affine points]\label{ellipticcurves:def-curve}
An \emph{elliptic curve} over \(K\) is a short Weierstrass equation~\eqref{ellipticcurves:eq-weierstrass} with nonzero discriminant, \(4a^3 + 27 b^2 \neq 0\). For any field extension \(L \supseteq K\) the set of \emph{\(L\)-rational affine points} is
\[
  E^{\mathrm{aff}}(L) = \{\, (x,y) \in L \times L : y^2 = x^3 + ax + b \,\} ,
\]
and the set of \emph{\(L\)-rational points} is
\[
  E(L) = E^{\mathrm{aff}}(L) \cup \{\mathcal{O}\},
\]
where \(\mathcal{O}\) is a single formal symbol called the \emph{point at infinity}. When \(L = K\) one writes simply \(E\) for the curve and \(E(K)\) for its rational points.
\end{definition}

\begin{remark}\label{ellipticcurves:rem-pasta-smooth}
For the Pasta curves \(a = 0\), \(b = 5\), so \(4a^3 + 27 b^2 = 27 \cdot 25 = 675\). As long as the characteristic does not divide \(675 = 3^3 \cdot 5^2\), i.e. the prime is not \(3\) or \(5\), the curve is nonsingular. The Pasta primes \(p, q\) are enormous, so this is automatic. More generally any curve \(y^2 = x^3 + b\) with \(b \neq 0\) is nonsingular away from characteristics \(2, 3\).
\end{remark}

\subsection{The point at infinity, geometrically}

To motivate the point at infinity, view the affine plane as sitting inside the \emph{projective plane}. A point of the projective plane \(\mathbb{P}^2(K)\) is an equivalence class \([X : Y : Z]\) of nonzero triples \((X, Y, Z) \in K^3 \setminus \{0\}\), where \((X,Y,Z) \sim (\lambda X, \lambda Y, \lambda Z)\) for every \(\lambda \in K^{\times}\). The affine points \((x, y)\) embed as \([x : y : 1]\); the points with \(Z = 0\), the \emph{line at infinity}, are the extra ``directions''. Homogenizing~\eqref{ellipticcurves:eq-weierstrass} by setting \(x = X/Z\), \(y = Y/Z\) and clearing denominators gives the projective Weierstrass equation
\begin{equation}\label{ellipticcurves:eq-projective}
  Y^2 Z = X^3 + a X Z^2 + b Z^3 .
\end{equation}
Setting \(Z = 0\) forces \(X^3 = 0\), hence \(X = 0\), leaving the single projective point \([0 : 1 : 0]\). This is the point at infinity \(\mathcal{O}\): it lies ``infinitely far up'' in the vertical direction, and every vertical line \(x = \text{const}\) passes through it. As we shall see, \(\mathcal{O}\) is the identity element of the group law. The reader may safely picture \(\mathcal{O}\) as a single point at the top (and bottom) of the \(y\)-axis where all vertical lines meet.

\subsection{The chord-and-tangent group law: geometry}

We now describe the procedure for ``adding'' two points of \(E(K)\). The construction is purely geometric and rests on a single fact: a line meets a cubic curve in exactly three points, counted with multiplicity. The point at infinity is engineered precisely so that this counting is always exact.

\begin{definition}[the group law, geometric form]\label{ellipticcurves:def-grouplaw-geom}
Let \(P, Q \in E(K)\). Define \(P + Q\) as follows.
\begin{itemize}
  \item If \(P = \mathcal{O}\), set \(P + Q = Q\); if \(Q = \mathcal{O}\), set \(P + Q = P\).
  \item Otherwise \(P = (x_1, y_1)\) and \(Q = (x_2, y_2)\) are affine. Draw the line \(\ell\) through \(P\) and \(Q\); if \(P = Q\), let \(\ell\) be the tangent line to the curve at \(P\). The line \(\ell\) meets the curve in a third point \(R'\) (counted with multiplicity), where if \(\ell\) is vertical the third point is \(\mathcal{O}\). Define \(P + Q\) to be the \emph{reflection of \(R'\) across the \(x\)-axis}: if \(R' = (x_3, -y_3)\) then \(P + Q = (x_3, y_3)\), and if \(R' = \mathcal{O}\) then \(P + Q = \mathcal{O}\).
\end{itemize}
\end{definition}

The reflection step is what converts ``the three intersection points sum to a fixed thing'' into a genuine group law with \(\mathcal{O}\) as identity. Concretely, the rule is: the three intersection points of any line with the curve add up to \(\mathcal{O}\). The reflection across the \(x\)-axis sends an affine point \((x, y)\) to \((x, -y)\), which (as we shall see) is its additive inverse; and it sends \(\mathcal{O}\) to \(\mathcal{O}\).

To convert this geometry into formulas one must (i) find the third intersection point of a line with the cubic and (ii) handle the degenerate cases. The key algebraic device is Vieta's formula relating the roots of a cubic to its coefficients.

\subsection{Explicit affine formulas}

Fix an elliptic curve \(E : y^2 = x^3 + ax + b\) over \(K\). Let \(P = (x_1, y_1)\) and \(Q = (x_2, y_2)\) be affine points.

\paragraph{Negation.} Define
\[
  -\mathcal{O} = \mathcal{O}, \qquad -(x_1, y_1) = (x_1, -y_1).
\]
Note \((x_1, -y_1)\) is again on the curve since the equation depends on \(y\) only through \(y^2\). The line through \(P\) and \(-P\) is the vertical line \(x = x_1\), which meets the curve at \(P\), \(-P\), and \(\mathcal{O}\).

\paragraph{Addition (the generic chord).} Suppose \(P, Q \neq \mathcal{O}\) and \(x_1 \neq x_2\). The line through \(P\) and \(Q\) is \(y = \lambda x + \mu\) with
\[
  \lambda = \frac{y_2 - y_1}{x_2 - x_1}, \qquad \mu = y_1 - \lambda x_1 .
\]
Substituting into the curve equation, \((\lambda x + \mu)^2 = x^3 + a x + b\), gives
\[
  x^3 - \lambda^2 x^2 + (a - 2\lambda\mu) x + (b - \mu^2) = 0 .
\]
This cubic in \(x\) has the three roots \(x_1, x_2, x_3\), where \(x_3\) is the \(x\)-coordinate of the third intersection point \(R'\). By Vieta's formula the sum of the roots equals the negative of the \(x^2\)-coefficient, so \(x_1 + x_2 + x_3 = \lambda^2\). Hence
\begin{equation}\label{ellipticcurves:eq-add}
  x_3 = \lambda^2 - x_1 - x_2, \qquad y_3 = \lambda (x_1 - x_3) - y_1,
\end{equation}
where the \(y\)-formula already incorporates the reflection across the \(x\)-axis: the third intersection point is \((x_3, \lambda x_3 + \mu)\), and reflecting gives \(y_3 = -(\lambda x_3 + \mu) = -(\lambda x_3 + y_1 - \lambda x_1) = \lambda(x_1 - x_3) - y_1\). Then \(P + Q = (x_3, y_3)\).

\paragraph{Doubling (the tangent).} Suppose \(P = Q = (x_1, y_1)\) with \(y_1 \neq 0\). Implicit differentiation of \(y^2 = x^3 + ax + b\) yields the slope of the tangent line: \(2 y \, dy = (3x^2 + a)\, dx\), so
\[
  \lambda = \frac{3 x_1^2 + a}{2 y_1} .
\]
The same Vieta computation now uses a double root at \(x_1\), giving \(x_1 + x_1 + x_3 = \lambda^2\), hence
\begin{equation}\label{ellipticcurves:eq-double}
  x_3 = \lambda^2 - 2 x_1, \qquad y_3 = \lambda(x_1 - x_3) - y_1, \qquad [2]P = (x_3, y_3).
\end{equation}

\paragraph{Edge cases.} The remaining cases are exactly where the line is vertical:
\begin{itemize}
  \item If \(P = \mathcal{O}\), then \(P + Q = Q\); symmetrically if \(Q = \mathcal{O}\).
  \item If \(x_1 = x_2\) but \(y_1 = -y_2\) (so \(Q = -P\), including the subcase \(y_1 = y_2 = 0\)), the line is vertical, the third point is \(\mathcal{O}\), and \(P + Q = \mathcal{O}\).
  \item If \(P = Q\) and \(y_1 = 0\), the tangent is vertical (the derivative slope formula has zero denominator), so \([2]P = \mathcal{O}\). These are exactly the points of order two.
\end{itemize}
Every case is covered: either one input is \(\mathcal{O}\); or both are affine with \(x_1 \neq x_2\) (use~\eqref{ellipticcurves:eq-add}); or both are affine with \(x_1 = x_2\), in which case either \(y_1 = -y_2\) (result \(\mathcal{O}\)) or \(y_1 = y_2 \neq 0\) (use~\eqref{ellipticcurves:eq-double}).

\begin{proposition}[closure]\label{ellipticcurves:prop-closure}
For all \(P, Q \in E(K)\), the point \(P + Q\) computed above lies in \(E(K)\). In particular \(E(K)\) is closed under \(+\).
\end{proposition}

\begin{proof}
If either input is \(\mathcal{O}\) the result is the other input, which lies in \(E(K)\) by hypothesis; if the result is \(\mathcal{O}\) it lies in \(E(K)\) by definition. Otherwise \((x_3, y_3)\) follows from~\eqref{ellipticcurves:eq-add} or~\eqref{ellipticcurves:eq-double}. By construction \((x_3, -y_3)\) (the third intersection point \(R'\)) satisfies both the line equation and the curve equation, hence lies on the curve; since the curve equation is invariant under \(y \mapsto -y\), the reflected point \((x_3, y_3)\) is on the curve too, so \(P+Q \in E^{\mathrm{aff}}(K) \subseteq E(K)\). All coordinates are rational expressions in \(x_i, y_i, a\) with denominators \(x_2 - x_1\) or \(2y_1\) that are nonzero in the relevant cases, so \(x_3, y_3 \in K\).
\end{proof}

\begin{remark}[Complete versus incomplete addition]\label{ellipticcurves:rem-complete}
The chord formula~\eqref{ellipticcurves:eq-add} is \emph{incomplete}: it fails whenever \(x_1 = x_2\) (the doubling and inverse cases), and an implementation must branch into the case analysis above. \emph{Complete} addition formulas --- single rational expressions valid for \emph{all} pairs of inputs, including \(P = Q\) and \(P = -Q\) --- exist for the curves used here, at the cost of more field multiplications per addition. The distinction is load-bearing in this monograph's later volumes: Halo~2's elliptic-curve gadgets use complete formulas where an adversarially chosen input could otherwise hit an exceptional case, while the Sinsemilla hash function of Orchard deliberately uses the \emph{incomplete} formulas and turns their exceptional cases into constraint violations, so that an exceptional case renders the statement unprovable rather than silently wrong; soundness thus rests on the proof system, while honest inputs avoid the exceptional cases except with negligible probability.
\end{remark}
\subsection{Identity, inverses, commutativity, and associativity}

\begin{proposition}[identity and inverses]\label{ellipticcurves:prop-id-inv}
The operation \(+\) on \(E(K)\) has the following properties.
\begin{enumerate}
  \item \(\mathcal{O}\) is a two-sided identity: \(P + \mathcal{O} = \mathcal{O} + P = P\) for all \(P\).
  \item Every \(P\) has a two-sided inverse \(-P\): \(P + (-P) = \mathcal{O}\), where \(-\mathcal{O} = \mathcal{O}\) and \(-(x,y) = (x,-y)\).
  \item The operation is commutative: \(P + Q = Q + P\).
\end{enumerate}
\end{proposition}

\begin{proof}
(1) follows immediately from the definition. (2): for \(P = (x,y)\) with \(y \neq 0\), the line through \(P\) and \(-P = (x,-y)\) is vertical and meets the curve at \(\mathcal{O}\) as the third point, so \(P + (-P) = \mathcal{O}\); for \(y = 0\), \(P = -P\) and \([2]P = \mathcal{O}\) directly; and \(\mathcal{O} + \mathcal{O} = \mathcal{O}\). (3): the chord through \(P\) and \(Q\) does not depend on their order, and neither does the tangent in the doubling case; algebraically, formulas~\eqref{ellipticcurves:eq-add} are symmetric under the swap \((x_1,y_1) \leftrightarrow (x_2,y_2)\) (the slope \(\lambda\) is unchanged, and \(x_3 = \lambda^2 - x_1 - x_2\) is symmetric).
\end{proof}

\begin{theorem}[group law]\label{ellipticcurves:thm-grouplaw}
Let \(E\) be an elliptic curve over a field \(K\). The set \(E(K)\) with the operation \(+\) of Definition~\ref{ellipticcurves:def-grouplaw-geom} is an abelian group with identity \(\mathcal{O}\).
\end{theorem}

The only group axiom not yet established is \emph{associativity}, \((P + Q) + R = P + (Q + R)\). This is the deep and surprising part of the theory: the definition gives no reason why reflecting third intersection points should be an associative operation, and a brute-force verification through the explicit formulas requires splitting into a large number of coordinate cases and checking lengthy rational-function identities. We sketch below the conceptual proof, which exhibits the reason associativity holds.

\begin{proof}[Proof of associativity (sketch)]
\emph{Sketch.} The proof uses the geometry of cubic curves. The key tool is the following \emph{Cayley--Bacharach}-type fact: if two cubic curves \(C_1, C_2\) in \(\mathbb{P}^2\) meet in exactly nine points, and a third cubic \(C\) passes through eight of those nine points, then \(C\) passes through the ninth as well. (This follows from a dimension count: cubics form a \(9\)-dimensional projective space of coefficients, and passing through a point is one linear condition; eight general points impose independent conditions, pinning down the pencil through them.)

To prove \((P + Q) + R = P + (Q + R)\), one writes both sides via the chord-and-tangent construction and arranges eight common intersection points among three suitably chosen lines (a degenerate cubic, namely a product of three lines) and the curve \(E\); the Cayley--Bacharach principle then forces the two candidate sums to share their final coordinate, hence to be equal. Translating ``reflect across the \(x\)-axis'' into ``the three points on a line sum to \(\mathcal{O}\)'' renders the bookkeeping uniform. We obtain a complete, case-free treatment by identifying \(E(K)\) with the degree-zero part \(\mathrm{Pic}^0(E)\) of the divisor class group of the curve, where the group law is manifestly associative because it is inherited from addition of divisors modulo principal divisors; the map \(P \mapsto [P] - [\mathcal{O}]\) is a bijection \(E(K) \to \mathrm{Pic}^0(E)\) intertwining the two operations. Either route yields a genuine proof; both lie beyond the elementary scope here, so we take associativity as established.
\end{proof}

\begin{remark}\label{ellipticcurves:rem-notation}
For \(n \in \Z\) and \(P \in E(K)\) we write \([n]P\) for the \(n\)-fold sum \(P + \cdots + P\) (with \([0]P = \mathcal{O}\) and \([-n]P = -([n]P)\)). This makes \(E(K)\) a module over \(\Z\), i.e. its abelian-group structure written multiplicatively-by-integers. Note that \([n]P\) is the \emph{scalar multiplication} developed below (Definition~\ref{ellipticcurves:def-dbladd}), not a coordinate-wise product.
\end{remark}

\subsection{Projective and Jacobian coordinates}

The affine formulas~\eqref{ellipticcurves:eq-add} and~\eqref{ellipticcurves:eq-double} each require a field \emph{inversion} (dividing by \(x_2 - x_1\) or \(2 y_1\)). In a large prime field \(\Fp\), an inversion costs roughly as much as dozens of multiplications (the implementation typically computes it via the extended Euclidean algorithm or as \(z^{p-2}\) by Fermat). When a long chain of additions arises, as in scalar multiplication, it is advantageous to defer all inversions to the very end by carrying points in a redundant coordinate system where the group law uses only additions, subtractions, and multiplications.

\begin{remark}[Inversion-free coordinate systems]\label{ellipticcurves:rem-coords}
Two redundant coordinate systems serve this purpose. In \emph{(homogeneous) projective coordinates}, a triple \([X : Y : Z]\) with \(Z \neq 0\) represents the affine point \((X/Z,\, Y/Z)\) on the projective curve~\eqref{ellipticcurves:eq-projective}, with \(\mathcal O = [0:1:0]\). In \emph{Jacobian coordinates}, \((X : Y : Z)\) with \(Z \neq 0\) represents \((X/Z^2,\, Y/Z^3)\), under the equivalence \((X,Y,Z) \sim (\lambda^2 X, \lambda^3 Y, \lambda Z)\), the curve equation becoming \(Y^2 = X^3 + aXZ^4 + bZ^6\). In either system the addition and doubling formulas can be rewritten with no divisions at all, at the cost of more multiplications per operation; a long chain of group operations then defers all inversions to a single final recovery of affine coordinates (compute \(Z^{-1}\) once, then the needed powers by multiplication). Jacobian coordinates usually yield the faster doubling and are the standard choice in cryptographic libraries, including the implementations used by Halo~2. For curves \(y^2 = x^3 + b\) with \(a = 0\), the term \(aZ^4\) in the Jacobian doubling formulas vanishes, making doubling especially cheap --- one reason the Pasta curves take \(a = 0\).
\end{remark}

\subsection{Scalar multiplication and double-and-add}

The fundamental operation in elliptic-curve cryptography is \emph{scalar multiplication}: given a point \(P\) and an integer \(n \geq 0\), compute
\[
  [n]P = \underbrace{P + P + \cdots + P}_{n\ \text{times}} .
\]
Computing this by \(n - 1\) successive additions is infeasible when \(n\) has hundreds of bits. Instead one uses the binary expansion of \(n\), exactly as in fast modular exponentiation.

\begin{definition}[double-and-add]\label{ellipticcurves:def-dbladd}
Write \(n = \sum_{i=0}^{k-1} n_i 2^i\) with \(n_i \in \{0,1\}\) and \(n_{k-1} = 1\). The \emph{left-to-right double-and-add} algorithm computes \([n]P\) as follows: set \(R \leftarrow \mathcal{O}\); for \(i = k-1\) down to \(0\), set \(R \leftarrow [2]R\), and if \(n_i = 1\) set \(R \leftarrow R + P\); return \(R\).
\end{definition}

\begin{proposition}\label{ellipticcurves:prop-dbladd}
Double-and-add returns \([n]P\) using at most \(k - 1\) doublings and at most \(k - 1\) additions, where \(k = \lfloor \log_2 n\rfloor + 1\) is the bit length of \(n\).
\end{proposition}

\begin{proof}
Let \(R_j\) be the value of \(R\) after processing bits \(n_{k-1}, \ldots, n_j\). We claim \(R_j = [\,\lfloor n / 2^{j}\rfloor\,]P\), and prove this by downward induction on \(j\). For the top bit, after \(i = k-1\) one has \(R = [2]\mathcal{O} + n_{k-1} P = [n_{k-1}]P = [1]P = [\lfloor n/2^{k-1}\rfloor]P\). Assuming \(R_{j+1} = [\lfloor n/2^{j+1}\rfloor]P\), the next iteration computes \([2]R_{j+1} + [n_j]P = [\,2\lfloor n/2^{j+1}\rfloor + n_j\,]P = [\lfloor n/2^{j}\rfloor]P\), since \(2\lfloor n/2^{j+1}\rfloor + n_j = \lfloor n/2^j\rfloor\). At \(j = 0\) one obtains \(R_0 = [n]P\). The loop runs \(k-1\) times after the initial top bit is consumed, giving the stated counts (the doubling of \(\mathcal{O}\) and addition for the top bit are free/trivial).
\end{proof}

Thus scalar multiplication by an \(n\)-bit scalar costs \(O(n)\) group operations --- hence \(O(n)\) field operations up to constant factors, or \(O(n \log^2 p)\) bit operations with schoolbook field arithmetic. This efficiency in the ``forward'' direction, contrasted with the apparent hardness of the inverse problem (\S\ref{ellipticcurves:subsec-ecdlp}), is the foundation of the cryptography.

\subsection{The group structure of \(E(\Fp)\)}

Now specialize \(K = \Fp\), the finite field with \(p\) elements (\(p\) prime, \(p > 3\)). Since \(\Fp\) is finite, so is \(E(\Fp)\): it has at most \(2p + 1\) points (for each of the \(p\) values of \(x\), at most two values of \(y\), plus \(\mathcal{O}\)). Theorem~\ref{ellipticcurves:thm-grouplaw} makes \(E(\Fp)\) a finite abelian group. One of the central theorems of the subject governs its order.

\begin{theorem}[Hasse]\label{ellipticcurves:thm-hasse}
Let \(E\) be an elliptic curve over \(\Fp\). Write
\[
  \#E(\Fp) = p + 1 - t,
\]
which \emph{defines} the integer \(t\), the \emph{trace of Frobenius}. Then
\[
  |t| \leq 2\sqrt{p}.
\]
Equivalently, \(\#E(\Fp)\) lies in the \emph{Hasse interval} \([\,p + 1 - 2\sqrt{p},\ p + 1 + 2\sqrt{p}\,]\).
\end{theorem}

\begin{proof}[Proof (sketch)]
\emph{Sketch.} Heuristically, as \(x\) ranges over \(\Fp\), the value \(x^3 + ax + b\) is a nonzero square about half the time (giving two points), a nonsquare about half the time (giving none), and zero occasionally (giving one). Using the Legendre symbol \(\legendre{\cdot}{p}\), the number of affine points is
\[
  \#E^{\mathrm{aff}}(\Fp) = \sum_{x \in \Fp}\left(1 + \legendre{x^3 + ax + b}{p}\right) = p + \sum_{x \in \Fp}\legendre{x^3 + ax + b}{p},
\]
so \(t = -\sum_{x}\legendre{x^3 + ax + b}{p}\). The sum is a character sum that exhibits square-root cancellation; the sharp bound \(|t| \leq 2\sqrt p\) is Hasse's theorem. The conceptual proof realizes \(t\) as the trace of the \(p\)-power \emph{Frobenius endomorphism} \(\pi\colon (x,y)\mapsto (x^p, y^p)\) acting on the curve. One shows \(\pi\) satisfies the characteristic equation \(\pi^2 - [t]\pi + [p] = 0\) in the endomorphism ring, and that the associated quadratic form (the degree) is positive definite; the discriminant condition \(t^2 - 4p \le 0\) for the (complex) roots then yields \(|t| \le 2\sqrt p\). The full argument requires the theory of isogenies and the Weil pairing, beyond the present scope.
\end{proof}

Hasse's theorem states that the order of \(E(\Fp)\) is \emph{very close} to \(p\): it cannot deviate from \(p + 1\) by more than \(2\sqrt p\), a relative error of order \(p^{-1/2}\). For cryptography this is crucial: it establishes that the group is large (about \(p\) elements, so about \(255\) bits for the Pasta curves) without any brute-force point count, and point-counting algorithms (Schoof--Elkies--Atkin) compute the \emph{exact} order in polynomial time.

For the cryptographic application one wants \(E(\Fp)\) to contain a large \emph{prime-order} subgroup. The general structure theory of \(E(\Fp)\) as an abelian group is not needed here, because the Pasta curves are designed so that \(\#E(\Fp)\) is itself prime: a group of prime order is cyclic (Corollary~\ref{groups:cor-prime-cyclic}), so \(E(\Fp)\) is cyclic of prime order, the optimal case.

\subsection{Torsion, the cofactor, and prime-order subgroups}

\begin{definition}[order of a point, torsion]\label{ellipticcurves:def-torsion}
The \emph{order} of a point \(P \in E(\Fp)\) is the least integer \(m \geq 1\) with \([m]P = \mathcal{O}\) (it exists and divides \(\#E(\Fp)\) by Lagrange's theorem). For an integer \(m \geq 1\), the \emph{\(m\)-torsion subgroup} is
\[
  E[m] = \{\, P \in E(\overline{\Fp}) : [m]P = \mathcal{O} \,\},
\]
and \(E(\Fp)[m] = E[m]\cap E(\Fp)\) is its rational part.
\end{definition}

\begin{definition}[cofactor]\label{ellipticcurves:def-cofactor}
Suppose \(\#E(\Fp) = h \cdot r\) where \(r\) is the (large) prime to be used and \(h\) is the remaining factor. The integer \(h\) is the \emph{cofactor}, and \(r\) the \emph{order} of the prime-order subgroup. A point \(P\) of order \(r\) generates a cyclic subgroup \(\langle P\rangle \cong \Z/r\Z\), the \emph{prime-order subgroup} used in cryptography.
\end{definition}

A nontrivial cofactor \(h > 1\) is undesirable: it admits ``small-subgroup'' points (of order dividing \(h\)) that can leak information or break protocol assumptions, so implementations either clear the cofactor by multiplying inputs by \(h\), or use prime-order curves where \(h = 1\). The Pasta curves have \emph{cofactor \(1\)}: \(\#E_p(\Fp) = q\) and \(\#E_q(\Fq) = p\) are themselves prime, so \(E_p(\Fp) \cong \Z/q\Z\) and \(E_q(\Fq) \cong \Z/p\Z\) are cyclic of prime order and no cofactor clearing is needed. This is one of the principal design advantages of the Pasta cycle for proof systems.

\begin{proposition}[two-torsion]\label{ellipticcurves:prop-twotorsion}
The affine points of order \(2\) in \(E(\Fp)\) are exactly the points \((x_0, 0)\) with \(f(x_0) = x_0^3 + a x_0 + b = 0\). Hence \(E(\Fp)\) has \(0\), \(1\), or \(3\) points of order \(2\) according as \(f\) has \(0\), \(1\), or \(3\) roots in \(\Fp\), and \(\#E(\Fp)\) is even iff \(f\) has a root in \(\Fp\).
\end{proposition}

\begin{proof}
\(P = (x,y)\) has order \(2\) iff \(P = -P = (x, -y)\) and \(P \neq \mathcal{O}\), i.e. \(y = -y\), i.e. \(y = 0\) (as \(\mathrm{char}\neq 2\)); the curve equation then forces \(f(x) = 0\). Together with \(\mathcal{O}\) these form the \(2\)-torsion subgroup \(E(\Fp)[2]\), which is trivial, \(\Z/2\Z\), or \((\Z/2\Z)^2\); its size is even iff a root exists, and by Lagrange evenness of the size is equivalent to evenness of \(\#E(\Fp)\).
\end{proof}

\begin{remark}\label{ellipticcurves:rem-pasta-odd}
For \(y^2 = x^3 + 5\) over \(\Fp\): the curve has a point of order \(2\) iff \(x^3 + 5 = 0\) has a solution, i.e. iff \(-5\) is a cube in \(\Fp\). For the Pasta primes \(\#E_p(\Fp) = q\) is odd (in fact prime), so there are \emph{no} rational \(2\)-torsion points; consistently, \(-5\) is not a cube in \(\Fp\).
\end{remark}

\subsection{The elliptic-curve discrete logarithm problem}\label{ellipticcurves:subsec-ecdlp}

\begin{definition}[ECDLP]\label{ellipticcurves:def-ecdlp}
Let \(P \in E(\Fp)\) be a point of large prime order \(r\), generating \(G = \langle P\rangle\). The \emph{elliptic-curve discrete logarithm problem} (ECDLP) is: given \(P\) and a point \(Q \in G\), find the unique integer \(n \in \{0, 1, \ldots, r-1\}\) with \(Q = [n]P\). One writes \(n = \log_P Q\).
\end{definition}

The map \(n \mapsto [n]P\) is a group isomorphism \(\Z/r\Z \to G\), so a discrete logarithm always exists and is unique; the content lies in its \emph{computational} difficulty. Scalar multiplication (the forward map) is fast by double-and-add, but no efficient algorithm is known to invert it on a well-chosen elliptic curve.

\begin{proposition}[square-root attacks]\label{ellipticcurves:prop-ecdlp-sqrt}
The baby-step giant-step method solves the ECDLP in a group of prime order \(r\) in \(O(\sqrt{r})\) group operations and \(O(\sqrt r)\) space; Pollard's rho method achieves the same expected time with \(O(1)\) space.
\end{proposition}

\begin{proof}[Proof (sketch)]
Baby-step giant-step writes \(n = i\lceil\sqrt r\rceil + j\) with \(0 \le i, j < \lceil\sqrt r\rceil\), precomputes the ``baby steps'' \([j]P\) in a table, and matches against the ``giant steps'' \(Q - [i\lceil\sqrt r\rceil]P\); a collision reveals \(n\). Pollard's rho achieves the same time with negligible space by a pseudo-random walk on \(G\) and the birthday paradox: a self-collision in the walk yields a relation \([a]P + [b]Q = [a']P + [b']Q\), from which \(n \equiv (a - a')(b' - b)^{-1} \pmod r\). Both are exponential in the bit length \(\log_2 r\).
\end{proof}

\begin{remark}[Absence of better attacks]\label{ellipticcurves:rem-ecdlp-hardness}
That nothing beats the square-root bound on a well-chosen curve is an empirical claim, not a theorem: after four decades of study, no algorithm faster than \(O(\sqrt r)\) is known for general elliptic curves. For \(r\) of cryptographic size (\(r \approx 2^{255}\) for the Pasta curves) this is \(\approx 2^{128}\) operations, infeasible. Curves with special structure can be weaker: if \(\#E(\Fp) = p\) (anomalous curves) there is a polynomial-time attack via the formal group, and the MOV/Frey--Rück attack transfers the ECDLP to a finite-field DLP when the \emph{embedding degree} is small. Cryptographically secure curves, including the Pasta curves, are chosen to have a large prime subgroup order, large embedding degree, and \(\#E(\Fp) \neq p\), defeating all known subexponential attacks.
\end{remark}

\begin{remark}\label{ellipticcurves:rem-index-calculus}
The contrast with multiplicative groups \(\Fp^{\times}\) is essential: in \(\Fp^{\times}\) the index-calculus method solves the discrete logarithm in subexponential time \(L_p[1/3]\), forcing those groups to be very large. On a well-chosen elliptic curve no index calculus is known, so the generic \(O(\sqrt r)\) bound stands and much smaller groups suffice (about \(256\) bits for \(128\)-bit security, versus thousands of bits for \(\Fp^{\times}\)). Smaller groups mean faster arithmetic and shorter keys, which is why elliptic curves dominate modern protocols, including Halo~2.
\end{remark}

\subsection{The \(j\)-invariant and the GLV endomorphism}

\begin{definition}[\(j\)-invariant]\label{ellipticcurves:def-jinvariant}
The \emph{\(j\)-invariant} of the curve \(y^2 = x^3 + ax + b\) (with \(4a^3 + 27 b^2 \neq 0\)) is
\[
  j = 1728\,\frac{4 a^3}{4 a^3 + 27 b^2}\ \in K .
\]
\end{definition}

\begin{proposition}\label{ellipticcurves:prop-jinvariant}
Two elliptic curves over \(\overline{K}\) are isomorphic (over \(\overline{K}\)) if and only if they have the same \(j\)-invariant. The admissible coordinate changes preserving short Weierstrass form are \((x, y) \mapsto (u^2 x, u^3 y)\) for \(u \in \overline{K}^{\times}\), under which \((a, b) \mapsto (u^{-4} a, u^{-6} b)\); the combination \(j\) is invariant.
\end{proposition}

\begin{proof}[Proof (sketch)]
\emph{Sketch.} A direct computation shows the only isomorphisms between short Weierstrass curves fixing \(\mathcal{O}\) are the scalings \((x,y)\mapsto (u^2 x, u^3 y)\); substituting into the equation gives \(a \mapsto u^{-4}a\), \(b\mapsto u^{-6}b\). Then \(4a^3/(4a^3 + 27b^2)\) is unchanged because numerator and denominator both scale by \(u^{-12}\). Conversely, given equal \(j\)-invariants one solves for a suitable \(u\) to construct an isomorphism (with separate easy cases \(j = 0\) and \(j = 1728\)).
\end{proof}

\begin{remark}\label{ellipticcurves:rem-pasta-j}
For \(y^2 = x^3 + b\) one has \(a = 0\), so \(j = 0\). The Pasta curves both have \(j = 0\). Curves with \(j = 0\) are special: they admit an extra automorphism of order \(3\), which underlies the efficient endomorphism described next. (At the other extreme, \(b = 0\) gives \(j = 1728\) and an order-\(4\) automorphism.)
\end{remark}

The vanishing \(j = 0\) is not accidental in the Pasta design; it is what makes the GLV (Gallant--Lambert--Vanstone) speedup available, halving the cost of scalar multiplication.

\begin{proposition}[the GLV endomorphism for \(y^2 = x^3 + b\)]\label{ellipticcurves:prop-glv}
Suppose \(\Fp\) contains a primitive cube root of unity \(\omega\), i.e. \(p \equiv 1 \pmod 3\), so \(\omega \in \Fp\) with \(\omega^2 + \omega + 1 = 0\). On \(E : y^2 = x^3 + b\) the map
\[
  \phi(x, y) = (\omega\, x,\ y), \qquad \phi(\mathcal{O}) = \mathcal{O},
\]
is a group endomorphism of \(E(\Fp)\). Let \(G = \langle P\rangle\) be a subgroup of prime order \(r\) with \(r \equiv 1 \pmod 3\) and \(\phi(G) \subseteq G\). Then \(\phi\) acts on \(G\) as multiplication by a scalar: there is \(\lambda \in \Z/r\Z\) with \(\lambda^2 + \lambda + 1 \equiv 0 \pmod r\) and
\[
  \phi(Q) = [\lambda]\,Q \quad\text{for all } Q \in G .
\]
For the Pasta curves the hypothesis \(\phi(G) \subseteq G\) is automatic, since \(E(\Fp)\) itself is the prime-order group \(G\), and \(r \equiv 1 \pmod 3\) holds.
\end{proposition}

\begin{proof}
First, \(\phi\) maps the curve to itself: if \(y^2 = x^3 + b\) then \((y)^2 = (\omega x)^3 + b\) because \(\omega^3 = 1\), so \((\omega x, y) \in E\). One can check that \(\phi\) is a homomorphism from the addition formulas~\eqref{ellipticcurves:eq-add}: scaling all \(x\)-coordinates by \(\omega\) and leaving \(y\)-coordinates fixed is compatible with the slope \(\lambda_{\mathrm{line}} = (y_2 - y_1)/(\omega x_2 - \omega x_1) = \omega^{-1}\lambda_{\mathrm{line}}^{\mathrm{old}}\) and with \(x_3 = \lambda^2 - x_1 - x_2 \mapsto \omega^{-2} \lambda^2 - \omega x_1 - \omega x_2 = \omega\lambda^2 - \omega x_1 - \omega x_2 = \omega(\ldots)\), using \(\omega^{-2} = \omega\); a short computation confirms \(\phi(P_1 + P_2) = \phi(P_1) + \phi(P_2)\). Now \(\phi\) has order \(3\) (since \(\omega^3 = 1\) and \(\omega \neq 1\)), and \(\phi^2 + \phi + 1 = 0\) as an endomorphism because \(\omega^2 + \omega + 1 = 0\). Restricted to the cyclic group \(G \cong \Z/r\Z\) (which \(\phi\) preserves by hypothesis), every endomorphism is multiplication by some \(\lambda \in \Z/r\Z\) (as \(\Hom(\Z/r\Z, \Z/r\Z) \cong \Z/r\Z\)); the relation \(\phi^2 + \phi + 1 = 0\) becomes \(\lambda^2 + \lambda + 1 \equiv 0 \pmod r\). Consistently, this congruence is solvable precisely when \(r \equiv 1 \pmod 3\), the hypothesis of the statement.
\end{proof}

\begin{remark}[how GLV speeds up scalar multiplication]\label{ellipticcurves:rem-glv-speedup}
Because computing \(\phi\) takes a single field multiplication \(x \mapsto \omega x\) (negligible compared to a group operation), one can decompose a scalar \(n \in \Z/r\Z\) as
\[
  [n]P = [n_1]P + [n_2]\,\phi(P), \qquad n \equiv n_1 + n_2\lambda \pmod r,
\]
where \(n_1, n_2\) are each only about \(\frac{1}{2}\log_2 r\) bits long (a lattice reduction in the lattice of \((a, b)\) with \(a + b\lambda \equiv 0 \pmod r\) finds them). The two half-length scalar multiplications, run simultaneously with shared doublings (Straus--Shamir's trick), cost about half the doublings of a single full-length double-and-add. This is a substantial speedup, and its availability is one reason the Pasta curves take the form \(y^2 = x^3 + b\) with \(p \equiv 1 \pmod 3\).
\end{remark}

\subsection{A worked toy example}

We illustrate the entire machinery below on a curve small enough to compute by hand, followed by the corresponding facts for the Pasta curves.

\begin{example}[the curve \(y^2 = x^3 + 5\) over \(\F_{11}\)]\label{ellipticcurves:exa-toy}
Take \(p = 11\) (note \(11 > 3\), and \(4\cdot 0 + 27\cdot 5^2 = 675 \equiv 675 - 61\cdot 11 = 4 \not\equiv 0 \pmod{11}\), so the curve is nonsingular). The squares in \(\F_{11}\) are
\[
  \{0^2, \ldots, 10^2\} = \{0, 1, 4, 9, 5, 3\},
\]
i.e. the nonzero quadratic residues are \(\{1, 3, 4, 5, 9\}\). Tabulating \(f(x) = x^3 + 5 \bmod 11\) and reading off the \(y\)-values:
\[
\begin{array}{c|c|c|l}
  x & x^3 \bmod 11 & f(x) = x^3 + 5 & y \text{ with } y^2 = f(x) \\\hline
  0 & 0 & 5  & \pm 4 \ (4^2 = 16 \equiv 5) \\
  1 & 1 & 6  & \text{none} \ (6 \notin \text{QR}) \\
  2 & 8 & 2  & \text{none} \\
  3 & 5 & 10 & \text{none} \\
  4 & 9 & 3  & \pm 5 \ (5^2 = 25 \equiv 3) \\
  5 & 4 & 9  & \pm 3 \\
  6 & 7 & 1  & \pm 1 \\
  7 & 2 & 7  & \text{none} \\
  8 & 6 & 0  & 0 \\
  9 & 3 & 8  & \text{none} \\
  10 & 10 & 4 & \pm 2 \\
\end{array}
\]
Counting: \(x \in \{0,4,5,6,10\}\) each give \(2\) points (\(10\) points), \(x = 8\) gives \(1\) point \((8, 0)\), and adding \(\mathcal{O}\) yields
\[
  \#E(\F_{11}) = 10 + 1 + 1 = 12 .
\]
Hasse's bound: \(p + 1 = 12\), so \(t = 0\) and indeed \(|t| = 0 \leq 2\sqrt{11} \approx 6.63\). The single point \((8,0)\) has order \(2\) (here \(-5 \equiv 6\) is a cube, namely \(8^3 = 512 \equiv 6\)), consistent with \(\#E\) being even.

\smallskip
\emph{An addition.} Let \(P = (0, 4)\) and \(Q = (5, 3)\). Since \(x_1 = 0 \neq 5 = x_2\), use the chord:
\[
  \lambda = \frac{3 - 4}{5 - 0} = \frac{-1}{5} = (-1)\cdot 5^{-1} = (-1)\cdot 9 = -9 \equiv 2 \pmod{11},
\]
since \(5\cdot 9 = 45 \equiv 1\). Then
\[
  x_3 = \lambda^2 - x_1 - x_2 = 4 - 0 - 5 = -1 \equiv 10, \qquad
  y_3 = \lambda(x_1 - x_3) - y_1 = 2(0 - 10) - 4 = -24 \equiv 9 \pmod{11}.
\]
So \(P + Q = (10, 9)\), which is on the curve: \(9^2 = 81 \equiv 4\) and \(10^3 + 5 = 1005 \equiv 4 \pmod{11}\). \checkmark

\smallskip
\emph{A doubling.} Take \(P = (0, 4)\). With \(a = 0\),
\[
  \lambda = \frac{3 x_1^2 + a}{2 y_1} = \frac{0}{8} = 0, \quad
  x_3 = \lambda^2 - 2 x_1 = 0, \quad
  y_3 = \lambda(x_1 - x_3) - y_1 = -4 \equiv 7 .
\]
So \([2](0,4) = (0, 7) = -(0,4)\); hence \([3](0,4) = \mathcal{O}\), so \(P = (0,4)\) has order \(3\).

\smallskip
\emph{The GLV endomorphism.} Since \(11 \equiv 2 \pmod 3\), there is \emph{no} primitive cube root of unity in \(\F_{11}\), so the GLV map is not available over \(\F_{11}\); it would require \(p \equiv 1 \pmod 3\). (The Pasta primes do satisfy \(p, q \equiv 1 \pmod 3\), so GLV applies there.) For a small GLV-friendly illustration one could instead take \(p = 13 \equiv 1 \pmod 3\), where \(\omega = 3\) satisfies \(3^2 + 3 + 1 = 13 \equiv 0\), and then \(\phi(x,y) = (3x, y)\) is a nontrivial endomorphism of \(y^2 = x^3 + b\) over \(\F_{13}\).
\end{example}

\begin{example}[the Pasta curves]\label{ellipticcurves:exa-pasta}
For the running example \(y^2 = x^3 + 5\):
\begin{itemize}
  \item Over \(\Fp\) (Pallas, with \(p\) the \(255\)-bit Pasta prime): \(\#E_p(\Fp) = q\), a prime; the group is cyclic, \(E_p(\Fp) \cong \Z/q\Z\); cofactor \(h = 1\); the trace is \(t = p + 1 - q\), with \(|t| \le 2\sqrt p\).
  \item Over \(\Fq\) (Vesta): \(\#E_q(\Fq) = p\), a prime; \(E_q(\Fq) \cong \Z/p\Z\); cofactor \(1\).
  \item Both have \(j = 0\) (since \(a = 0\)), admit the order-\(3\) endomorphism \(\phi(x,y) = (\omega x, y)\) with \(\omega\) a primitive cube root of unity (which exists since \(p, q \equiv 1 \pmod 3\)), and so support the GLV speedup.
  \item The ``cycle'' property \(\#E_p(\Fp) = q\) and \(\#E_q(\Fq) = p\) means the scalar field of one curve is the base field of the other and vice versa; this is exactly what lets recursive proof systems chain proofs across the two curves without expensive non-native field arithmetic. The base field of Pallas is \(\Fp\) and its scalar field (the order of the group) is \(\Fq\), matching the base field of Vesta, and symmetrically.
\end{itemize}
The discrete logarithm problem in \(E_p(\Fp)\) and \(E_q(\Fq)\) is believed to require about \(2^{128}\) operations (half the bit length of \(r \approx 2^{255}\)), the security level targeted by Halo~2 and Zcash Orchard.
\end{example}
\subsection{Summary}

We have constructed the elliptic curve group from the ground up: a nonsingular short Weierstrass cubic over a field of characteristic \(\neq 2,3\), its affine points together with a point at infinity, and the chord-and-tangent addition with explicit formulas~\eqref{ellipticcurves:eq-add}--\eqref{ellipticcurves:eq-double} and all edge cases. The point at infinity is the identity, \((x,y)\) and \((x,-y)\) are inverses, commutativity is immediate, and associativity, though geometrically subtle, makes \(E(K)\) an abelian group. Projective and Jacobian coordinates remove inversions from the inner loop, and double-and-add computes \([n]P\) in \(O(\log n)\) group operations. Over \(\Fp\) the group has order \(p + 1 - t\) with \(|t| \le 2\sqrt p\) (Hasse); choosing the order to be prime (cofactor \(1\), as for the Pasta curves) yields a clean cyclic group in which the discrete logarithm problem is hard. Finally, the \(j = 0\) curves \(y^2 = x^3 + b\) admit the efficiently computable GLV endomorphism. The Pasta curves \(y^2 = x^3 + 5\) over \(\Fp\) and \(\Fq\) realize all of these features simultaneously and form the \(2\)-cycle that carries the recursive proofs of Halo~2 and Zcash Orchard.

\clearpage

\section{Probability, asymptotics, and computation}\label{sec:probability}

Cryptography is the discipline of rendering certain computations easy for legitimate parties and infeasible for adversaries. Stating and proving such guarantees requires three intertwined languages: the language of \emph{probability}, which permits reasoning about random keys, random challenges, and the probability that an attack succeeds; the language of \emph{asymptotics}, which permits precise description of how resources and success probabilities scale; and the language of \emph{computation}, which fixes the notion of an efficient algorithm and of what it means for a problem to be hard. This section develops all three from elementary foundations. The only prerequisites are familiarity with finite sets, sums, and the elementary calculus of limits; we construct everything else here.

Throughout, all probability spaces are finite. This is not a serious restriction for the present purposes: keys, messages, group elements, and the random coins of any terminating algorithm all reside in finite sets. Working with finite spaces avoids measure theory entirely while sacrificing nothing we require.

\subsection{Finite probability spaces and events}

The basic object is a finite set of \emph{outcomes} together with an assignment of weights summing to one.

\begin{definition}[Finite probability space]\label{probability:def-space}
A \emph{finite probability space} is a pair \((\Omega, \Pr)\) where \(\Omega\) is a nonempty finite set, called the \emph{sample space}, and \(\Pr\colon \Omega \to [0,1]\) is a function, called the \emph{probability mass function}, satisfying the normalisation condition
\[
\sum_{\omega \in \Omega} \Pr[\omega] = 1.
\]
We call an element \(\omega \in \Omega\) an \emph{outcome} or an \emph{elementary event}.
\end{definition}

\begin{definition}[Event]\label{probability:def-event}
An \emph{event} is a subset \(A \subseteq \Omega\). Its \emph{probability} is
\[
\Pr[A] = \sum_{\omega \in A} \Pr[\omega],
\]
with the convention \(\Pr[\varnothing] = 0\) (an empty sum). The event \(A\) \emph{occurs} on outcome \(\omega\) if \(\omega \in A\).
\end{definition}

Because \(\Omega\) is finite, every subset is an event; we require no separate \(\sigma\)-algebra. The following elementary facts hold, and we use them freely.

\begin{proposition}[Basic axioms]\label{probability:prop-axioms}
For any finite probability space \((\Omega,\Pr)\) and events \(A, B \subseteq \Omega\):
\begin{enumerate}
\item \(0 \le \Pr[A] \le 1\), \(\Pr[\varnothing] = 0\), and \(\Pr[\Omega] = 1\).
\item \emph{(Complement.)} \(\Pr[\Omega \setminus A] = 1 - \Pr[A]\).
\item \emph{(Monotonicity.)} If \(A \subseteq B\) then \(\Pr[A] \le \Pr[B]\).
\item \emph{(Inclusion--exclusion for two events.)} \(\Pr[A \cup B] = \Pr[A] + \Pr[B] - \Pr[A \cap B]\).
\item \emph{(Finite additivity.)} If \(A_1, \dots, A_n\) are pairwise disjoint then \(\Pr\!\left[\bigcup_{i=1}^n A_i\right] = \sum_{i=1}^n \Pr[A_i]\).
\end{enumerate}
\end{proposition}

\begin{proof}
Every claim follows by manipulating the defining sums. For (1), each summand \(\Pr[\omega]\) lies in \([0,1]\) and the total is \(1\) by normalisation; an event's probability is a partial sum of nonnegative terms, hence in \([0,1]\). For (2), \(\Pr[A] + \Pr[\Omega \setminus A] = \sum_{\omega \in \Omega}\Pr[\omega] = 1\) since \(A\) and \(\Omega \setminus A\) partition \(\Omega\). For (3), \(\Pr[B] = \sum_{\omega \in B}\Pr[\omega] = \sum_{\omega \in A}\Pr[\omega] + \sum_{\omega \in B\setminus A}\Pr[\omega] \ge \Pr[A]\), the last term being a sum of nonnegative numbers. For (5), disjointness implies that each \(\omega\) in the union lies in exactly one \(A_i\), so the double sum \(\sum_i \sum_{\omega \in A_i}\Pr[\omega]\) counts each \(\omega \in \bigcup_i A_i\) exactly once. For (4), write \(A \cup B = A \sqcup (B \setminus A)\) (disjoint), so by (5) \(\Pr[A\cup B] = \Pr[A] + \Pr[B \setminus A]\); also \(B = (A\cap B) \sqcup (B \setminus A)\), giving \(\Pr[B \setminus A] = \Pr[B] - \Pr[A \cap B]\). Substituting yields the claim.
\end{proof}

The single most important example, ubiquitous in cryptography, is the uniform distribution.

\begin{definition}[Uniform distribution]\label{probability:def-uniform}
Let \(S\) be a nonempty finite set. The \emph{uniform distribution} on \(S\) is the probability space \((S, \Pr)\) with \(\Pr[\omega] = 1/|S|\) for every \(\omega \in S\). The notation \(x \leftarrow S\) (or \(x \xleftarrow{\$} S\)) denotes that we sample \(x\) according to the uniform distribution on \(S\). Under the uniform distribution, for any event \(A \subseteq S\),
\[
\Pr[A] = \frac{|A|}{|S|}.
\]
\end{definition}

Thus, for uniform sampling, probability is exactly the ratio of favourable outcomes to total outcomes, recovering the classical "counting" notion of probability.

\subsection{Conditional probability and independence}

Frequently we learn that one event has occurred, and we must update the probability of another accordingly. Conditioning captures this.

\begin{definition}[Conditional probability]\label{probability:def-cond}
Let \(A, B\) be events with \(\Pr[B] > 0\). The \emph{conditional probability of \(A\) given \(B\)} is
\[
\Pr[A \mid B] = \frac{\Pr[A \cap B]}{\Pr[B]}.
\]
When \(\Pr[B] = 0\) we leave the quantity undefined.
\end{definition}

Conditioning on \(B\) restricts the sample space to \(B\) and renormalises so that \(B\) has probability \(1\): the map \(A \mapsto \Pr[A \mid B]\) is itself a probability mass function on \(B\). Two immediate and indispensable consequences follow.

\begin{proposition}[Chain rule and total probability]\label{probability:prop-chain}
Let \(A_1, \dots, A_n\) be events with \(\Pr[A_1 \cap \cdots \cap A_{n-1}] > 0\). Then
\[
\Pr[A_1 \cap \cdots \cap A_n] = \Pr[A_1]\,\Pr[A_2 \mid A_1]\,\Pr[A_3 \mid A_1 \cap A_2]\cdots \Pr[A_n \mid A_1 \cap \cdots \cap A_{n-1}].
\]
Moreover, if \(B_1, \dots, B_k\) partition \(\Omega\) (pairwise disjoint with union \(\Omega\)) and each \(\Pr[B_i] > 0\), then for any event \(A\),
\[
\Pr[A] = \sum_{i=1}^{k} \Pr[A \mid B_i]\,\Pr[B_i].
\]
\end{proposition}

\begin{proof}
The chain rule follows by telescoping: each factor \(\Pr[A_{j} \mid A_1 \cap \cdots \cap A_{j-1}]\) equals \(\Pr[A_1 \cap \cdots \cap A_j]/\Pr[A_1 \cap \cdots \cap A_{j-1}]\) by Definition~\ref{probability:def-cond}, and the product collapses to \(\Pr[A_1 \cap \cdots \cap A_n]/1\). (Monotonicity ensures all intermediate intersections have positive probability, so every factor is defined.) For total probability, the events \(A \cap B_i\) are pairwise disjoint with union \(A\) (since the \(B_i\) partition \(\Omega\)), so by finite additivity \(\Pr[A] = \sum_i \Pr[A \cap B_i] = \sum_i \Pr[A \mid B_i]\,\Pr[B_i]\), using \(\Pr[A \cap B_i] = \Pr[A \mid B_i]\Pr[B_i]\).
\end{proof}

Independence is the formal statement that one event carries no information about another.

\begin{definition}[Independence of events]\label{probability:def-indep}
Two events \(A, B\) are \emph{independent} if
\[
\Pr[A \cap B] = \Pr[A]\,\Pr[B].
\]
A family \(\{A_i\}_{i \in I}\) of events is \emph{mutually independent} if for every finite subset \(J \subseteq I\),
\[
\Pr\!\left[\bigcap_{i \in J} A_i\right] = \prod_{i \in J} \Pr[A_i].
\]
\end{definition}

\begin{remark}\label{probability:rem-indep}
When \(\Pr[B] > 0\), independence of \(A\) and \(B\) is equivalent to \(\Pr[A \mid B] = \Pr[A]\): conditioning on \(B\) does not change the probability of \(A\). Mutual independence is strictly stronger than pairwise independence. The standard counterexample is the following: toss two fair coins, and let \(A\) = "first is heads", \(B\) = "second is heads", \(C\) = "the two coins differ". Each pair is independent, yet \(\Pr[A \cap B \cap C] = 0 \neq \tfrac18 = \Pr[A]\Pr[B]\Pr[C]\), so the three are not mutually independent.
\end{remark}

\subsection{Random variables and expectation}

A random variable assigns a value to each outcome; it summarises a random experiment by a number (or a vector, etc.).

\begin{definition}[Random variable]\label{probability:def-rv}
Let \((\Omega, \Pr)\) be a finite probability space and \(T\) any set. A \emph{(\(T\)-valued) random variable} is a function \(X\colon \Omega \to T\). When \(T \subseteq \R\), we call \(X\) \emph{real-valued}. For \(t \in T\),
\[
\Pr[X = t] = \Pr[\{\omega \in \Omega : X(\omega) = t\}],
\]
and more generally \(\Pr[X \in B]\) denotes the probability that \(X\) takes a value in \(B \subseteq T\). The function \(t \mapsto \Pr[X=t]\) on the finite image of \(X\) is the \emph{distribution} (or \emph{law}) of \(X\).
\end{definition}

\begin{definition}[Independent random variables]\label{probability:def-indep-rv}
Random variables \(X_1, \dots, X_n\) (with values in sets \(T_1, \dots, T_n\)) are \emph{mutually independent} if for all \(t_1 \in T_1, \dots, t_n \in T_n\),
\[
\Pr[X_1 = t_1, \dots, X_n = t_n] = \prod_{i=1}^n \Pr[X_i = t_i].
\]
\end{definition}

\begin{definition}[Expectation]\label{probability:def-exp}
Let \(X\colon \Omega \to \R\) be a real-valued random variable on a finite space. Its \emph{expectation} (or \emph{mean}) is
\[
\E[X] = \sum_{\omega \in \Omega} X(\omega)\,\Pr[\omega] = \sum_{t \in \img(X)} t \cdot \Pr[X = t],
\]
where \(\img(X)\) denotes the (finite) image of \(X\). The two expressions agree by grouping outcomes according to their \(X\)-value.
\end{definition}

A convenient and frequently used special case is the following: the expectation of the \emph{indicator} of an event is exactly its probability.

\begin{definition}[Indicator]\label{probability:def-indicator}
For an event \(A\), its \emph{indicator random variable} \(\mathbf{1}_A\colon \Omega \to \{0,1\}\) takes the value \(\mathbf{1}_A(\omega) = 1\) if \(\omega \in A\) and \(0\) otherwise.
\end{definition}

\begin{proposition}[Indicator expectation]\label{probability:prop-indicator}
For any event \(A\), \(\E[\mathbf{1}_A] = \Pr[A]\).
\end{proposition}

\begin{proof}
\(\E[\mathbf{1}_A] = \sum_{\omega} \mathbf{1}_A(\omega)\Pr[\omega] = \sum_{\omega \in A} \Pr[\omega] = \Pr[A]\), since the indicator vanishes off \(A\).
\end{proof}

The principal property of expectation is its linearity, which holds with no independence assumption whatsoever. This renders expectation considerably more tractable than probability.

\begin{theorem}[Linearity of expectation]\label{probability:thm-linearity}
Let \(X, Y\colon \Omega \to \R\) be random variables on a finite space and \(a, b \in \R\). Then
\[
\E[aX + bY] = a\,\E[X] + b\,\E[Y].
\]
More generally, for random variables \(X_1, \dots, X_n\) and scalars \(a_1, \dots, a_n\),
\[
\E\!\left[\sum_{i=1}^n a_i X_i\right] = \sum_{i=1}^n a_i\,\E[X_i].
\]
This holds whether or not the \(X_i\) are independent.
\end{theorem}

\begin{proof}
Directly from the definition,
\[
\E[aX + bY] = \sum_{\omega \in \Omega} \big(aX(\omega) + bY(\omega)\big)\Pr[\omega]
= a\sum_{\omega} X(\omega)\Pr[\omega] + b\sum_{\omega} Y(\omega)\Pr[\omega] = a\E[X] + b\E[Y],
\]
splitting the finite sum and extracting the constants. The \(n\)-term statement follows by induction on \(n\).
\end{proof}

\subsection{The union bound and a birthday calculation}

The most frequently used inequality in cryptographic proofs is the union bound: the probability that at least one of several events occurs is at most the sum of their probabilities. It is crude but requires no independence and almost always suffices.

\begin{theorem}[Union bound / Boole's inequality]\label{probability:thm-union}
For any events \(A_1, \dots, A_n\) in a finite probability space,
\[
\Pr\!\left[\bigcup_{i=1}^n A_i\right] \le \sum_{i=1}^n \Pr[A_i].
\]
\end{theorem}

\begin{proof}
Define the disjoint "first occurrence" events \(B_1 = A_1\) and \(B_i = A_i \setminus (A_1 \cup \cdots \cup A_{i-1})\) for \(i \ge 2\). The \(B_i\) are pairwise disjoint, \(B_i \subseteq A_i\), and \(\bigcup_i B_i = \bigcup_i A_i\). By finite additivity and monotonicity (Proposition~\ref{probability:prop-axioms}),
\[
\Pr\!\left[\bigcup_i A_i\right] = \Pr\!\left[\bigcup_i B_i\right] = \sum_i \Pr[B_i] \le \sum_i \Pr[A_i]. \qedhere
\]
\end{proof}

The complementary inequality for \emph{conjunctions} is equally useful and follows by taking complements.

\begin{corollary}[Union bound for the good event]\label{probability:cor-union}
If each event \(A_i\) holds except with probability at most \(\varepsilon_i\) (that is, \(\Pr[\Omega \setminus A_i] \le \varepsilon_i\)), then all of them hold simultaneously except with probability at most \(\sum_i \varepsilon_i\):
\[
\Pr\!\left[\bigcap_{i=1}^n A_i\right] \ge 1 - \sum_{i=1}^n \varepsilon_i.
\]
\end{corollary}

\begin{proof}
By De Morgan, \(\Omega \setminus \bigcap_i A_i = \bigcup_i (\Omega \setminus A_i)\). Apply the union bound to the complements and take the complement of both sides.
\end{proof}

We now derive the \emph{birthday bound}. It governs how soon collisions appear when sampling uniformly and underlies the generic attacks against \(n\)-bit hash functions and the security loss in many reductions.

\begin{proposition}[Birthday bound]\label{probability:prop-birthday}
Sample \(x_1, \dots, x_k\) independently and uniformly from a set of size \(N\). Let \(C\) be the event that some two of them collide, i.e.\ \(x_i = x_j\) for some \(i < j\). Then
\[
\Pr[C] \le \binom{k}{2}\frac{1}{N} = \frac{k(k-1)}{2N},
\]
and, on the other side,
\[
\Pr[C] \ge 1 - \prod_{i=0}^{k-1}\Big(1 - \frac{i}{N}\Big) \ge 1 - e^{-k(k-1)/(2N)} .
\]
In particular a collision becomes likely once \(k = \Theta(\sqrt{N})\).
\end{proposition}

\begin{proof}
\emph{Upper bound.} For \(i < j\) let \(A_{ij}\) be the event \(x_i = x_j\). Since \(x_i, x_j\) are independent and uniform, \(\Pr[A_{ij}] = \sum_{v} \Pr[x_i = v]\Pr[x_j = v] = N \cdot (1/N)^2 = 1/N\). Now \(C = \bigcup_{i<j} A_{ij}\), a union of \(\binom{k}{2}\) events, so the union bound (Theorem~\ref{probability:thm-union}) gives \(\Pr[C] \le \binom{k}{2}/N\).

\emph{Lower bound.} The complement \(\overline{C}\) is the event that all \(k\) samples are distinct. Counting ordered tuples of distinct values, \(\Pr[\overline{C}] = \prod_{i=0}^{k-1}(1 - i/N)\): having fixed \(i\) distinct values, the next sample avoids them with probability \(1 - i/N\). Using \(1 - x \le e^{-x}\) for all \(x \in \R\),
\[
\Pr[\overline{C}] = \prod_{i=0}^{k-1}\Big(1 - \frac iN\Big) \le \prod_{i=0}^{k-1} e^{-i/N} = \exp\!\Big(-\tfrac1N\sum_{i=0}^{k-1} i\Big) = e^{-k(k-1)/(2N)} .
\]
Taking complements gives \(\Pr[C] = 1 - \Pr[\overline C] \ge 1 - e^{-k(k-1)/(2N)}\).
\end{proof}

\begin{remark}\label{probability:rem-birthday}
The threshold \(k \approx \sqrt{N}\) is the reason an \(n\)-bit hash function (so \(N = 2^n\)) offers only about \(2^{n/2}\) collision resistance: an adversary expects a collision after roughly \(2^{n/2}\) random evaluations. For this reason one chooses output lengths twice the desired security level. The same square-root phenomenon recurs in random-walk discrete-log attacks (Pollard's rho), so designers size the curves used in Halo~2 and Orchard so that \(\sqrt{N}\) is itself astronomically large.
\end{remark}

\subsection{Statistical distance}

Comparing distributions---for example a real key-generation procedure against an idealised uniform one---requires a notion of how far apart two distributions are. The appropriate notion for "no statistical test can distinguish them well" is statistical (total variation) distance.

\begin{definition}[Statistical distance]\label{probability:def-sd}
Let \(P\) and \(Q\) be probability distributions on the same finite set \(S\) (i.e.\ probability mass functions \(P, Q\colon S \to [0,1]\)). Their \emph{statistical distance} (or \emph{total variation distance}) is
\[
\Delta(P, Q) = \frac12 \sum_{s \in S} \big|P(s) - Q(s)\big|.
\]
For random variables \(X, Y\) on \(S\), \(\Delta(X, Y)\) denotes the statistical distance of their distributions.
\end{definition}

The factor \(\tfrac12\) makes \(\Delta\) range over \([0,1]\) and gives it a clean interpretation as the maximum advantage of a distinguisher, which we establish next.

\begin{theorem}[Properties of statistical distance]\label{probability:thm-sd}
Let \(P, Q, R\) be distributions on a finite set \(S\). Then:
\begin{enumerate}
\item \(0 \le \Delta(P,Q) \le 1\); and \(\Delta(P,Q) = 0\) iff \(P = Q\).
\item \emph{(Metric.)} \(\Delta\) is symmetric and satisfies the triangle inequality \(\Delta(P,R) \le \Delta(P,Q) + \Delta(Q,R)\).
\item \emph{(Variational characterisation.)} \(\displaystyle \Delta(P,Q) = \max_{A \subseteq S}\big(P(A) - Q(A)\big) = \max_{A \subseteq S}\big|P(A) - Q(A)\big|\), where \(P(A) = \sum_{s\in A}P(s)\).
\item \emph{(Data-processing / post-processing.)} For any (possibly randomised) function \(f\) applied to a sample, \(\Delta(f(P), f(Q)) \le \Delta(P, Q)\). In particular no algorithm can increase statistical distance.
\end{enumerate}
\end{theorem}

\begin{proof}
(1) Nonnegativity is clear. For the upper bound, let \(S^+ = \{s : P(s) \ge Q(s)\}\) and \(S^- = S \setminus S^+\). Since \(\sum_s P(s) = \sum_s Q(s) = 1\), we have \(\sum_{s}(P(s)-Q(s)) = 0\), so \(\sum_{s \in S^+}(P(s)-Q(s)) = \sum_{s \in S^-}(Q(s)-P(s)) =: D \ge 0\). Hence \(\Delta(P,Q) = \tfrac12\big(\sum_{S^+}(P-Q) + \sum_{S^-}(Q-P)\big) = D\). As \(D = \sum_{S^+}(P(s) - Q(s)) \le \sum_{S^+}P(s) \le 1\), we get \(\Delta \le 1\). If \(\Delta = 0\) every \(|P(s)-Q(s)| = 0\), so \(P=Q\); the converse is immediate.

(2) Symmetry is evident. For the triangle inequality, \(|P(s)-R(s)| \le |P(s)-Q(s)| + |Q(s)-R(s)|\) termwise by the triangle inequality on \(\R\); summing and halving gives the claim.

(3) With \(D\) as above, take \(A = S^+\): then \(P(A) - Q(A) = \sum_{S^+}(P-Q) = D = \Delta(P,Q)\), so the maximum is at least \(\Delta\). Conversely, for \emph{any} \(A\), \(P(A)-Q(A) = \sum_{s\in A}(P(s)-Q(s)) \le \sum_{s\in A,\, P(s)\ge Q(s)}(P(s)-Q(s)) \le D = \Delta\), since dropping the negative terms only increases the sum and extending from \(A\cap S^+\) to all of \(S^+\) only increases it further. Thus the maximum over \(A\) of \(P(A)-Q(A)\) is exactly \(\Delta\); replacing \(A\) by its complement swaps the sign, giving the same value for \(\max|P(A)-Q(A)|\).

(4) First take \(f\) deterministic, \(f\colon S \to U\). For \(u \in U\), \(f(P)(u) = P(f^{-1}(u))\). Then
\[
\Delta(f(P), f(Q)) = \tfrac12\sum_{u}\big|P(f^{-1}(u)) - Q(f^{-1}(u))\big| \le \tfrac12 \sum_u \sum_{s \in f^{-1}(u)} |P(s)-Q(s)| = \Delta(P,Q),
\]
using \(|\sum_s a_s| \le \sum_s |a_s|\) on each fibre, and that the fibres \(f^{-1}(u)\) partition \(S\). A randomised \(f\) is a deterministic function of \(s\) together with independent fresh coins \(r\); applying the deterministic case to the joint variable \((s, r)\) (whose two distributions \(P \times U_r\) and \(Q \times U_r\) have the same statistical distance \(\Delta(P,Q)\), the shared coin distribution contributing nothing) yields the bound.
\end{proof}

\begin{remark}[Distinguishing interpretation]\label{probability:rem-sd}
Part (3) states that \(\Delta(P,Q)\) is exactly the best advantage of any (even computationally unbounded) test that, given a single sample, must guess whether it came from \(P\) or \(Q\): the optimal test outputs "\(P\)" on the set \(A = S^+\), and its advantage \(\Pr_{P}[A] - \Pr_Q[A]\) equals \(\Delta(P,Q)\). Part (4) is what makes statistical distance compose in proofs: applying the same procedure to two close inputs keeps the outputs close. If \(\Delta(P,Q) \le \varepsilon\), then \(P\) and \(Q\) are interchangeable in any analysis at the cost of an additive \(\varepsilon\) in every probability. When \(\varepsilon\) is negligible (defined below), we call \(P\) and \(Q\) \emph{statistically indistinguishable}.
\end{remark}
\subsection{Uniform sampling and the bias of modular reduction}

Cryptography requires uniform elements of a finite set throughout: a uniform scalar in \(\Z/q\Z\), a uniform field element of \(\Fp\), a uniform point in a group. The most tractable source of randomness, however, is a stream of independent uniform bits. We must therefore convert uniform bitstrings into uniform elements of \(S\), and we must understand the bias this conversion introduces.

If \(|S| = 2^\ell\) is a power of two, the conversion is exact: read \(\ell\) uniform bits as a number in \(\{0, \dots, 2^\ell - 1\}\) and index into \(S\). The nontrivial case arises when \(|S|\) is not a power of two---most importantly when \(S = \Z/q\Z\) for a prime \(q\), which is the scalar field of an elliptic curve.

A common and elementary method is \emph{reduce-modulo}: sample a long uniform integer and reduce it mod \(q\). The result is not perfectly uniform---some residues receive one more preimage than others---but the bias is negligible provided the integer is long enough. We make this precise below.

\begin{proposition}[Bias of modular reduction]\label{probability:prop-modbias}
Let \(q \ge 2\) be an integer and let \(K = 2^L\) for some integer \(L\). Sample \(U\) uniformly from \(\{0, 1, \dots, K-1\}\) (equivalently, read \(L\) uniform bits as an integer), and set \(R = U \bmod q \in \{0, 1, \dots, q-1\}\). Let \(\mathcal{U}_q\) denote the uniform distribution on \(\{0,\dots,q-1\}\). Then the statistical distance of \(R\) from uniform satisfies
\[
\Delta(R, \mathcal{U}_q) \le \frac{q}{K} = \frac{q}{2^{L}}.
\]
Consequently, if \(q < 2^n\) and one takes \(L = n + s\) bits, then \(\Delta(R, \mathcal{U}_q) < 2^{-s}\).
\end{proposition}

\begin{proof}
Write \(K = mq + t\) with \(m = \lfloor K/q\rfloor\) and remainder \(t = K \bmod q\), so \(0 \le t < q\). For a residue \(r \in \{0,\dots,q-1\}\), the number of integers in \(\{0,\dots,K-1\}\) congruent to \(r\) modulo \(q\) is \(m+1\) if \(r < t\) and \(m\) if \(r \ge t\). Hence
\[
\Pr[R = r] = \begin{cases} (m+1)/K, & r < t,\\[2pt] m/K, & r \ge t. \end{cases}
\]
Comparing with the target probability \(1/q\) and using \(m/K \le 1/q \le (m+1)/K\) (since \(mq \le K \le (m+1)q\)),
\[
\Big|\Pr[R=r] - \tfrac1q\Big| \le \frac{m+1}{K} - \frac{m}{K} = \frac1K \quad\text{for every } r.
\]
Therefore
\[
\Delta(R, \mathcal{U}_q) = \frac12\sum_{r=0}^{q-1}\Big|\Pr[R=r] - \tfrac1q\Big| \le \frac12 \cdot q \cdot \frac1K = \frac{q}{2K} \le \frac{q}{K}.
\]
(The computation yields the sharper bound \(q/(2K)\); \(q/K\) is the usual stated form.) For the final claim, \(q < 2^n\) and \(K = 2^{n+s}\) give \(q/K < 2^n/2^{n+s} = 2^{-s}\).
\end{proof}

\begin{remark}\label{probability:rem-modbias}
The consequence is as follows: to sample a scalar mod a \(\sim 256\)-bit prime \(q\) with statistical distance below \(2^{-128}\) from uniform, draw \(256 + 128 = 384\) uniform bits and reduce. The bias \(2^{-128}\) is negligible (Definition~\ref{probability:def-negl}), so the reduce-modulo sample is statistically indistinguishable from a perfectly uniform scalar, and by the data-processing inequality (Theorem~\ref{probability:thm-sd}(4)) substituting one for the other anywhere costs at most \(2^{-128}\) per use. An alternative, \emph{rejection sampling}---draw \(\lceil \log_2 q\rceil\) bits, accept if the result is \(< q\), else retry---gives \emph{exactly} uniform output at the cost of a variable but, with overwhelming probability, small number of retries; for the present purposes the negligible bias of reduce-modulo is acceptable and simpler to analyse. This is precisely the kind of distribution that the Orchard specification employs when deriving scalars from hash outputs.
\end{remark}

\subsection{Asymptotic notation}

The discussion now turns from probability to growth rates. To address efficiency and hardness, we must compare the manner in which functions of an integer parameter grow as that parameter tends to infinity. The Bachmann--Landau notation makes such comparisons precise. Throughout, functions map \(\N\) (or a cofinite subset thereof) to the nonnegative reals \(\R_{\ge 0}\), and "for sufficiently large \(n\)" means "for all \(n\) beyond some threshold \(n_0\)".

\begin{definition}[Big-O, Omega, Theta, little-o, little-omega]\label{probability:def-asymp}
Let \(f, g\colon \N \to \R_{\ge 0}\) with \(g\) eventually positive. We define the following:
\begin{itemize}
\item \(f = O(g)\) (read "\(f\) is big-O of \(g\)", an asymptotic upper bound) if there exist constants \(c > 0\) and \(n_0\) with \(f(n) \le c\,g(n)\) for all \(n \ge n_0\).
\item \(f = \Omega(g)\) (asymptotic lower bound) if there exist \(c > 0\) and \(n_0\) with \(f(n) \ge c\,g(n)\) for all \(n \ge n_0\). Equivalently, \(f = \Omega(g)\) iff \(g = O(f)\).
\item \(f = \Theta(g)\) (tight bound) if \(f = O(g)\) and \(f = \Omega(g)\); i.e.\ there are \(c_1, c_2 > 0\) and \(n_0\) with \(c_1 g(n) \le f(n) \le c_2 g(n)\) for \(n \ge n_0\).
\item \(f = o(g)\) (read "little-o", strictly smaller order) if for every \(c > 0\) there is \(n_0\) with \(f(n) \le c\,g(n)\) for all \(n \ge n_0\); equivalently \(\lim_{n\to\infty} f(n)/g(n) = 0\).
\item \(f = \omega(g)\) (strictly larger order) if \(g = o(f)\); equivalently \(\lim_{n\to\infty} f(n)/g(n) = \infty\).
\end{itemize}
\end{definition}

\begin{remark}\label{probability:rem-asymp-eq}
The "\(=\)" in "\(f = O(g)\)" is traditional but denotes set membership: \(O(g)\) is the class of all functions bounded above by a constant multiple of \(g\), and \(f = O(g)\) means \(f \in O(g)\). One never reads such equations right to left. We follow the customary abuse of notation here.
\end{remark}

\begin{example}\label{probability:exp-asymp}
\(3n^2 + 10 n + 7 = \Theta(n^2)\); indeed for \(n \ge 1\) one has \(3n^2 \le 3n^2 + 10n + 7 \le 20 n^2\). Also \(n^2 = o(n^3)\) since \(n^2/n^3 = 1/n \to 0\); \(\log_2 n = o(n)\); \(n^{100} = o(2^n)\) (any polynomial is little-o of any exponential with base \(> 1\)); and \(n! = \omega(2^n)\). Note that \(2^{n} = o(2^{n+1})\) is false, but \(2^{n} = \Theta(2^{n+1})\) is true, since \(2^{n+1} = 2 \cdot 2^n\) is within a constant factor.
\end{example}

\begin{proposition}[Useful closure facts]\label{probability:prop-asymp-closure}
Let \(f_1 = O(g_1)\) and \(f_2 = O(g_2)\). Then \(f_1 + f_2 = O(\max(g_1, g_2)) = O(g_1 + g_2)\) and \(f_1 f_2 = O(g_1 g_2)\). The same closure statements hold with \(O\) replaced uniformly by \(\Omega\) or \(\Theta\). Moreover \(O(\cdot)\) is transitive: \(f = O(g)\) and \(g = O(h)\) imply \(f = O(h)\).
\end{proposition}

\begin{proof}
Choose \(c_1, n_1\) and \(c_2, n_2\) witnessing the two hypotheses, and let \(n_0 = \max(n_1, n_2)\). For \(n \ge n_0\): \(f_1(n) + f_2(n) \le c_1 g_1(n) + c_2 g_2(n) \le (c_1 + c_2)\max(g_1(n), g_2(n))\), and \(\max(g_1, g_2) \le g_1 + g_2 \le 2\max(g_1,g_2)\), so the two right-hand sides agree up to a constant; this proves the sum rule. For products, \(f_1(n) f_2(n) \le c_1 c_2\, g_1(n) g_2(n)\). Transitivity: \(f \le c\,g\) and \(g \le c' h\) eventually give \(f \le cc'\,h\) eventually. The \(\Omega\) statements follow by the equivalence \(f = \Omega(g) \Leftrightarrow g = O(f)\), and \(\Theta\) by combining the two.
\end{proof}

\subsection{Polynomial, exponential, and negligible functions}

We now classify the growth rates relevant to cryptography. Fix a distinguished variable \(\lambda \in \N\), the \emph{security parameter} (discussed in detail in Subsection~\ref{probability:sub-secparam}); informally larger \(\lambda\) means more security and larger keys. Resources (running time, key length) and success probabilities are functions of \(\lambda\).

\begin{definition}[Polynomial and super-polynomial]\label{probability:def-poly}
A function \(f\colon \N \to \R_{\ge 0}\) is \emph{polynomially bounded}, written \(f = \poly(\lambda)\), if there is a constant \(c\) (and a threshold) with \(f(\lambda) \le \lambda^c + c\) for all large \(\lambda\); equivalently \(f(\lambda) = O(\lambda^c)\) for some constant \(c\). It is \emph{super-polynomial} if it grows faster than every polynomial, i.e.\ \(f(\lambda) = \omega(\lambda^c)\) for every constant \(c\) (equivalently \(\lambda^c = o(f(\lambda))\) for all \(c\)).
\end{definition}

\begin{definition}[Exponential]\label{probability:def-exp-fn}
A function \(f\) is \emph{exponential} if \(f(\lambda) = 2^{\Omega(\lambda^{\epsilon})}\) for some constant \(\epsilon > 0\); the prototypical case is \(f(\lambda) = 2^{c\lambda}\) for a constant \(c > 0\). Every exponential function is super-polynomial, but the converse fails: e.g.\ \(\lambda^{\log \lambda} = 2^{(\log\lambda)^2}\) is super-polynomial yet sub-exponential.
\end{definition}

The central cryptographic notion is that of a \emph{negligible} function: one that shrinks faster than the reciprocal of every polynomial. A negligible probability is one we may ignore, because no efficient (polynomially many) repetition of an experiment can amplify it to a noticeable chance.

\begin{definition}[Negligible function]\label{probability:def-negl}
A function \(\nu\colon \N \to \R_{\ge 0}\) is \emph{negligible}, written \(\nu = \negl(\lambda)\), if for every positive constant \(c\) there exists \(\lambda_0\) such that
\[
\nu(\lambda) < \frac{1}{\lambda^{c}} \qquad \text{for all } \lambda \ge \lambda_0.
\]
Equivalently, \(\nu(\lambda) = o(\lambda^{-c})\) for every constant \(c > 0\); equivalently, \(\lambda^{c}\,\nu(\lambda) \to 0\) as \(\lambda \to \infty\) for every constant \(c\). A function is \emph{noticeable} (or non-negligible) if it is not negligible, i.e.\ \(f(\lambda) \ge \lambda^{-c}\) infinitely often for some constant \(c\); it is \emph{overwhelming} if \(1 - f\) is negligible.
\end{definition}

\begin{example}\label{probability:exp-negl}
\(2^{-\lambda}\), \(2^{-\sqrt{\lambda}}\), and \(\lambda^{-\log\lambda}\) are negligible. Indeed for \(2^{-\lambda}\): given any \(c\), since \(2^{-\lambda}\lambda^{c} = \lambda^c / 2^{\lambda} \to 0\) (a polynomial over an exponential), the function meets the definition. By contrast, \(1/\lambda^{100}\) is \emph{not} negligible---taking \(c = 101\) shows \(1/\lambda^{100} > 1/\lambda^{101}\) for all \(\lambda\)---it is noticeable, even though it tends to \(0\). This is the key subtlety: tending to zero does not suffice; a negligible function must beat every inverse polynomial.
\end{example}

The reason "negligible" is the correct closure of "ignorable" is the following stability under the operations that appear in security proofs.

\begin{proposition}[Closure properties of negligible functions]\label{probability:prop-negl-closure}
Let \(\nu, \mu\) be negligible and let \(p\) be a polynomially bounded function. Then:
\begin{enumerate}
\item \(\nu + \mu\) is negligible (the class is closed under addition, hence under any fixed finite sum).
\item \(p \cdot \nu\) is negligible (a polynomial times a negligible function is negligible).
\item If \(f(\lambda) \le \nu(\lambda)\) for all large \(\lambda\), then \(f\) is negligible (domination by a negligible function).
\item Summing \(p(\lambda)\) functions, each bounded by a single negligible function \(\nu\), yields a negligible function; in particular a union bound over polynomially many negligible events is negligible.
\end{enumerate}
\end{proposition}

\begin{proof}
(1) Fix \(c\). Applying the definition of negligibility to the exponent \(c+1\), there is a threshold beyond which \(\nu(\lambda), \mu(\lambda) < \lambda^{-(c+1)} \le \tfrac12 \lambda^{-c}\) (the latter once \(\lambda \ge 2\)); adding the two bounds gives \(\nu(\lambda) + \mu(\lambda) < \lambda^{-c}\) for all large \(\lambda\).

(2) Suppose \(p(\lambda) \le \lambda^{d}\) for large \(\lambda\) (and some constant \(d\)). Fix \(c\). By negligibility applied to the exponent \(c + d\), \(\nu(\lambda) < \lambda^{-(c+d)}\) for large \(\lambda\). Then \(p(\lambda)\nu(\lambda) < \lambda^{d}\lambda^{-(c+d)} = \lambda^{-c}\) for large \(\lambda\). Hence \(p\nu\) is negligible.

(3) Immediate: if \(f \le \nu\) eventually and \(\nu < \lambda^{-c}\) eventually, then \(f < \lambda^{-c}\) eventually, for every \(c\).

(4) Write the sum \(\Sigma(\lambda) = \sum_{i=1}^{p(\lambda)} \nu_i(\lambda)\) where each \(\nu_i \le \nu\) by hypothesis (in security proofs all the \(\nu_i\) arise from one bound). Then \(\Sigma(\lambda) \le p(\lambda)\,\nu(\lambda)\), which is negligible by (2); apply (3).
\end{proof}

\begin{remark}[Why negligible is the threshold of security]\label{probability:rem-negl}
Suppose an adversary wins a game with negligible probability \(\nu(\lambda)\). If it repeats the attack \(p(\lambda)\) times (polynomially many, the most an efficient adversary can afford), the probability that some repetition succeeds is, by the union bound, at most \(p(\lambda)\nu(\lambda)\)---still negligible by Proposition~\ref{probability:prop-negl-closure}(2). Thus negligible success probability remains negligible under any efficient amplification: it is genuinely ignorable. Conversely a noticeable success probability \(\ge \lambda^{-c}\) can be boosted to a constant by \(\lambda^{c}\) repetitions. This dichotomy is the formal backbone of the phrase "the scheme is secure except with negligible probability".
\end{remark}

\subsection{Algorithms, running time, and PPT}

It remains to make "efficient algorithm" and "infeasible" precise. We adopt the standard model: algorithms are Turing machines (or, equivalently up to polynomial factors, programs in any reasonable model), inputs and outputs are finite binary strings, and the number of elementary steps as a function of input length is the resource we measure. Write \(\{0,1\}^*\) for the set of all finite binary strings and \(|x|\) for the length of \(x \in \{0,1\}^*\).

\begin{definition}[Deterministic polynomial-time algorithm]\label{probability:def-pt}
An algorithm \(M\) runs in \emph{(deterministic) polynomial time} if there is a polynomial \(p\) such that, on every input \(x \in \{0,1\}^*\), \(M\) halts within at most \(p(|x|)\) steps and produces an output. The class of decision problems solvable by such machines is \(\mathsf{P}\).
\end{definition}

Determinism is insufficient for cryptography: key generation, sampling, and many algorithms require randomness. We model randomness by granting the machine access to an auxiliary tape of independent uniform random bits (the "coins").

\begin{definition}[Probabilistic algorithm and PPT]\label{probability:def-ppt}
A \emph{probabilistic} (or \emph{randomised}) algorithm \(M\) is a Turing machine with an additional read-only \emph{random tape} initialised with independent uniformly random bits \(r \in \{0,1\}^*\). For a fixed input \(x\), the output \(M(x)\) is a random variable over the choice of \(r\); write \(M(x; r)\) when one wishes to display the coins explicitly, and then \(M(x;\cdot)\) is a deterministic function. The algorithm runs in \emph{probabilistic polynomial time} (it is \emph{PPT}) if there is a polynomial \(p\) such that, for every input \(x\) \emph{and every choice of coins} \(r\), \(M\) halts within \(p(|x|)\) steps. (Bounding the time for every \(r\), not merely in expectation, is the standard "strict" notion and ensures \(M\) ever reads only \(p(|x|)\) coins.)
\end{definition}

\begin{remark}[Non-uniformity and circuits]\label{probability:rem-nonuniform}
There are two standard ways to formalise an efficient adversary. The \emph{uniform} model is a single PPT machine that works for all input lengths, as above. The \emph{non-uniform} model allows a separate algorithm (equivalently, a Boolean circuit) of size \(\poly(\lambda)\) for each \(\lambda\), which may be hard-wired with an arbitrary polynomial-length "advice" string depending on \(\lambda\). Non-uniform adversaries are at least as powerful as uniform ones and are often more convenient in reduction proofs because the advice can encode a "best" attack. The asymptotic theory below is identical for both; security definitions in the Halo~2/Orchard setting typically state the security against non-uniform PPT adversaries, which is the more conservative choice.
\end{remark}

With efficiency pinned down, we can state what it means for a computational task to be hard. The guiding principle is the Cobham--Edmonds thesis: "efficiently computable" = "computable in polynomial time". Its negation gives infeasibility.

\begin{definition}[Computational infeasibility]\label{probability:def-infeasible}
A computational task parameterised by the security parameter \(\lambda\) is \emph{(computationally) feasible} if there is a PPT algorithm solving it (with overwhelming success probability). It is \emph{infeasible} if for \emph{every} PPT algorithm \(\Adv\) the probability that \(\Adv\) solves it is negligible in \(\lambda\):
\[
\Pr\big[\Adv \text{ solves the task on parameter } \lambda\big] = \negl(\lambda),
\]
where the probability is over the task's randomness and \(\Adv\)'s coins. Equivalently, no efficient algorithm does noticeably better than the trivial strategy.
\end{definition}

\begin{remark}[Reading a hardness assumption]\label{probability:rem-hardness}
Cryptographic assumptions are statements of this form. For instance, the discrete-logarithm assumption in a family of groups \(\{\G_\lambda\}\) of order \(q(\lambda)\) asserts: for every PPT \(\Adv\),
\[
\Pr\Big[\Adv(g, g^{x}) = x \;\Big|\; g \leftarrow \text{generator},\; x \leftarrow \Z/q\Z\Big] = \negl(\lambda).
\]
The entire edifice of Halo~2 and Orchard rests on such infeasibility statements, which one makes true in practice by choosing \(\lambda\) (and hence the group size \(\approx 2^{\lambda}\)) large enough that the best known attack---running in time \(\Theta(\sqrt{q}) = \Theta(2^{\lambda/2})\) by the birthday/rho phenomenon of Remark~\ref{probability:rem-birthday}---is astronomically slow. The asymptotic language of this section is what permits the terms "negligible" and "polynomial" without committing to a particular \(\lambda\); the concrete-security viewpoint then instantiates \(\lambda\) (e.g.\ \(\lambda = 128\)) to read off actual bit-counts.
\end{remark}

\subsection{The security parameter}\label{probability:sub-secparam}

This section closes by making explicit the parameter that threads through every preceding definition.

\begin{definition}[Security parameter]\label{probability:def-secparam}
The \emph{security parameter} is a positive integer \(\lambda \in \N\) that one supplies to all algorithms of a cryptographic scheme, conventionally in unary as the string \(1^\lambda\) (so that "polynomial in the input length" coincides with "polynomial in \(\lambda\)"). It governs simultaneously:
\begin{itemize}
\item the \emph{sizes} of objects (key lengths, group order \(\approx 2^\lambda\), field size), each \(\poly(\lambda)\);
\item the \emph{running time} of honest algorithms, which must be \(\poly(\lambda)\); and
\item the \emph{security level}: every PPT adversary's advantage must be \(\negl(\lambda)\).
\end{itemize}
\end{definition}

\begin{remark}[Why unary, and what \(\lambda\) buys]\label{probability:rem-unary}
Passing \(1^\lambda\) (a string of \(\lambda\) ones, of length \(\lambda\)) rather than the integer \(\lambda\) in binary (length \(\log_2\lambda\)) is a deliberate convention: it forces "polynomial-time in the input length" to mean polynomial in \(\lambda\) itself, so an honest key generator may take time \(\poly(\lambda)\) to produce \(\lambda\)-bit keys. As \(\lambda\) grows, two things happen in lockstep: honest parties pay a polynomial cost, while every adversary's success probability falls negligibly and (concretely) the best attack's running time grows exponentially. Asymptotic security is the statement that this gap holds for all sufficiently large \(\lambda\); the practitioner then fixes a single value (commonly \(\lambda = 128\), targeting \(2^{128}\)-operation attacks) and verifies that the concrete sizes---a \(\sim 255\)-bit Pallas/Vesta scalar field, hash outputs of \(\geq 256\) bits per the birthday bound of Proposition~\ref{probability:prop-birthday}, and sampling slack of \(\sim 128\) extra bits per Proposition~\ref{probability:prop-modbias}---deliver it. Every "negligible" in the remainder of this monograph is implicitly a function of this \(\lambda\).
\end{remark}
\end{document}
